% Lesson 51 — Vocabulary: Words and expressions related to exploiting service packages (telephone/internet services) — Anglais Tle A — Cours signé OnBuch+
% Programme officiel MINESEC. Compiler avec : tectonic EN51-vocabulary-words-and-expressions-related-to-exploiting.tex
\newcommand{\DOCMATIERE}{Anglais}
\newcommand{\DOCNIVEAU}{Tle A}
\newcommand{\DOCMODULE}{Chapter 5 : Media and Communication}
\newcommand{\DOCLECON}{EN51}
\newcommand{\DOCTITRE}{Lesson 51 — Vocabulary: Words and expressions related to exploiting service packages (telephone/internet services)}
% OnBuch+ — préambule « cours signés » ENRICHI (compilé avec tectonic / XeLaTeX)
% Système visuel « Pop » d'OnBuch+ : fond crème (jamais blanc), bordures encre
% épaisses, ombres « dures » décalées, titres Archivo Black en capitales.
% Version MATHÉMATIQUES : graphes (pgfplots/tikz), boîtes pédagogiques
% (définition, propriété, méthode, attention, le savais-tu, exercice résolu,
% exercice, TP, bilan), page de garde et signature OnBuch+ ancrée en bas.
%
% Macros attendues dans main.tex AVANT \input{preamble} :
%   \DOCMATIERE  \DOCNIVEAU  \DOCMODULE  \DOCLECON (numéro)  \DOCTITRE
\documentclass[11pt]{article}

\usepackage{fontspec}
\usepackage[english,french]{babel} % français = langue principale ; l'anglais via \eng{..} / textebox
\usepackage[a4paper,margin=2cm,top=3.4cm,bottom=2.8cm,headheight=1.4cm,footskip=1.3cm]{geometry}
\usepackage{amsmath,amssymb,amsfonts}
\usepackage{mathtools} % \xmapsto, \coloneqq... souvent utilisés par les modèles
\usepackage{cancel} % simplifications barrées (bilans de forces)
% \abs{z} (module, valeur absolue) et \norm{u} (norme), utilisés par les modèles
\providecommand{\abs}{}\let\abs\relax\DeclarePairedDelimiter{\abs}{\lvert}{\rvert}
\providecommand{\norm}{}\let\norm\relax\DeclarePairedDelimiter{\norm}{\lVert}{\rVert}
\usepackage[table]{xcolor} % [table] : fournit aussi \rowcolors (alternance de lignes)
\usepackage{pagecolor}
\usepackage{tikz}
\usetikzlibrary{arrows.meta,positioning,calc,shapes.geometric,shapes.misc,shapes.arrows,decorations.pathmorphing,decorations.pathreplacing,decorations.markings,patterns,shadows,mindmap,trees,fit,backgrounds,matrix,intersections,angles,quotes}
\usepackage{pgfplots}
\usepackage[version=4]{mhchem} % équations nucléaires \ce{^{235}_{92}U}
\usepackage[european,siunitx]{circuitikz} % schémas électriques
\pgfplotsset{compat=1.18}
\usepgfplotslibrary{fillbetween,groupplots} % aires entre courbes (calcul intégral)
\usepackage{enumitem}
\usepackage{fancyhdr}
% [most] charge toutes les bibliothèques tcolorbox (listings, minted, raster,
% documentation...) et gonfle inutilement la mémoire TeX sur un document long
% (~300 boîtes) : on ne charge que ce qui est réellement utilisé.
\usepackage[skins,breakable]{tcolorbox}
\usepackage{graphicx}
\usepackage{colortbl}
\usepackage{booktabs}
\usepackage{tabularx}
\usepackage{diagbox} % case d'en-tête diagonale des tableaux à double entrée
% Colonnes extensibles centrée / gauche / droite (C, L, R), souvent utilisées par les modèles
\newcolumntype{C}{>{\centering\arraybackslash}X}
\newcolumntype{L}{>{\raggedright\arraybackslash}X}
\newcolumntype{R}{>{\raggedleft\arraybackslash}X}
\usepackage{array}
\usepackage{multicol}
\usepackage{multirow}
\usepackage{siunitx}
% parse-numbers=false : accepte aussi \SI{2,5 \times 10^{-3}}{...}
\sisetup{locale=FR,output-decimal-marker={,},group-minimum-digits=4,parse-numbers=false}
\DeclareSIUnit\ohm{\ensuremath{\symup{\Omega}}} % U+2126 absent de la police
\DeclareSIUnit\division{div} % réglages d'oscilloscope (ms/div, V/div)
\DeclareSIUnit\tr{tr} % tours (vitesse de rotation, ex. \SI{1500}{\tr\per\minute})
\DeclareSIUnit\molar{M} % concentration molaire (ex. \SI{3}{\molar}, \SI{1}{\micro\molar})
\DeclareSIUnit\an{an} % année (le modèle confond parfois avec \year, un compteur TeX) — ex. \SI{5}{\kilo\meter\per\an}
\DeclareSIUnit\FCFA{FCFA} % franc CFA (ex. \SI{500}{\FCFA\per\kilo\gram})
\DeclareSIUnit\degreeBrix{\textdegree Brix} % teneur en sucre d'un sirop/jus (réfractométrie)
\DeclareSIUnit\month{mois} % le modèle utilise parfois \month, un compteur TeX, comme unité de durée
\usepackage{qrcode}
% \degree : commande fréquemment utilisée par les modèles pour un angle ou une
% température en dehors de \SI{}{} (non fournie par les paquets chargés).
\providecommand{\degree}{^{\circ}}
% \qed (fin de démonstration), utilisé par les modèles sans amsthm
\providecommand{\qed}{\hfill$\square$}
% \barre : double barre des valeurs interdites dans un tableau de variations (tkz-tab)
\providecommand{\barre}{\ensuremath{\|}}
% environnement proof (amsthm non chargé) utilisé par les modèles
\ifdefined\proof\else\newenvironment{proof}[1][Démonstration]{\par\noindent\textit{#1.}\ }{\qed\par}\fi
\usepackage{lastpage}
\usepackage{hyperref}
\hypersetup{hidelinks}

% ============================================================
% Palette « Pop » OnBuch+ — mêmes hex que la classe P (plus_ui.dart)
% ============================================================
\definecolor{poporange}{HTML}{F59321}
\definecolor{popdark}{HTML}{E07A0C}
\definecolor{popcream}{HTML}{FFF6E8}
\definecolor{popink}{HTML}{1C1714}
\definecolor{poppurple}{HTML}{7A5AE0}
\definecolor{popgreen}{HTML}{2E9E6A}
\definecolor{poppink}{HTML}{E0457B}
\definecolor{popblue}{HTML}{2F7DE1}
\definecolor{popgold}{HTML}{C98A12}
\definecolor{popmuted}{HTML}{8A7F76}
\definecolor{popline}{HTML}{EADFD2}
\definecolor{popgray}{HTML}{8A7F76}
\colorlet{popgrey}{popgray}
% Alias tolérants (le modèle nomme parfois les couleurs différemment)
\colorlet{popwhite}{white}
\colorlet{popblack}{popink}
\colorlet{popred}{poppink}
\colorlet{popyellow}{popgold}
% Teintes claires (fonds de tableaux/encadrés) — le modèle nomme parfois
% les couleurs avec un suffixe L (ex. popblueL) sans les avoir définies.
\colorlet{poporangeL}{poporange!20}
\colorlet{popdarkL}{popdark!20}
\colorlet{poppurpleL}{poppurple!20}
\colorlet{popgreenL}{popgreen!20}
\colorlet{poppinkL}{poppink!20}
\colorlet{popblueL}{popblue!20}
\colorlet{popgoldL}{popgold!20}
\colorlet{popredL}{popred!20}
\colorlet{popgrayL}{popgray!20}
% Teintes claires pour fonds de tableaux / zones
\colorlet{popblueL}{popblue!10!white}
\colorlet{poporangeL}{poporange!14!white}
\colorlet{popgreenL}{popgreen!12!white}
\colorlet{poppurpleL}{poppurple!12!white}
\colorlet{poppinkL}{poppink!12!white}
\colorlet{popgoldL}{popgold!14!white}

\pagecolor{popcream}

% ============================================================
% Typographie
% ============================================================
\defaultfontfeatures{Ligatures=TeX}
\setmainfont{PlusJakartaSans-Medium.ttf}[Path=../../../fonts/,BoldFont=PlusJakartaSans-Bold.ttf]
% Symboles, API et emojis absents de la police principale : repli sur DejaVu Sans (caractères actifs)
\newfontface\symfont{DejaVuSans.ttf}[Path=../../../fonts/]
\makeatletter
\newcommand\symactive[1]{\catcode#1=13 \begingroup\lccode`\~=#1 \lowercase{\endgroup\def~}{{\symfont\char#1}}}
\newcommand\symblank[1]{\catcode#1=13 \begingroup\lccode`\~=#1 \lowercase{\endgroup\def~}{}}
\newcommand\symhyph[1]{\catcode#1=13 \begingroup\lccode`\~=#1 \lowercase{\endgroup\def~}{\mbox{-}}}
\newcount\symc
\newcommand\symrange[2]{\symc=#1 \@whilenum\symc<#2 \do{\expandafter\symactive\expandafter{\the\symc}\advance\symc by 1 }}
\newcommand\symblankrange[2]{\symc=#1 \@whilenum\symc<#2 \do{\expandafter\symblank\expandafter{\the\symc}\advance\symc by 1 }}
\makeatother
\symrange{"0250}{"02FF}\symactive{"2713}\symactive{"2714}\symactive{"2717}\symactive{"2718}\symactive{"2610}\symactive{"2611}\symactive{"2612}
\symrange{"2600}{"27BF}
\catcode"2705=13 \begingroup\lccode`\~="2705 \lowercase{\endgroup\def~}{{\symfont\char"2713}}\symblank{"26BD}\symblank{"26A1}
\symhyph{"2011}\symhyph{"2E1A}\symblankrange{"1F300}{"1FAFF}\symblank{"FE0F}
% Maths en sans-serif (Fira Math) avec chiffres et lettres droites de Plus
% Jakarta, pour que les formules se fondent dans le texte.
\usepackage[math-style=ISO,bold-style=ISO]{unicode-math}
\setmathfont{FiraMath-Regular.otf}[Path=../../../fonts/]
\setmathfont{PlusJakartaSans-Medium.ttf}[Path=../../../fonts/,range={up/num,up/{Latin,latin},\mathup}]
\usepackage{icomma} % virgule décimale sans espace en mode math
% Caractères mathématiques Unicode absents des polices de texte (Plus Jakarta,
% Archivo Black) : rendus par leur équivalent mathématique, en texte comme en
% formule — sinon ils s'affichent en carré vide (ex. « Arithmétique dans ℤ »).
\usepackage{newunicodechar}
\newunicodechar{ℝ}{\ensuremath{\mathbb{R}}}
\newunicodechar{ℤ}{\ensuremath{\mathbb{Z}}}
\newunicodechar{ℂ}{\ensuremath{\mathbb{C}}}
\newunicodechar{ℕ}{\ensuremath{\mathbb{N}}}
\newunicodechar{ℚ}{\ensuremath{\mathbb{Q}}}
\newunicodechar{𝔻}{\ensuremath{\mathbb{D}}}
\newunicodechar{α}{\ensuremath{\alpha}}
\newunicodechar{β}{\ensuremath{\beta}}
\newunicodechar{γ}{\ensuremath{\gamma}}
\newunicodechar{δ}{\ensuremath{\delta}}
\newunicodechar{ε}{\ensuremath{\varepsilon}}
\newunicodechar{θ}{\ensuremath{\theta}}
\newunicodechar{λ}{\ensuremath{\lambda}}
\newunicodechar{μ}{\ensuremath{\mu}}
\newunicodechar{σ}{\ensuremath{\sigma}}
\newunicodechar{τ}{\ensuremath{\tau}}
\newunicodechar{φ}{\ensuremath{\varphi}}
\newunicodechar{ψ}{\ensuremath{\psi}}
\newunicodechar{ω}{\ensuremath{\omega}}
\newunicodechar{Σ}{\ensuremath{\Sigma}}
% majuscules produites par \MakeUppercase dans les titres (α-aminés -> Α)
\newunicodechar{Α}{\ensuremath{\alpha}}
\newunicodechar{Β}{\ensuremath{\beta}}
\newunicodechar{Γ}{\ensuremath{\Gamma}}
\newunicodechar{Δ}{\ensuremath{\Delta}}
\newunicodechar{Ε}{\ensuremath{\varepsilon}}
\newunicodechar{Θ}{\ensuremath{\Theta}}
\newunicodechar{Λ}{\ensuremath{\Lambda}}
\newunicodechar{Μ}{\ensuremath{\mu}}
\newunicodechar{Τ}{\ensuremath{\tau}}
\newunicodechar{Φ}{\ensuremath{\Phi}}
\newunicodechar{Ψ}{\ensuremath{\Psi}}
\newunicodechar{Ω}{\ensuremath{\Omega}}
\newunicodechar{↦}{\ensuremath{\mapsto}}
\newunicodechar{∈}{\ensuremath{\in}}
\newunicodechar{∉}{\ensuremath{\notin}}
\newunicodechar{∧}{\ensuremath{\wedge}}
\newunicodechar{⇔}{\ensuremath{\Leftrightarrow}}
\newunicodechar{⟺}{\ensuremath{\Longleftrightarrow}}
\newunicodechar{⇒}{\ensuremath{\Rightarrow}}
\newunicodechar{⟹}{\ensuremath{\Longrightarrow}}
\newunicodechar{∘}{\ensuremath{\circ}}
\newunicodechar{∪}{\ensuremath{\cup}}
\newunicodechar{∩}{\ensuremath{\cap}}
\newunicodechar{≡}{\ensuremath{\equiv}}
\newunicodechar{⊂}{\ensuremath{\subset}}
\newunicodechar{⊥}{\ensuremath{\perp}}
\newunicodechar{∃}{\ensuremath{\exists}}
\newunicodechar{∀}{\ensuremath{\forall}}
\newunicodechar{∛}{\ensuremath{\sqrt[3]{\,}}}
\newunicodechar{∜}{\ensuremath{\sqrt[4]{\,}}}
\newunicodechar{⃗}{\ensuremath{{}^{\to}}}
\newunicodechar{ⁿ}{\ifmmode^{n}\else\textsuperscript{\ensuremath{n}}\fi}
\newunicodechar{ˣ}{\ifmmode^{x}\else\textsuperscript{\ensuremath{x}}\fi}
\newunicodechar{ᵇ}{\ifmmode^{b}\else\textsuperscript{\ensuremath{b}}\fi}
\newunicodechar{ʳ}{\ifmmode^{r}\else\textsuperscript{\ensuremath{r}}\fi}
\newunicodechar{ᵘ}{\ifmmode^{u}\else\textsuperscript{\ensuremath{u}}\fi}
\newunicodechar{ʸ}{\ifmmode^{y}\else\textsuperscript{\ensuremath{y}}\fi}
\newunicodechar{ᵃ}{\ifmmode^{a}\else\textsuperscript{\ensuremath{a}}\fi}
\newunicodechar{ᵉ}{\ifmmode^{e}\else\textsuperscript{\ensuremath{e}}\fi}
\newunicodechar{ᵏ}{\ifmmode^{k}\else\textsuperscript{\ensuremath{k}}\fi}
\newunicodechar{ᵖ}{\ifmmode^{p}\else\textsuperscript{\ensuremath{p}}\fi}
\newunicodechar{ᴺ}{\ifmmode^{N}\else\textsuperscript{\ensuremath{N}}\fi}
\newunicodechar{ᶜ}{\ifmmode^{c}\else\textsuperscript{\ensuremath{c}}\fi}
\newunicodechar{ʰ}{\ifmmode^{h}\else\textsuperscript{\ensuremath{h}}\fi}
\newunicodechar{ᵠ}{\ifmmode^{q}\else\textsuperscript{\ensuremath{q}}\fi}
\newunicodechar{ₙ}{\ifmmode_{n}\else\textsubscript{\ensuremath{n}}\fi}
\newunicodechar{ₐ}{\ifmmode_{a}\else\textsubscript{\ensuremath{a}}\fi}
\newunicodechar{ᵢ}{\ifmmode_{i}\else\textsubscript{\ensuremath{i}}\fi}
\newunicodechar{ᵣ}{\ifmmode_{r}\else\textsubscript{\ensuremath{r}}\fi}
\newunicodechar{ₖ}{\ifmmode_{k}\else\textsubscript{\ensuremath{k}}\fi}
\newunicodechar{ᵦ}{\ifmmode_{\beta}\else\textsubscript{\ensuremath{\beta}}\fi}
\newunicodechar{ₓ}{\ifmmode_{x}\else\textsubscript{\ensuremath{x}}\fi}
\newfontfamily\popdisplay{ArchivoBlack-Regular.ttf}[Path=../../../fonts/]
\newfontfamily\poplogo{SpaceGrotesk-Bold.ttf}[Path=../../../fonts/]
\newfontfamily\popsemi{PlusJakartaSans-SemiBold.ttf}[Path=../../../fonts/]
\newfontfamily\popxbold{PlusJakartaSans-ExtraBold.ttf}[Path=../../../fonts/]
\newfontfamily\popxboldls{PlusJakartaSans-ExtraBold.ttf}[Path=../../../fonts/,LetterSpace=3.0]

\setlist[enumerate]{leftmargin=1.7em,itemsep=5pt,topsep=4pt}
\setlist[itemize]{leftmargin=1.7em,itemsep=4pt,topsep=4pt,label={\color{poporange}\textbullet}}
\setlength{\parindent}{0pt}
\setlength{\parskip}{5pt}
\renewcommand{\arraystretch}{1.25}

% pgfplots : style Pop par défaut
\pgfplotsset{
  every axis/.append style={axis line style={popink,line width=1pt},tick style={popink},
    grid style={popline,line width=0.5pt},label style={font=\small},tick label style={font=\footnotesize}},
  popcurve/.style={poporange,line width=1.8pt,no markers,smooth},
}
% Même style disponible dans un \draw TikZ simple (hors pgfplots)
\tikzset{popcurve/.style={poporange,line width=1.8pt}}
\tikzset{popgrid/.style={popline,very thin}}
\colorlet{popcurve}{poporange} % le nom du style est parfois utilisé comme couleur

% ============================================================
% Pill / squiggle
% ============================================================
\newcommand{\poppill}[3]{%
  \tikz[baseline=-0.55ex]\node[draw=popink,line width=1.2pt,rounded corners=2.6pt,%
    fill=#2,inner xsep=6.5pt,inner ysep=3pt]{%
    \popxboldls\fontsize{7}{7}\selectfont\color{#3}\MakeUppercase{#1}};%
}
\newcommand{\popsquiggle}[1]{%
  \tikz[baseline=-0.4ex]{
    \draw[#1,line width=2.6pt,line cap=round]
      plot [smooth] coordinates {(0,0.09) (0.34,0.01) (0.68,0.21) (1.02,0.01) (1.36,0.10)};
  }%
}
% Mot-clé surligné (marqueur orange sous le texte)
\newcommand{\cle}[1]{{\popxbold\color{popdark}#1}}

% ============================================================
% En-tête / pied
% ============================================================
\pagestyle{fancy}
\fancyhf{}
\renewcommand{\headrulewidth}{0pt}
\renewcommand{\footrulewidth}{0pt}
\lhead{\raisebox{-0.15ex}{\poplogo\fontsize{13}{13}\selectfont\color{popink}OnBuch\color{poporange}+}}
\rhead{\poppill{\DOCMATIERE\ \textbullet\ \DOCNIVEAU\ \textbullet\ Leçon \DOCLECON}{white}{popink}}
\lfoot{\popxbold\fontsize{7.6}{7.6}\selectfont\color{popmuted}\MakeUppercase{Cours signé OnBuch+}\ \textbullet\ {\popsemi\DOCTITRE}}
\rfoot{\tikz[baseline=-0.6ex]\node[rounded corners=8.5pt,fill=popink,minimum height=17pt,minimum width=17pt,inner xsep=5pt,inner sep=0]{\popxbold\fontsize{8}{8}\selectfont\color{popcream}\thepage\,/\,\pageref*{LastPage}};}

% ============================================================
% Titres
% ============================================================
\newcommand{\chaptitre}[2]{%
  \begin{tcolorbox}[enhanced,colback=poporange,colframe=popink,boxrule=2.6pt,arc=11pt,
      drop shadow=popink,left=15pt,right=15pt,top=12pt,bottom=11pt,before skip=3pt,after skip=18pt]
    \poppill{#2}{popink}{popcream}\par\vspace{9pt}
    {\raggedright\hyphenpenalty=10000\exhyphenpenalty=10000\popdisplay\fontsize{22}{24}\selectfont\color{popcream}\MakeUppercase{#1}\par}
  \end{tcolorbox}
}
% Section numérotée : pastille orange avec le numéro + titre Archivo
\newcounter{popsec}
\newcommand{\coursec}[1]{%
  \stepcounter{popsec}\setcounter{popsub}{0}%
  \par\vspace{16pt}\noindent
  \begin{minipage}{\linewidth}
  \tikz[baseline=-0.7ex]\node[draw=popink,line width=1.6pt,fill=poporange,rounded corners=4pt,minimum size=22pt,inner sep=0]{\popdisplay\fontsize{12}{12}\selectfont\color{popcream}\thepopsec};%
  \hspace{8pt}{\popdisplay\fontsize{15}{17}\selectfont\color{popink}\MakeUppercase{#1}}\par
  \vspace{4pt}\noindent{\color{poporange}\rule{36pt}{3.6pt}}
  \end{minipage}\par\nopagebreak\vspace{8pt}
}
\newcounter{popsub}
\newcommand{\courssub}[1]{%
  \stepcounter{popsub}%
  \par\vspace{9pt}\noindent{\popxbold\fontsize{11.5}{13}\selectfont\color{popink}{\color{poporange}\thepopsec.\thepopsub}\hspace{6pt}#1}\par\nopagebreak\vspace{4pt}
}

% ============================================================
% Boîtes pédagogiques — même recette : fond blanc, bordure épaisse colorée,
% ombre dure encre, badge plein. Titre optionnel [#1].
% ============================================================
\tcbset{popbase/.style={breakable,enhanced,colback=white,boxrule=2.3pt,arc=9pt,
  shadow={3pt}{-3pt}{0pt}{popink},left=14pt,right=14pt,top=11pt,bottom=11pt,before skip=11pt,after skip=15pt,
  fontupper=\popsemi\fontsize{10.6}{14.5}\selectfont\color{popink},
  fontlower=\popsemi\fontsize{10.6}{14.5}\selectfont\color{popink}}}
% Coche dessinée (le glyphe ✓ n'existe pas dans les polices du document)
\newcommand{\popcheck}[1]{\tikz[baseline=-0.5ex]\draw[#1,line width=1.6pt,line cap=round,line join=round](-0.13,0.0)--(-0.04,-0.09)--(0.13,0.1);}
\providecommand{\coche}{\popcheck{popgreen}}
\providecommand{\resultat}[1]{\par\smallskip\noindent\fcolorbox{popgreen}{popgreenL}{\parbox{\dimexpr\linewidth-2\fboxsep-2\fboxrule\relax}{\textbf{Résultat.} #1}}\par\smallskip}
\newcommand{\popboxhead}[3]{\poppill{#1}{#2}{white}\ifx\relax#3\relax\else\hspace{6pt}{\popxbold\fontsize{10.6}{12}\selectfont\color{popink}#3}\fi\par\vspace{7pt}}

\newtcolorbox{aretenir}[1][]{popbase,colframe=popgreen,before upper={\popboxhead{À retenir}{popgreen}{#1}}}
% Texte en anglais (document de lecture, dialogue, extrait) : cadre
% d'étude. Un titre optionnel (source, auteur) s'affiche dans l'en-tête.
\newtcolorbox{textebox}[1][]{popbase,colframe=popblue,before upper={\popboxhead{Text}{popblue}{#1}\begin{otherlanguage}{english}},after upper={\end{otherlanguage}}}
% Marques ✓ / ✗ (forme correcte / fautive) souvent utilisées par les modèles
\providecommand{\cross}{\textcolor{poppink}{\ensuremath{\times}}}
\providecommand{\xmark}{\cross}\providecommand{\crossmark}{\cross}\providecommand{\cmark}{\ensuremath{\checkmark}}
% Ligne de réponse / justification à compléter (inventée par les modèles)
\providecommand{\justification}[1]{\par\noindent\textit{Justification~:} \dotfill\par}
% \textipa (package tipa non chargé) : la police ne couvre pas l'API -> simple passage
\providecommand{\textipa}[1]{#1}
% Soulignement (ulem non chargé) : \uline, \ul
\providecommand{\uline}[1]{\underline{#1}}\providecommand{\ul}[1]{\underline{#1}}
% \glossaire{\item ...} inventée par les modèles : liste de glossaire
\providecommand{\glossaire}[1]{\begin{itemize}#1\end{itemize}}
% \alert (beamer) inventée par les modèles : texte mis en évidence
\providecommand{\alert}[1]{\textbf{\textcolor{poppink}{#1}}}
% variantes ASCII d'ordinaux écrites en macro
\providecommand{\iere}{\textsuperscript{re}}\providecommand{\ieme}{\textsuperscript{e}}\providecommand{\ere}{\textsuperscript{re}}\providecommand{\eme}{\textsuperscript{e}}\providecommand{\er}{\textsuperscript{er}}\providecommand{\ie}{\textsuperscript{e}}
% Ordinaux français écrits en macro par les modèles : 1\ière, 3\ième
\newcommand{\ière}{\textsuperscript{re}}\newcommand{\ième}{\textsuperscript{e}}
% \exemple (inventée par les modèles) : étiquette « Exemple : »
\providecommand{\exemple}{\textbf{Exemple~:}\ }
% \square utilisable aussi hors mode math (cases à cocher)
\AtBeginDocument{\let\squareMath\square\renewcommand{\square}{\ensuremath{\squareMath}}}
% Mot ou phrase en anglais à l'intérieur d'un paragraphe français
\newcommand{\eng}[1]{\ifmmode\text{\textit{#1}}\else\begingroup\selectlanguage{english}\itshape #1\endgroup\fi}
\let\esp\eng % alias toléré
% flèches d'intonation inventées par les modèles
\providecommand{\textsearrow}{}\renewcommand{\textsearrow}{\ensuremath{\searrow}}\providecommand{\textnearrow}{}\renewcommand{\textnearrow}{\ensuremath{\nearrow}}
% \fr{..} : mot français (inventé par les modèles)
\providecommand{\fr}[1]{\textit{#1}}
\newtcolorbox{exemplebox}[1][]{popbase,colframe=poporange,before upper={\popboxhead{Exemple}{poporange}{#1}}}
\newtcolorbox{definition}[1][]{popbase,colframe=popblue,before upper={\popboxhead{Définition}{popblue}{#1}}}
\newtcolorbox{propriete}[1][]{popbase,colframe=poppurple,before upper={\popboxhead{Propriété}{poppurple}{#1}}}
\newtcolorbox{methode}[1][]{popbase,colframe=popgold,before upper={\popboxhead{Méthode}{popgold}{#1}}}
\newtcolorbox{attention}[1][]{popbase,colframe=poppink,before upper={\popboxhead{Attention}{poppink}{#1}}}
\newtcolorbox{experience}[1][]{popbase,colframe=popblue,colback=popblueL,before upper={\popboxhead{Expérience}{popblue}{#1}}}
% « Le savais-tu ? » : carte violette pleine (contexte, culture, Cameroun)
\newtcolorbox{savaistu}[1][]{popbase,colback=poppurple,colframe=popink,
  fontupper=\popsemi\fontsize{10.4}{14.2}\selectfont\color{white},
  before upper={\poppill{Le savais-tu\,?}{white}{poppurple}\ifx\relax#1\relax\else\hspace{6pt}{\popxbold\color{white}#1}\fi\par\vspace{7pt}}}
% Exercice résolu : énoncé en haut, \tcblower puis la solution sur fond crème
\newcounter{popexo}\newcounter{popexr}
\newtcolorbox{exoresolu}[1][]{popbase,colframe=poporange,colbacklower=popcream,
  segmentation style={popink,line width=1.2pt,dashed},
  before upper={\stepcounter{popexr}\popboxhead{Exercice résolu \thepopexr}{poporange}{#1}},
  before lower={\poppill{Solution}{popink}{popcream}\par\vspace{6pt}}}
% Exercice (non résolu en place) — la correction est dans la partie Corrigés
\newtcolorbox{exercice}[1][]{popbase,colframe=popink,
  before upper={\stepcounter{popexo}\popboxhead{Exercice \thepopexo}{popink}{#1}}}
% Corrigé
\newtcolorbox{corrige}[1][]{popbase,colframe=popgreen,colback=popgreenL,
  before upper={\popboxhead{Corrigé}{popgreen}{#1}}}
% Objectifs / prérequis / bilan
\newtcolorbox{objectifs}{popbase,colframe=popink,colback=poporangeL,
  before upper={\popboxhead{Objectifs de la leçon}{popink}{}}}
\newtcolorbox{prerequis}{popbase,colframe=popink,colback=popblueL,
  before upper={\popboxhead{Prérequis}{popblue}{}}}
\newtcolorbox{bilan}{popbase,colframe=popink,colback=white,boxrule=2.8pt,
  before upper={\popboxhead{Fiche bilan}{poporange}{L'essentiel à connaître pour l'examen}}}
% Figure encadrée avec légende
\newtcolorbox{popfigure}[1][]{enhanced,breakable=false,colback=white,colframe=popline,boxrule=1.4pt,arc=8pt,
  left=10pt,right=10pt,top=10pt,bottom=8pt,before skip=10pt,after skip=12pt,halign=center,
  lower separated=false,halign lower=center,
  fontlower=\popsemi\fontsize{9}{11.5}\selectfont\color{popmuted},
  #1}
\newcommand{\legende}[1]{\par\vspace{6pt}{\popsemi\fontsize{9}{11.5}\selectfont\color{popmuted}#1}}

% ============================================================
% Signature OnBuch+
% ============================================================
\newcommand{\poplogoinline}[1]{{\poplogo\fontsize{#1}{#1}\selectfont\color{popink}OnBuch\color{poporange}+}}
\newcommand{\signatureonbuch}{%
  \begin{tcolorbox}[enhanced,colback=popink,colframe=popink,boxrule=2.2pt,arc=10pt,
      drop shadow=poporange,left=14pt,right=14pt,top=10pt,bottom=10pt,before skip=0pt,after skip=0pt]
    \tikz[baseline=-0.7ex]\node[circle,fill=popgreen,draw=popcream,line width=1.3pt,minimum size=22pt,inner sep=0]{\popcheck{white}};%
    \hspace{9pt}\begin{minipage}[c]{0.62\linewidth}
      {\popxbold\fontsize{12}{13}\selectfont\color{popcream}Signé\ \textbullet\ {\poplogo OnBuch\color{poporange}+}}\par\vspace{2pt}
      {\raggedright\popsemi\fontsize{8}{10.5}\selectfont\color{popline}\MakeUppercase{Équipe pédagogique OnBuch+}\\\MakeUppercase{Conforme au programme officiel MINESEC}\par}
    \end{minipage}\hfill
    {\poplogo\fontsize{22}{22}\selectfont\color{popcream}OnBuch\color{poporange}+}
  \end{tcolorbox}%
}

% ============================================================
% Page de garde
% #1 titre · #2 matière · #3 niveau · #4 module · #5 n° leçon · #6 durée ·
% #7 description / objectifs (texte ou itemize)
% ============================================================
\newcommand{\pagegarde}[7]{%
  \thispagestyle{empty}
  \newgeometry{a4paper,margin=1.7cm,top=1.6cm,bottom=1.4cm}%
  \begin{tikzpicture}[remember picture,overlay]
    \fill[poporange!18] ([xshift=-3.2cm,yshift=-3.6cm]current page.north east) circle (4.6cm);
    \fill[poppurple!14] ([xshift=2.4cm,yshift=7.6cm]current page.south west) circle (3.8cm);
  \end{tikzpicture}%
  {\poplogo\fontsize{30}{30}\selectfont\color{popink}OnBuch\color{poporange}+}\hfill
  \raisebox{0.6ex}{\poppill{Cours signé \textbullet\ Premium}{popink}{popcream}}\par
  \vspace{1pt}\popsquiggle{poppurple}\par
  \vspace{8pt}
  \begin{tcolorbox}[enhanced,colback=poporange,colframe=popink,boxrule=2.8pt,arc=13pt,
      drop shadow=popink,left=18pt,right=18pt,top=12pt,bottom=13pt,after skip=10pt]
    \poppill{#2 \textbullet\ #3}{popink}{popcream}\hspace{5pt}\poppill{#4}{white}{popink}\par\vspace{8pt}
    {\popdisplay\fontsize{14}{14}\selectfont\color{popink}LEÇON #5}\par\vspace{4pt}
    {\raggedright\hyphenpenalty=10000\exhyphenpenalty=10000\popdisplay\fontsize{25}{27}\selectfont\color{popcream}\MakeUppercase{#1}\par}
    \vspace{8pt}
    {\popxbold\fontsize{9.5}{9.5}\selectfont\color{popink}\MakeUppercase{Durée conseillée : #6}}
  \end{tcolorbox}
  \begin{tcolorbox}[enhanced,colback=white,colframe=popink,boxrule=2.2pt,arc=10pt,
      drop shadow=popink,left=16pt,right=16pt,top=10pt,bottom=10pt,after skip=8pt]
    \poppill{Dans cette leçon}{poppurple}{white}\par\vspace{5pt}
    \popsemi\fontsize{9.3}{12.6}\selectfont\color{popink}#7
  \end{tcolorbox}
  \noindent\begin{minipage}[t]{0.487\linewidth}
    \begin{tcolorbox}[enhanced,colback=popgreen,colframe=popink,boxrule=2.2pt,arc=10pt,
        drop shadow=popink,left=11pt,right=11pt,top=10pt,bottom=10pt,height=3.1cm,valign=center]
      \tcbox[colback=white,colframe=popink,boxrule=1.3pt,arc=5pt,boxsep=0pt,nobeforeafter,
        left=3pt,right=3pt,top=3pt,bottom=3pt,valign=center]{\qrcode[height=1.9cm]{https://play.google.com/store/apps/details?id=com.luvvix.onbuch&hl=fr}}%
      \hspace{8pt}\begin{minipage}[c]{0.5\linewidth}
        {\popdisplay\fontsize{11}{12.5}\selectfont\color{white}\MakeUppercase{Télécharge OnBuch}}\par\vspace{4pt}
        {\raggedright\popsemi\fontsize{8.4}{11}\selectfont\color{white}Gratuit \textbullet\ Google Play\\Tuteur IA, annales, concours, résultats\par}
      \end{minipage}
    \end{tcolorbox}
  \end{minipage}\hfill
  \begin{minipage}[t]{0.487\linewidth}
    \begin{tcolorbox}[enhanced,colback=poppurple,colframe=popink,boxrule=2.2pt,arc=10pt,
        drop shadow=popink,left=11pt,right=11pt,top=10pt,bottom=10pt,height=3.1cm,valign=center]
      \tikz[baseline=-0.7ex]\node[circle,fill=white,draw=popink,line width=1.3pt,minimum size=1.6cm,inner sep=0]{\popdisplay\fontsize{24}{24}\selectfont\color{poppurple}+};%
      \hspace{8pt}\begin{minipage}[c]{0.6\linewidth}
        {\popdisplay\fontsize{11}{12.5}\selectfont\color{white}\MakeUppercase{Passe à OnBuch+}}\par\vspace{4pt}
        {\raggedright\popsemi\fontsize{8.4}{11}\selectfont\color{white}Tous les cours signés, Tuteur IA illimité, fiches PDF — onglet OnBuch+ de l'app\par}
      \end{minipage}
    \end{tcolorbox}
  \end{minipage}
  \vspace{8pt}
  \popcontenu
  \vfill
  \signatureonbuch
  \restoregeometry
  \newpage
}

% Rangée « Ce cours contient » (page de garde)
\newcommand{\poptile}[3]{% #1 couleur #2 gros texte #3 légende
  \begin{tcolorbox}[enhanced,colback=#1,colframe=popink,boxrule=2pt,arc=9pt,shadow={2.5pt}{-2.5pt}{0pt}{popink},
      left=6pt,right=6pt,top=9pt,bottom=9pt,halign=center,valign=center,height=1.75cm,width=\linewidth,nobeforeafter]
    {\popdisplay\fontsize{12.5}{13}\selectfont\color{white}\MakeUppercase{#2}}\par\vspace{3pt}
    {\centering\popsemi\fontsize{7.8}{9.5}\selectfont\color{white}#3\par}
  \end{tcolorbox}}
\newcommand{\popcontenu}{%
  \par\noindent{\popxboldls\fontsize{7.5}{7.5}\selectfont\color{popmuted}CE COURS SIGNÉ CONTIENT}\par\vspace{7pt}
  \noindent\begin{minipage}[t]{0.235\linewidth}\poptile{popblue}{Cours}{complet, illustré, pas à pas}\end{minipage}\hfill
  \begin{minipage}[t]{0.235\linewidth}\poptile{poporange}{Méthodes}{et exercices résolus}\end{minipage}\hfill
  \begin{minipage}[t]{0.235\linewidth}\poptile{poppink}{Exercices}{type examen + corrigés}\end{minipage}\hfill
  \begin{minipage}[t]{0.235\linewidth}\poptile{popgreen}{Fiche bilan}{l'essentiel à retenir}\end{minipage}\par}

% Sommaire
\newtcolorbox{sommaire}{enhanced,colback=white,colframe=popink,boxrule=2.2pt,arc=10pt,
  drop shadow=popink,left=18pt,right=18pt,top=13pt,bottom=13pt,before skip=6pt,after skip=20pt,
  before upper={\poppill{Sommaire}{poppurple}{white}\par\vspace{9pt}\popsemi\fontsize{10.6}{14.5}\selectfont\color{popink}}}

% Signature de fin de cours — poussée tout en bas de la dernière page
\newcommand{\findecours}{%
  \par\vfill\nopagebreak
  \begin{center}\popsquiggle{poporange}\end{center}\vspace{4pt}
  \signatureonbuch
}
\AtBeginDocument{\def\llbracket{\mathopen{[\mkern-3.5mu[}}\def\rrbracket{\mathclose{]\mkern-3.5mu]}}} % intervalles d'entiers (glyphe ⟦ absent de la police)
% Glyphes absents de Fira Math : redéfinis à partir de glyphes disponibles
\AtBeginDocument{%
  \def\setminus{\mathbin{\backslash}}\let\smallsetminus\setminus
  \def\vdots{\mathord{\vbox{\baselineskip3.2pt\lineskiplimit0pt\kern4pt\hbox{.}\hbox{.}\hbox{.}}}}%
  \def\ddots{\mathinner{\mkern1mu\raise6.5pt\hbox{.}\mkern2mu\raise3.6pt\hbox{.}\mkern2mu\raise0.7pt\hbox{.}\mkern1mu}}%
  \def\triangle{\mathord{\scalebox{0.9}{$\symup{\Delta}$}}}%
  \def\nabla{\mathord{\raisebox{\depth}{\scalebox{1}[-1]{$\symup{\Delta}$}}}}%
  \def\checkmark{\popcheck{popdark}}%
  \def\star{\mathbin{\ast}}\def\bigstar{\mathord{\ast}}%
  \def\diamond{\mathbin{\scalebox{0.6}{$\Diamond$}}}%
  \def\bigcup{\mathop{\mathchoice{\raisebox{-0.3em}{\scalebox{1.7}{$\cup$}}}{\scalebox{1.25}{$\cup$}}{\cup}{\cup}}\displaylimits}%
  \def\bigcap{\mathop{\mathchoice{\raisebox{-0.3em}{\scalebox{1.7}{$\cap$}}}{\scalebox{1.25}{$\cap$}}{\cap}{\cap}}\displaylimits}%
  \def\lesseqqgtr{\mathrel{\lessgtr}}\def\bigoplus{\mathop{\oplus}\displaylimits}%
}
\providecommand{\ssi}{si et seulement si} % abréviation fréquente
\AtBeginDocument{\providecommand{\wideparen}{\overparen}} % arc de cercle

%% FIGFIX-BEGIN v5 — correctifs globaux des figures (docs/onbuch-generation/tools/figfix)
\usepackage{adjustbox}\usepackage{etoolbox}\tcbuselibrary{raster}
% toute figure TikZ/pgfplots plus large que la ligne est réduite à la largeur disponible
\BeforeBeginEnvironment{tikzpicture}{\begin{adjustbox}{max width=\linewidth}}
\AfterEndEnvironment{tikzpicture}{\end{adjustbox}}
% légendes pgfplots lisibles par-dessus les courbes
\pgfplotsset{every axis/.append style={legend style={fill=white,fill opacity=0.92,text opacity=1,draw=black!15,inner sep=3pt}}}
% étiquettes d'arêtes (syntaxe edge["..."]) avec fond blanc
\tikzset{every edge quotes/.append style={fill=white,inner sep=1.2pt}}
%% FIGFIX-END
\begin{document}

\pagegarde{Lesson 51 — Vocabulary: Words and expressions related to exploiting service packages (telephone/internet services)}{Anglais}{Tle A}{Chapter 5 : Media and Communication}{EN51}{2 h}{Cette leçon de vocabulaire (Tle A, chapitre 5 : Media and Communication) fournit tout le lexique nécessaire pour comprendre et exploiter les offres de services téléphoniques et internet (forfaits, recharges, données, facturation). Elle insiste sur les collocations, les faux amis et la réutilisation du lexique dans des dialogues et courts textes authentiques.
\vspace{4pt}
\begin{multicols}{2}\small
\begin{itemize}
\item À la fin de cette leçon, je sais identifier et définir le vocabulaire des forfaits téléphoniques et internet (bundle, plan, airtime, data, subscription...)
\item À la fin de cette leçon, je sais employer les collocations correctes avec ce lexique (top up my phone, run out of credit, subscribe to a plan...)
\item À la fin de cette leçon, je sais distinguer les mots de la même famille (subscribe/subscription/subscriber, connect/connection/connectivity...)
\item À la fin de cette leçon, je sais éviter les faux amis et calques du français (consommer, crédit, forfait, actualité...)
\item À la fin de cette leçon, je sais prononcer et orthographier correctement les mots-clés du thème
\item À la fin de cette leçon, je sais comprendre un court texte ou dialogue commercial sur les offres télécoms
\end{itemize}\end{multicols}}

\begin{sommaire}
\begin{multicols}{2}
\begin{enumerate}
\item Champ lexical 1 : Téléphone et appels
\item Champ lexical 2 : Internet et données
\item Champ lexical 3 : Forfaits, abonnements et facturation
\item Faux amis et calques du français
\item Mots de la même famille et formation des mots
\item Le lexique en contexte : dialogue et texte authentiques
\item Activité d'intégration
\item Méthodes et exercices résolus
\item Exercices
\item Corrigés des exercices
\item Fiche bilan
\end{enumerate}
\end{multicols}
\end{sommaire}

\begin{objectifs}
\begin{itemize}
\item À la fin de cette leçon, je sais identifier et définir le vocabulaire des forfaits téléphoniques et internet (bundle, plan, airtime, data, subscription...)
\item À la fin de cette leçon, je sais employer les collocations correctes avec ce lexique (top up my phone, run out of credit, subscribe to a plan...)
\item À la fin de cette leçon, je sais distinguer les mots de la même famille (subscribe/subscription/subscriber, connect/connection/connectivity...)
\item À la fin de cette leçon, je sais éviter les faux amis et calques du français (consommer, crédit, forfait, actualité...)
\item À la fin de cette leçon, je sais prononcer et orthographier correctement les mots-clés du thème
\item À la fin de cette leçon, je sais comprendre un court texte ou dialogue commercial sur les offres télécoms
\item À la fin de cette leçon, je sais comparer deux offres de services et conseiller un client à l'oral et à l'écrit
\item À la fin de cette leçon, je sais rédiger un court paragraphe ou dialogue utilisant au moins huit mots du champ lexical
\end{itemize}
\end{objectifs}

\begin{prerequis}
\begin{itemize}
\item Vocabulaire général des médias et de la communication (Première : phone, call, message, internet, network)
\item Comparatifs et superlatifs (cheaper than, the best offer) pour comparer des forfaits
\item Présent simple et present continuous pour décrire des habitudes de consommation
\item Nombres, prix et unités (CFA francs, megabytes, gigabytes, per month)
\item Structures de conseil (should, you'd better, I recommend...)
\end{itemize}
\end{prerequis}

\newpage

\coursec{Champ lexical 1 : Téléphone et appels}

\courssub{Situation d'accroche et problématique}

Imaginez la scène. Amina, a student in Yaoundé, wants to subscribe to an internet bundle to prepare her Bac research online. At the MTN service centre, the agent speaks of \eng{data bundles}, \eng{airtime}, \eng{top-up}, \eng{rollover} and \eng{fair usage policy} — words she hears every day but cannot explain in English. Meanwhile, her penfriend in London talks about his \eng{monthly plan}, \eng{unlimited calls} and \eng{Wi-Fi hotspot}. How can Amina master this vocabulary to compare offers, avoid being cheated and talk fluently about telephone and internet services?

C'est exactement l'objectif de cette leçon : construire, pas à pas, le vocabulaire des services téléphoniques et internet, en commençant par le champ lexical du \cle{téléphone et des appels}. À la fin de cette section, vous saurez parler de votre téléphone, de votre crédit, de vos appels et de vos problèmes de réseau en anglais naturel et correct.

\begin{textebox}[A phone call problem — reading text]
Amina woke up early on Saturday morning. She wanted to call her uncle in Bamenda to ask him about the documents she needed for her scholarship application. She picked up her mobile phone and dialled his number, but the line was engaged. She tried again five minutes later. This time the phone rang, but nobody answered, so she left a message on his voicemail.

Then she looked at the screen and saw a missed call from her friend Clarisse. She tried to call her back, but a recorded voice said: “You have insufficient credit to make this call.” Amina had run out of airtime. She decided to buy a scratch card at the small shop near her house in Mvog-Ada. The shopkeeper gave her a 1,000 FCFA scratch card. She scratched the silver panel, dialled the recharge code, and topped up her account in less than a minute.

Unfortunately, the network coverage in her neighbourhood was very poor that morning. There was almost no signal, so her call to Clarisse kept dropping. “The network is really bad today,” she complained to her mother. “I will wait until the signal is stronger, or I will send her a text message instead.” Her mother smiled and said: “Why don't you use instant messaging? It is cheaper, and it works even when the signal is weak.” Amina opened her messaging application, typed a short message, and forwarded the photo to Clarisse. Problem solved — without spending a single franc of airtime.
\end{textebox}

\textbf{Glossaire}\par
\begin{itemize}
\item \eng{scholarship application} (n.) : demande de bourse
\item \eng{to dial a number} (v.) : composer un numéro
\item \eng{engaged} (adj.) : occupé (ligne téléphonique)
\item \eng{missed call} (n.) : appel manqué
\item \eng{insufficient credit} (n.) : crédit insuffisant
\item \eng{to run out of airtime} (v.) : être à court de crédit
\item \eng{scratch card} (n.) : carte à gratter (recharge)
\item \eng{to top up} (v.) : recharger (un compte téléphonique)
\item \eng{network coverage} (n.) : couverture réseau
\item \eng{signal} (n.) : signal, réseau
\item \eng{the call keeps dropping} : l'appel n'arrête pas de couper
\item \eng{to forward} (v.) : transférer (un message)
\end{itemize}

\textbf{Questions de compréhension}\par
\begin{enumerate}
\item Why did Amina want to call her uncle?
\item What happened the first time she dialled his number?
\item Why could she not call Clarisse back?
\item How did she top up her phone?
\item What solution did her mother suggest, and why?
\end{enumerate}

\courssub{Le téléphone et ses usages : devices and network}

Observons d'abord les mots qui désignent l'appareil et son environnement technique.

\begin{definition}[Les mots de l'appareil et du réseau]
\begin{itemize}
\item \cle{mobile phone} (GB) / \cle{cell phone} (US) : téléphone portable. Prononciation : « mo-baïl faune » / « sèl faune ».
\item \cle{handset} : le combiné, l'appareil lui-même (terme technique utilisé par les vendeurs et les opérateurs).
\item \cle{SIM card} : carte SIM. \eng{SIM} = \emph{Subscriber Identity Module}.
\item \cle{network provider} / \cle{operator} : opérateur / fournisseur de réseau (MTN, Orange, Nexttel...).
\item \cle{coverage} : couverture (du réseau). Mot invariable, toujours au singulier.
\item \cle{signal} : signal, réseau. On dit \eng{a strong signal} (un bon signal) et \eng{a weak signal} (un signal faible).
\end{itemize}
\end{definition}

\begin{exemplebox}[Exemples traduits]
\eng{I bought a new handset because my old phone was too slow.} « J'ai acheté un nouveau combiné parce que mon vieux téléphone était trop lent. »

\eng{You need a valid SIM card to use this network.} « Il faut une carte SIM valide pour utiliser ce réseau. »

\eng{Which network provider offers the best coverage in your area?} « Quel opérateur offre la meilleure couverture dans ta zone ? »

\eng{There is no signal in the village; we have to climb the hill to make a call.} « Il n'y a pas de réseau au village ; il faut monter la colline pour passer un appel. »
\end{exemplebox}

\begin{savaistu}[Mobile, cell ou hand phone ?]
Au Royaume-Uni, on dit \eng{mobile phone} ; aux États-Unis, on dit \eng{cell phone} (abréviation de \emph{cellular phone}). Au Cameroun anglophone, on entend surtout \eng{mobile phone}, et on parle souvent de \eng{credit} ou \eng{airtime} pour désigner la recharge, comme dans \eng{Please send me some credit!} (« Envoie-moi du crédit, s'il te plaît ! »). En Inde et en Asie du Sud-Est, on entend parfois \eng{hand phone} — à éviter dans un examen camerounais, où \eng{mobile phone} est la forme attendue.
\end{savaistu}

\courssub{Le crédit et la recharge : airtime and top-ups}

\begin{definition}[Airtime]
\cle{Airtime} (mot invariable) : \emph{the amount of money available on your phone to make calls or send messages} — c'est-à-dire le crédit téléphonique. Exemple : \eng{I have no airtime left.} « Je n'ai plus de crédit. »
\end{definition}

\begin{definition}[Le vocabulaire de la recharge]
\begin{itemize}
\item \cle{credit} : le crédit (argent disponible sur le compte téléphonique).
\item \cle{to top up} / \cle{to recharge} : recharger. \eng{top-up} (nom) : une recharge.
\item \cle{scratch card} : carte à gratter.
\item \cle{mobile money transfer} : transfert d'argent mobile (MTN Mobile Money, Orange Money).
\item \cle{recharge code} : code de recharge.
\end{itemize}
\end{definition}

\begin{propriete}[Collocations essentielles avec le crédit]
En anglais, on ne « met » pas du crédit (\emph{put credit} est un calque du français) : on utilise des verbes précis.
\begin{itemize}
\item \eng{to top up my phone} : recharger mon téléphone
\item \eng{to recharge my account} : recharger mon compte
\item \eng{to buy a top-up / a scratch card} : acheter une recharge / une carte à gratter
\item \eng{to run out of airtime / credit} : être à court de crédit
\item \eng{to transfer credit to someone} : transférer du crédit à quelqu'un
\end{itemize}
\end{propriete}

\begin{exemplebox}[Exemples traduits]
\eng{She topped up her phone with 1,000 FCFA.} « Elle a rechargé son téléphone de 1 000 FCFA. »

\eng{I need to top up my phone — I've run out of airtime and I have a missed call from my uncle in Bamenda. The network coverage is poor here, so I'll buy a scratch card downtown.} « Je dois recharger mon téléphone — je n'ai plus de crédit et j'ai un appel manqué de mon oncle à Bamenda. La couverture réseau est mauvaise ici, donc j'achèterai une carte à gratter en ville. »

\eng{My father sent me credit by mobile money transfer.} « Mon père m'a envoyé du crédit par transfert d'argent mobile. »
\end{exemplebox}

\begin{attention}[Credit ou network ? Ne confondez pas !]
Ne dites jamais \eng{I don't have credit} pour dire « je n'ai pas de réseau ». \cle{Credit} signifie l'argent disponible sur votre compte, pas le réseau ! Pour parler du réseau, dites : \eng{I have no signal} ou \eng{there is no network coverage here}.

\begin{center}
\begin{tabularx}{\linewidth}{|X|X|X|}
\hline
\rowcolor{popblueL} Français & English & Remarque \\
\hline
Je n'ai pas de crédit. & I have no credit / I've run out of airtime. & \eng{credit} = argent sur le compte \\
\hline
Je n'ai pas de réseau. & I have no signal / no network coverage. & \eng{signal} / \eng{coverage} = réseau \\
\hline
Il n'y a pas de réseau ici. & There is no coverage here. & \eng{coverage} est invariable \\
\hline
\end{tabularx}
\end{center}
\end{attention}

\courssub{Les appels et les messages : calls and messages}

\begin{definition}[Le vocabulaire des appels]
\begin{itemize}
\item \cle{to make a call} : passer un appel ; \cle{to receive a call} : recevoir un appel.
\item \cle{to dial a number} : composer un numéro. Prononciation : « daïal ».
\item \cle{to pick up (the phone)} : décrocher ; \cle{to hang up} : raccrocher.
\item \cle{missed call} : appel manqué.
\item \cle{voicemail} : messagerie vocale. \eng{to leave a message on someone's voicemail} : laisser un message sur la messagerie de quelqu'un.
\item \cle{call rate} : tarif d'appel.
\item \cle{roaming} : itinérance (utilisation du téléphone à l'étranger). \eng{Roaming charges are very high.} « Les frais d'itinérance sont très élevés. »
\item \cle{engaged} (GB) / \cle{busy} (US) : occupé (ligne).
\end{itemize}
\end{definition}

\begin{definition}[Le vocabulaire des messages]
\begin{itemize}
\item \cle{text message} / \cle{SMS} : SMS, message texte.
\item \cle{instant messaging} : messagerie instantanée (WhatsApp, Telegram...).
\item \cle{to send a message} : envoyer un message ; \cle{to forward a message} : transférer un message ; \cle{to reply to a message} : répondre à un message.
\end{itemize}
\end{definition}

\begin{propriete}[Les verbes à particule du téléphone]
Attention à la place du complément avec les \emph{phrasal verbs} :
\begin{itemize}
\item \eng{She picked up the phone} ou \eng{she picked it up} — jamais \emph{she picked up it}.
\item \eng{He hung up the phone} ou \eng{he hung up on me} (« il m'a raccroché au nez » : \eng{to hang up on someone}).
\end{itemize}
\end{propriete}

\begin{exemplebox}[Exemples traduits]
\eng{The line is engaged.} « La ligne est occupée. »

\eng{I received a missed call from an unknown number.} « J'ai reçu un appel manqué d'un numéro inconnu. »

\eng{Please hold the line; I'll put you through to the manager.} « Veuillez patienter ; je vous passe le directeur. »

\eng{She forwarded my message to the whole class group.} « Elle a transféré mon message à tout le groupe de la classe. »
\end{exemplebox}

\courssub{Les problèmes : when things go wrong}

\begin{definition}[Exprimer les problèmes téléphoniques]
\begin{itemize}
\item \cle{to run out of credit} : être à court de crédit. \eng{Sorry, I couldn't call you — I ran out of credit.}
\item \cle{low battery} : batterie faible. \eng{My battery is low; my phone is about to die.} « Ma batterie est faible ; mon téléphone va s'éteindre. »
\item \cle{no signal} / \cle{poor network} : pas de signal / mauvais réseau.
\item \cle{engaged / busy line} : ligne occupée.
\item \cle{the call drops} : l'appel coupe. \eng{The call keeps dropping.} « L'appel n'arrête pas de couper. »
\item \cle{my phone is off / dead} : mon téléphone est éteint / déchargé.
\end{itemize}
\end{definition}

\begin{center}
\begin{tabularx}{\linewidth}{|X|X|X|}
\hline
\rowcolor{popblueL} Situation & Expression anglaise & Traduction \\
\hline
Plus de crédit & I've run out of airtime. & Je n'ai plus de crédit. \\
\hline
Batterie faible & My battery is low. & Ma batterie est faible. \\
\hline
Pas de réseau & There's no signal here. & Il n'y a pas de réseau ici. \\
\hline
Ligne occupée & The line is engaged. & La ligne est occupée. \\
\hline
Appel qui coupe & The call keeps dropping. & L'appel coupe sans arrêt. \\
\hline
Téléphone déchargé & My phone is dead. & Mon téléphone est déchargé. \\
\hline
\end{tabularx}
\end{center}

\courssub{Carte mentale récapitulative}

\begin{popfigure}
\begin{center}\tcbox[colback=popdark,colframe=popdark,coltext=white,arc=9pt,boxsep=4pt,left=6pt,right=6pt]{\bfseries\small TELEPHONE SERVICES}\end{center}
\vspace{-2pt}
\begin{tcbraster}[raster columns=2,raster equal height=rows,raster column skip=6pt,raster row skip=6pt,enhanced,blankest]
\begin{tcolorbox}[enhanced,colback=white,colframe=popblue,boxrule=0.9pt,arc=3pt,title={Devices},fonttitle=\bfseries\scriptsize,coltitle=white,colbacktitle=popblue,left=4pt,right=4pt,top=2pt,bottom=2pt,boxsep=1pt,toptitle=1.5pt,bottomtitle=1.5pt,halign=left]
\begin{itemize}[nosep,leftmargin=*,itemsep=1pt,font=\scriptsize]
\item SIM card
\item handset
\end{itemize}
\end{tcolorbox}
\begin{tcolorbox}[enhanced,colback=white,colframe=poporange,boxrule=0.9pt,arc=3pt,title={Credit},fonttitle=\bfseries\scriptsize,coltitle=white,colbacktitle=poporange,left=4pt,right=4pt,top=2pt,bottom=2pt,boxsep=1pt,toptitle=1.5pt,bottomtitle=1.5pt,halign=left]
\begin{itemize}[nosep,leftmargin=*,itemsep=1pt,font=\scriptsize]
\item airtime
\item top-up
\end{itemize}
\end{tcolorbox}
\begin{tcolorbox}[enhanced,colback=white,colframe=popgreen,boxrule=0.9pt,arc=3pt,title={Calls},fonttitle=\bfseries\scriptsize,coltitle=white,colbacktitle=popgreen,left=4pt,right=4pt,top=2pt,bottom=2pt,boxsep=1pt,toptitle=1.5pt,bottomtitle=1.5pt,halign=left]
\begin{itemize}[nosep,leftmargin=*,itemsep=1pt,font=\scriptsize]
\item dial
\item voicemail
\end{itemize}
\end{tcolorbox}
\begin{tcolorbox}[enhanced,colback=white,colframe=poppink,boxrule=0.9pt,arc=3pt,title={Problems},fonttitle=\bfseries\scriptsize,coltitle=white,colbacktitle=poppink,left=4pt,right=4pt,top=2pt,bottom=2pt,boxsep=1pt,toptitle=1.5pt,bottomtitle=1.5pt,halign=left]
\begin{itemize}[nosep,leftmargin=*,itemsep=1pt,font=\scriptsize]
\item no signal
\item run out of credit
\end{itemize}
\end{tcolorbox}
\end{tcbraster}

\legende{Carte mentale du champ lexical des services téléphoniques : quatre familles de mots à mémoriser ensemble.}
\end{popfigure}

\courssub{Exercice résolu et entraînement}

\begin{exoresolu}[Complétez avec le mot qui convient]
Complétez chaque phrase avec un mot de la liste : \eng{airtime — engaged — signal — scratch card — voicemail — dial}.

1. I can't call you now; there is no \ldots{} in my quarter. \quad 2. She bought a \ldots{} of 500 FCFA to recharge her phone. \quad 3. I tried his number three times, but the line was \ldots{}. \quad 4. He didn't answer, so I left a message on his \ldots{}. \quad 5. You must \ldots{} the country code before the number. \quad 6. I have run out of \ldots{}; please send me some credit.
\tcblower
\textbf{Corrigé détaillé :}
\begin{enumerate}
\item \textbf{signal} — il s'agit du réseau, pas du crédit (piège classique : \eng{credit} = argent).
\item \textbf{scratch card} — une carte à gratter de 500 FCFA pour recharger.
\item \textbf{engaged} — la ligne est occupée (GB ; US : \eng{busy}).
\item \textbf{voicemail} — laisser un message sur la messagerie vocale.
\item \textbf{dial} — composer l'indicatif du pays avant le numéro.
\item \textbf{airtime} — être à court de crédit : \eng{to run out of airtime}.
\end{enumerate}
\end{exoresolu}

\begin{exercice}[À vous de jouer]
Traduisez en anglais :
\begin{enumerate}
\item Je dois recharger mon téléphone ; je n'ai plus de crédit.
\item Il n'y a pas de réseau dans ce village.
\item Elle a composé le numéro, mais la ligne était occupée.
\item Mon oncle m'a laissé un message sur ma messagerie vocale.
\item La couverture réseau de cet opérateur est excellente à Douala.
\end{enumerate}
\end{exercice}

\begin{corrige}[Exercice — Traduction]
\begin{enumerate}
\item \eng{I need to top up my phone; I've run out of credit / airtime.}
\item \eng{There is no signal / no network coverage in this village.}
\item \eng{She dialled the number, but the line was engaged.}
\item \eng{My uncle left a message on my voicemail.}
\item \eng{This operator's network coverage is excellent in Douala.}
\end{enumerate}
\end{corrige}

\begin{aretenir}[L'essentiel de la section]
\begin{itemize}
\item \eng{mobile phone} (GB) / \eng{cell phone} (US) ; \eng{SIM card}, \eng{handset}, \eng{network provider}, \eng{coverage}, \eng{signal}.
\item \eng{airtime} = crédit téléphonique ; \eng{to top up / recharge} = recharger ; \eng{scratch card} = carte à gratter.
\item \eng{to make / receive a call}, \eng{to dial a number}, \eng{missed call}, \eng{voicemail}, \eng{engaged line}.
\item Piège majeur : \eng{credit} = argent sur le compte ; \eng{signal / coverage} = réseau. Ne les confondez jamais !
\item Collocations à connaître par cœur : \eng{top up my phone}, \eng{run out of airtime}, \eng{pick up / hang up the phone}.
\end{itemize}
\end{aretenir}

\coursec{Champ lexical 2 : Internet et données}

Après avoir exploré le vocabulaire de la téléphonie vocale et des appels (Section 1), nous entrons maintenant dans l'univers du \eng{Internet} et des \eng{data}. Au Cameroun comme partout, l'usage du smartphone a basculé : on « consomme » davantage de données (streaming, réseaux sociaux, téléchargements) que de minutes de communication vocale. Maîtriser ce lexique, c'est pouvoir comprendre son forfait, configurer son téléphone, signaler une panne ou négocier un abonnement — en anglais, langue technique par excellence du numérique.

\courssub{L'accès à Internet : équipement et connexion}

Commençons par observer un court dialogue authentique entre deux élèves de Terminale à Buea, qui tentent de se connecter pour réviser leur bac.

\begin{textebox}{Dialogue : Connexion au lycée}
\textbf{Alice:} \eng{The school Wi-Fi is down again. I can't log in to the portal.} \\
\textbf{Bruno:} \eng{Try the mobile hotspot on my phone. I still have 500 MB left on my bundle.} \\
\textbf{Alice:} \eng{Okay, what's the password?} \\
\textbf{Bruno:} \eng{It's "Bac2025". Connect to "Bruno\_Hotspot". The broadband at home is faster, but this 4G router works fine.} \\
\textbf{Alice:} \eng{Great, I'm online now. Let me download the past papers.}
\end{textebox}

\noindent
\emph{Traduction :} \textbf{Alice :} Le Wi-Fi du lycée ne marche plus. Je n'arrive pas à me connecter au portail. \textbf{Bruno :} Essaie le point d'accès mobile de mon téléphone. Il me reste encore 500 Mo sur mon forfait. \textbf{Alice :} D'accord, c'est quoi le mot de passe ? \textbf{Bruno :} C'est "Bac2025". Connecte-toi à "Bruno\_Hotspot". La fibre à la maison est plus rapide, mais ce routeur 4G fait l'affaire. \textbf{Alice :} Super, je suis en ligne. Je vais télécharger les annales.

\begin{definition}[Vocabulaire clé : L'accès]
\begin{itemize}
    \item \cle{Internet connection} / \cle{connectivity} : connexion Internet (qualité générale).
    \item \cle{Wi-Fi} (ou \cle{WiFi}) : réseau sans fil local. Prononciation : \eng{/waɪ faɪ/} (« waï-faï »).
    \item \cle{Hotspot} : point d'accès Wi-Fi (public ou partage de connexion depuis un téléphone).
    \item \cle{Broadband} : haut débit (connexion fixe rapide : fibre, ADSL, câble).
    \item \cle{Modem} : modem (appareil qui module/démodule le signal de l'opérateur).
    \item \cle{Router} (aussi \cle{Wi-Fi router}) : routeur / box qui distribue le Wi-Fi. Prononciation GB : \eng{/ruːtə/} (« routeur »), US : \eng{/raʊtər/} (« raou-teur »).
    \item \cle{To log in} (ou \cle{log on}) / \cle{to log out} (ou \cle{log off}) : se connecter / se déconnecter (à un compte, un portail, un réseau).
\end{itemize}
\end{definition}

\begin{propriete}[Grammaire : Verbes à particule (Phrasal verbs) — Log in / Log out]
Ces verbes sont \textbf{inséparables} quand l'objet est un nom, mais \textbf{séparables} avec un pronom.
\begin{center}
\begin{tabularx}{\linewidth}{|X|X|X|}
\hline
\rowcolor{popblueL} Forme & Exemple & Traduction \\
\hline
\eng{log in} (intransitif) & \eng{Please log in.} & Connectez-vous. \\
\eng{log in to} + nom & \eng{Log in to the portal.} & Connectez-vous au portail. \\
\eng{log out} (intransitif) & \eng{Don't forget to log out.} & N'oubliez pas de vous déconnecter. \\
\eng{log out of} + nom & \eng{Log out of your account.} & Déconnectez-vous de votre compte. \\
\hline
\end{tabularx}
\end{center}
\emph{Attention :} On dit \eng{log in TO} (pas \eng{into} pour un site/portail), mais \eng{log into} est possible pour un système/ordinateur.
\end{propriete}

\begin{exemplebox}{Exemples d'usage — Accès}
\eng{The \textbf{router} needs a restart; the \textbf{Wi-Fi} signal is weak.} \\
« Le routeur a besoin d'un redémarrage ; le signal Wi-Fi est faible. »

\eng{I use my phone as a \textbf{hotspot} when the \textbf{broadband} fails.} \\
« J'utilise mon téléphone comme point d'accès quand la fibre/l'ADSL coupe. »

\eng{You must \textbf{log in} with your student ID to access the \textbf{connection}.} \\
« Vous devez vous connecter avec votre matricule pour accéder à la connexion. »
\end{exemplebox}

\courssub{Les données : unités, forfaits et actions}

Au Cameroun, on achète des \eng{data bundles} (forfaits data) chez MTN, Orange, Camtel ou Nexttel. Comprendre les unités et les verbes d'action est vital pour ne pas se faire « couper » en plein téléchargement.

\begin{definition}[Vocabulaire clé : Données et volumes]
\begin{itemize}
    \item \cle{Data} (invariable) : données (internet). \textbf{Jamais \eng{datas}.}
    \item \cle{Megabyte (MB)} : mégaoctet (Mo). 1 MB $\approx$ 1 Mo.
    \item \cle{Gigabyte (GB)} : gigaoctet (Go). 1 GB = 1024 MB.
    \item \cle{Data bundle} (ou \cle{data plan}, \cle{internet package}) : forfait data (quantité de data achetée pour une durée donnée).
    \item \cle{To browse} / \cle{to surf the internet} / \cle{to browse the web} : naviguer sur le web.
    \item \cle{To go online} / \cle{to come online} : se connecter / être connecté. Antonyme : \cle{to be offline} / \cle{to go offline}.
    \item \cle{Download} (n/v) : téléchargement / télécharger (réception : serveur $\to$ téléphone).
    \item \cle{Upload} (n/v) : envoi / téléverser / uploader (envoi : téléphone $\to$ serveur).
    \item \cle{Streaming} : lecture en continu (vidéo, musique) sans stockage permanent.
    \item \cle{Social media} (pluriel usuel) : réseaux sociaux (Facebook, WhatsApp, TikTok, X/Twitter).
    \item \cle{App} (abrév. de \cle{application}) : application mobile.
\end{itemize}
\end{definition}

\begin{attention}[Piège majeur : \textbf{Data} et \textbf{Information} sont invariables]
En français, on dit « des données » (pluriel). En anglais courant (surtout télécom), \cle{data} est un \textbf{nom massif (uncountable)} $\to$ verbe au singulier.
\begin{center}
\begin{tabularx}{\linewidth}{|X|X|}
\hline
\rowcolor{poporangeL} \textbf{Correct} & \textbf{Incorrect (calque français)} \\
\hline
\eng{My \textbf{data is} finished.} & \eng{My datas are finished.} \\
\eng{This \textbf{data costs} 500 FCFA.} & \eng{These data cost...} \\
\eng{Do you have \textbf{enough data}?} & \eng{Do you have enough datas?} \\
\hline
\end{tabularx}
\end{center}
\emph{Même règle pour} \cle{information} : \eng{This information is useful} (pas \eng{informations are}).
\end{attention}

\begin{exemplebox}{Collocations clés — Verbes + Data}
\begin{itemize}
    \item \cle{Buy / Purchase a data bundle} : acheter un forfait. \eng{I bought a 2 GB bundle.}
    \item \cle{Use up / Consume / Exhaust one's data} : épuiser son forfait. \eng{He used up all his data in two days.}
    \item \cle{Top up / Recharge one's data} : recharger (souvent via *code\# ou Mobile Money). \eng{I need to top up my data.}
    \item \cle{Share a hotspot / Share one's connection} : partager sa connexion (tethering). \eng{Can you share your hotspot?}
    \item \cle{Connect to Wi-Fi} : se connecter au Wi-Fi. \eng{Connect to the free Wi-Fi at the airport.}
    \item \cle{Browse the internet / the web} : naviguer. \eng{She browses the web for hours.}
\end{itemize}
\end{exemplebox}

\courssub{La qualité de connexion : débit, bande passante, pannes}

\begin{textebox}{Témoignage : La galère du débit à Yaoundé}
\eng{Last night, I tried to download a 500 MB software update. My \textbf{connection} was terribly \textbf{slow} — the \textbf{bandwidth} must have been throttled. The \textbf{download speed} showed 50 kb/s! I gave up and went to bed. This morning, the \textbf{network} is stable, \textbf{connectivity} is perfect, and the \textbf{upload} of my video project took only two minutes. I love it when the \textbf{network} doesn't \textbf{fail} during a Zoom call.}
\end{textebox}

\noindent
\emph{Glossaire :} \textbf{throttled} (bridé, limité intentionnellement par l'opérateur) ; \textbf{kb/s} (kilobits par seconde) ; \textbf{stable} (stable) ; \textbf{project} (projet, devoir).

\begin{definition}[Vocabulaire clé : Qualité et pannes]
\begin{itemize}
    \item \cle{Fast / High-speed connection} : connexion rapide / haut débit.
    \item \cle{Slow / Low-speed connection} : connexion lente / bas débit.
    \item \cle{Bandwidth} : bande passante (capacité max de transfert, en Mbps/Gbps).
    \item \cle{Download speed} / \cle{Upload speed} : débit descendant / montant.
    \item \cle{Network failure} / \cle{Outage} : panne de réseau, coupure.
    \item \cle{Network congestion} : saturation du réseau (heures de pointe).
    \item \cle{Signal strength} / \cle{Signal bars} : force du signal / barres de réseau.
    \item \cle{Latency} / \cle{Ping} : latence / ping (temps de réponse, crucial pour les jeux/video call).
    \item \cle{To buffer} : mettre en tampon (la vidéo s'arrête pour charger). \eng{The video keeps buffering.}
\end{itemize}
\end{definition}

\begin{popfigure}
\begin{tikzpicture}[
    node distance=1.5cm and 2cm,
    box/.style={draw=popblue, thick, rounded corners, fill=popblueL, minimum width=2.2cm, minimum height=1cm, align=center, font=\small},
    arrow/.style={-Stealth, thick, draw=popdark},
    labelarrow/.style={midway, fill=white, inner sep=2pt, font=\tiny, text=popdark}
]
% Nœuds
\node[box] (phone) {Smartphone};
\node[box, right=of phone, align=center] (router){Wi-Fi Router \\ / Mobile Data};
\node[box, right=of router, align=center] (net){Internet \\ (Servers / Cloud)};

% Flèches Download (descendante visuelle : Internet -> Phone)
\draw[arrow, popblue] (net.north) -- ++(0,0.8) -| node[labelarrow, align=center]{\textbf{Download}\\(Recevoir)} (phone.north);
% Flèches Upload (montante visuelle : Phone -> Internet)
\draw[arrow, poporange] (phone.south) -- ++(0,-0.8) -| node[labelarrow, align=center]{\textbf{Upload}\\(Envoyer)} (net.south);

% Légendes centrales
\node[align=center, font=\small, text=popmuted] at ($(router.north)!0.5!(phone.north) + (0,1.2)$) {\textbf{Download} : Video, Apps, Pages Web};
\node[align=center, font=\small, text=popmuted] at ($(router.south)!0.5!(phone.south) - (0,1.2)$) {\textbf{Upload} : Photos, Vidéos, Fichiers, Requêtes};
\end{tikzpicture}
\legende{Schéma : Flux de données — Download (réception) vs Upload (envoi)}
\end{popfigure}

\begin{propriete}[Différence Download vs Upload — Mémo visuel]
\begin{center}
\begin{tabularx}{\linewidth}{|X|X|X|X|}
\hline
\rowcolor{popgreenL} Action & Verbe / Nom & Direction & Exemple typique \\
\hline
\eng{Download} & \eng{to download} / \eng{a download} & Serveur $\to$ Appareil & Regarder Netflix, télécharger une app, recevoir un mail \\
\eng{Upload} & \eng{to upload} / \eng{an upload} & Appareil $\to$ Serveur & Poster sur TikTok, envoyer un mail avec pièce jointe, sauvegarder sur Google Drive \\
\hline
\end{tabularx}
\end{center}
\emph{Astuce :} \eng{Down} = bas (ça descend vers toi) ; \eng{Up} = haut (ça monte vers le nuage).
\end{propriete}

\courssub{Mots de la même famille et formation de mots (Word Formation)}

Le lexique technique s'enrichit par dérivation. C'est un vivier pour l'expression écrite (Writing) et le Use of English.

\begin{propriete}[Familles de mots — Internet \& Data]
\begin{center}
\begin{tabular}{|l|l|l|l|}
\hline
\rowcolor{poppurpleL} \textbf{Verbe} & \textbf{Nom (action/état)} & \textbf{Nom (agent/outil)} & \textbf{Adjectif / Participe} \\
\hline
\eng{connect} & \eng{connection / connectivity} & \eng{connector} & \eng{connected / disconnected} \\
\eng{download} & \eng{download (n)} & \eng{downloader} (rare) & \eng{downloaded / downloading} \\
\eng{upload} & \eng{upload (n)} & \eng{uploader} (rare) & \eng{uploaded / uploading} \\
\eng{browse} & \eng{browsing} & \eng{browser} (navigateur) & \eng{browsable} \\
\eng{stream} & \eng{streaming} & \eng{streamer} & \eng{streamed / streaming} \\
\eng{install} & \eng{installation} & \eng{installer} & \eng{installed} \\
\eng{subscribe} & \eng{subscription} & \eng{subscriber} (abonné) & \eng{subscribed} \\
\eng{configure} & \eng{configuration} & \eng{configurator} & \eng{configured} \\
\hline
\end{tabular}
\end{center}
\emph{Note orthographe :} \eng{browse} (1 seul \eng{w}) $\to$ \eng{browser} (garder le \eng{e} muet devant \eng{r}) ; \eng{router} (pas de \eng{e} final au singulier !).
\end{propriete}

\begin{exemplebox}{Exemples en contexte — Familles de mots}
\eng{A stable \textbf{connection} requires a good \textbf{configuration} of the \textbf{router}.} \\
« Une connexion stable nécessite une bonne configuration du routeur. »

\eng{Heavy \textbf{streaming} \textbf{consumes} your \textbf{data bundle} quickly.} \\
« Le streaming intensif consomme votre forfait data rapidement. »

\eng{The \textbf{subscriber} must \textbf{subscribe} to a new \textbf{subscription} plan.} \\
« L'abonné doit s'abonner à un nouveau plan d'abonnement. »
\end{exemplebox}

\courssub{Prononciation et orthographe : points de vigilance}

\begin{aretenir}[Prononciation (notation française approximative)]
\begin{center}
\begin{tabular}{|l|l|}
\hline
\rowcolor{poppinkL} \textbf{Mot} & \textbf{Prononciation (FR)} \\
\hline
\eng{Wi-Fi} & « waï-faï » \\
\eng{Data} (GB) & « déï-ta » \\
\eng{Data} (US/Cameroun fréquent) & « da-ta » \\
\eng{Router} (GB) & « routeur » \\
\eng{Router} (US) & « raou-teur » \\
\eng{Broadband} & « brode-bande » \\
\eng{Bandwidth} & « bande-uide » \\
\eng{Browse} & « brauz » (1 seul w !) \\
\eng{Browser} & « brau-zeur » \\
\eng{Megabyte / Gigabyte} & « méga-baït / giga-baït » \\
\hline
\end{tabular}
\end{center}
\emph{Conseil :} Au Cameroun, l'anglais standard britannique est la norme scolaire, mais l'américain domine dans la tech. Acceptez les deux, mais soyez cohérent à l'oral (ex. \eng{router} « routeur » en classe).
\end{aretenir}

\begin{attention}[Faux amis et calques fréquents — Spécial Internet]
\begin{center}
\begin{tabularx}{\linewidth}{|X|X|X|}
\hline
\rowcolor{poporangeL} \textbf{Erreur (Français $\to$ Anglais)} & \textbf{Correction} & \textbf{Explication} \\
\hline
\eng{« Je n'ai plus de datas »} & \eng{I have no data left.} / \eng{My data is finished.} & \textbf{Data} = invariable (uncountable). \\
\eng{« Je vais downloader »} (Franglais) & \eng{I'll download it.} & Verbe anglais pur : \textbf{download} (pas \eng{downloader}). \\
\eng{« La connexion coupe »} & \eng{The connection drops / cuts out.} & \eng{To cut} = couper (ciseaux) ; \eng{drop} = tomber, couper (réseau). \\
\eng{« Ouvrir le data »} & \eng{Turn on mobile data / Enable data.} & \eng{Open} = ouvrir (porte, fichier) ; pour activer une fonction : \eng{turn on / enable / activate}. \\
\eng{« Je suis sur le net »} & \eng{I'm online.} / \eng{I'm on the internet.} & \eng{On the net} possible mais \eng{online} est l'adjectif/adverbe standard. \\
\eng{« Un routeur » (orthographe)} & \eng{A router.} & Pas de \eng{e} final en anglais (contrairement au fr \emph{routeur}). \\
\hline
\end{tabularx}
\end{center}
\end{attention}

\courssub{Le lexique en contexte : dialogue et texte authentique}

Voici un texte plus long (style article de blog tech / témoignage d'élève) qui réinvestit l'essentiel du vocabulaire de cette section.

\begin{textebox}{Article : Managing my Data Bundle in Douala}
\eng{By Esther, Terminale A, Douala}

\eng{Last month, my father subscribed to a new \textbf{broadband} plan with Camtel — a \textbf{fibre} connection promising 20 Mbps \textbf{download speed}. At first, the \textbf{connectivity} was amazing. I could \textbf{stream} 4K videos on YouTube without \textbf{buffering}, \textbf{upload} my heavy coding projects to GitHub in seconds, and \textbf{browse} social media smoothly.}

\eng{But two weeks ago, the \textbf{router} started overheating. The \textbf{Wi-Fi} signal became weak in my bedroom. I had to \textbf{log in} to the router's admin page (192.168.1.1) to \textbf{configure} the channel and \textbf{reboot} it. It worked for a while, then a major \textbf{network failure} hit the whole neighbourhood. \textbf{Offline} for 48 hours!}

\eng{During the outage, I relied on my phone's \textbf{mobile data}. I turned on the \textbf{hotspot} to share the connection with my laptop. Bad idea: I \textbf{used up} my 5 GB \textbf{data bundle} in three days just by \textbf{downloading} Windows updates. Now I \textbf{top up} daily with 200 FCFA \textbf{daily bundles}. I've learned to \textbf{monitor} my \textbf{data usage} in the settings and to \textbf{restrict background data} for apps like TikTok. Next payday, I'll ask Dad to \textbf{upgrade} to an \textbf{unlimited} plan!}

\medskip
\noindent
\textbf{Glossaire :}
\begin{itemize}
    \item \eng{subscribe to} : s'abonner à
    \item \eng{fibre} (GB) / \eng{fiber} (US) : fibre optique
    \item \eng{overheating} : surchauffe
    \item \eng{admin page} : page d'administration
    \item \eng{reboot} : redémarrer
    \item \eng{outage} : panne (coupure de service)
    \item \eng{rely on} : compter sur, dépendre de
    \item \eng{monitor} : surveiller, contrôler
    \item \eng{restrict background data} : limiter les données en arrière-plan
    \item \eng{upgrade} : passer à un forfait supérieur
    \item \eng{unlimited plan} : forfait illimité
\end{itemize}
\end{textebox}

\begin{exoresolu}{Compréhension et réemploi lexical (Texte : Esther)}
\textbf{Questions :}
\begin{enumerate}
    \item Quel type de connexion fixe Esther a-t-elle à la maison ? (broadband / fibre)
    \item Pourquoi a-t-elle dû accéder à la page d'administration du routeur ? (overheating, weak signal)
    \item Quel verbe décrit l'action d'épuiser son forfait ? (\eng{use up})
    \item Trouvez dans le texte l'équivalent de : « surveiller sa consommation » (\eng{monitor data usage}), « forfait journalier » (\eng{daily bundle}), « illimité » (\eng{unlimited}).
    \item Traduisez : « I've learned to restrict background data for apps like TikTok. » $\to$ « J'ai appris à limiter les données en arrière-plan pour des applis comme TikTok. »
\end{enumerate}

\tcblower

\textbf{Solution détaillée :}
\begin{enumerate}
    \item \eng{Broadband} (haut débit) via \eng{fibre} (fibre optique).
    \item Le routeur surchauffait (\eng{overheating}) et le signal Wi-Fi était faible (\eng{weak signal}) dans sa chambre.
    \item \eng{Used up} (phrasal verb : épuiser, consumer entièrement). \eng{She used up her 5 GB bundle.}
    \item \eng{Monitor data usage} / \eng{daily bundle} / \eng{unlimited (plan)}.
    \item \eng{Restrict} = limiter/réduire ; \eng{background data} = données en arrière-plan (mises à jour auto, notifications).
\end{enumerate}
\end{exoresolu}

\courssub{Méthode pour mémoriser et réemployer le lexique « Internet \& Data »}

\begin{methode}[Stratégie : « Mon téléphone, mon dictionnaire »]
\begin{enumerate}
    \item \textbf{Changez la langue de votre smartphone en English (UK).} Tous les menus (Settings $\to$ Network \& Internet $\to$ Data usage $\to$ Hotspot \& tethering) deviennent des fiches de vocabulaire vivantes.
    \item \textbf{Associez chaque mot à une action concrète :}
    \begin{itemize}
        \item \eng{Top-up} = la carte à gratter (scratch card) ou le code Mobile Money (*150\#).
        \item \eng{Hotspot} = l'icône « partage de connexion » (deux flèches circulaires).
        \item \eng{Bandwidth} = la jauge de vitesse (Speedtest).
        \item \eng{Buffering} = le cercle qui tourne sur YouTube/Netflix.
    \end{itemize}
    \item \textbf{Créez une « Data Phrase » personnelle} (une phrase qui contient 5 mots cibles). Exemple : \eng{“When my \textbf{data bundle} \textbf{runs out}, I \textbf{top up} via \textbf{Mobile Money} to \textbf{stream} music.”} Récitez-la chaque matin.
    \item \textbf{Écoutez / Lisez en anglais :} Suivez une chaîne tech camerounaise (ex. \eng{TechVoices Cameroon} sur YouTube) ou lisez les FAQ de MTN/Orange Cameroon en version anglaise.
\end{enumerate}
\end{methode}

\begin{experience}[Activité orale : « Help Desk Simulation » (Binôme)]
\textbf{Rôle A (Client) :} Vous êtes à Douala. Votre connexion fibre est lente (\eng{slow}), le routeur clignote en rouge (\eng{red light blinking}). Vous appelez le service client.
\textbf{Rôle B (Agent) :} Vous travaillez pour Camtel/MTN. Posez les questions de diagnostic : \eng{“Have you rebooted the router? What is the download speed? How many devices are connected?”}
\textbf{Objectif :} Utiliser au minimum 8 mots du champ lexical (\eng{bandwidth, connectivity, log in, router, broadband, troubleshoot, upgrade, technician}).
\textbf{Débrief :} Échangez les rôles. Notez les erreurs de prononciation (\eng{router, Wi-Fi, data}) et corrigez-vous mutuellement.
\end{experience}

\begin{savaistu}[Culture numérique au Cameroun : « Le crédit data, nouvelle monnaie »]
Au Cameroun, le \eng{data bundle} est devenu une quasi-monnaie d'échange. On « transfère du data » (\eng{data transfer / data gifting}) via *150\# (MTN) ou *141\# (Orange) pour rendre service, payer un petit service (moto-benskin, courses) ou aider un ami bloqué. Les opérateurs lancent des \eng{night bundles} (forfaits nuit, ex. 1 Go pour 100 FCFA de 23h à 6h) très prisés des étudiants pour les \eng{downloads} lourds. Le \eng{WhatsApp bundle} (forfait WhatsApp illimité) existe aussi, mais attention : il ne couvre pas le \eng{browsing} web ni le \eng{streaming} vidéo ! Le vocabulaire de cette leçon, c'est votre survie numérique quotidienne.
\end{savaistu}

\begin{exercice}[Entraînement : Vocabulaire Internet \& Data]
\begin{enumerate}
    \item \textbf{Complétez} avec le mot juste (choisissez dans la liste) :\\
    \eng{broadband, hotspot, bandwidth, upload, buffer, log out, bundle, router}
    \begin{enumerate}
        \item My \_\_\_\_\_\_\_\_ is almost finished; I need to buy a new 2 GB \_\_\_\_\_\_\_\_.
        \item The \_\_\_\_\_\_\_\_ speed is 50 Mbps, but the \_\_\_\_\_\_\_\_ speed is only 5 Mbps.
        \item Don't forget to \_\_\_\_\_\_\_\_ of your banking app when you finish.
        \item My phone creates a \_\_\_\_\_\_\_\_ so my laptop can connect via Wi-Fi.
        \item The video keeps \_\_\_\_\_\_\_\_ because the \_\_\_\_\_\_\_\_ is too low.
        \item We need to restart the \_\_\_\_\_\_\_\_ to fix the \_\_\_\_\_\_\_\_ connection.
    \end{enumerate}
    \item \textbf{Traduisez en anglais :}
    \begin{enumerate}
        \item J'ai épuisé mon forfait internet en regardant des films en streaming.
        \item La connexion coupe souvent ; il y a une panne de réseau.
        \item Peux-tu partager ta connexion ? Je n'ai plus de data.
        \item Je dois téléverser (uploader) cette vidéo sur YouTube, mais le débit montant est nul.
    \end{enumerate}
    \item \textbf{Corrigez les erreurs (grammaire / vocabulaire) :}
    \begin{enumerate}
        \item I have no more datas on my phone.
        \item Open your data, I want to browse.
        \item The Wi-Fi is very weak, I cannot download the application.
        \item He used up his bundle in two day.
        \item The router is not working, we must reboot him.
    \end{enumerate}
\end{enumerate}
\end{exercice}

\begin{corrige}[Exercice 1]
\textbf{1. Complétion :}
\begin{enumerate}
    \item \eng{data / bundle} (\eng{My \textbf{data} is almost finished; I need to buy a new 2 GB \textbf{bundle}.})
    \item \eng{download / upload} (\eng{The \textbf{download} speed is 50 Mbps, but the \textbf{upload} speed is only 5 Mbps.})
    \item \eng{log out} (\eng{Don't forget to \textbf{log out} of your banking app...})
    \item \eng{hotspot} (\eng{My phone creates a \textbf{hotspot}...})
    \item \eng{buffering / bandwidth} (\eng{The video keeps \textbf{buffering} because the \textbf{bandwidth} is too low.})
    \item \eng{router / broadband} (\eng{We need to restart the \textbf{router} to fix the \textbf{broadband} connection.})
\end{enumerate}

\textbf{2. Traduction :}
\begin{enumerate}
    \item \eng{I used up my data bundle streaming movies.} (ou \eng{by streaming movies})
    \item \eng{The connection keeps dropping / cutting out; there is a network failure / outage.}
    \item \eng{Can you share your hotspot / connection? I've run out of data.}
    \item \eng{I need to upload this video to YouTube, but the upload speed is terrible / zero.}
\end{enumerate}

\textbf{3. Correction d'erreurs :}
\begin{enumerate}
    \item \eng{I have no more \textbf{data} on my phone.} (\eng{data} invariable)
    \item \eng{\textbf{Turn on / Enable} your data, I want to browse.} (\eng{Open} $\to$ \eng{Turn on/Enable} pour une fonction)
    \item \eng{The Wi-Fi is very weak\textbf{, so} I cannot download the \textbf{app}.} (Virgule splice corrigée par \eng{so} ; \eng{application} $\to$ \eng{app} à l'oral/courant)
    \item \eng{He used up his bundle in two \textbf{days}.} (pluriel \eng{days})
    \item \eng{The router is not working, we must reboot \textbf{it}.} (objet inanimé $\to$ \eng{it}, pas \eng{him})
\end{enumerate}
\end{corrige}

\coursec{Champ lexical 3 : Forfaits, abonnements et facturation}

Après avoir exploré le vocabulaire technique de l'internet et des données mobiles (\emph{data, bandwidth, hotspot, tethering}), nous abordons désormais l'aspect commercial et contractuel : comment achète-t-on, gère-t-on et paie-t-on ces services ? Au Cameroun comme ailleurs, le choix entre une carte \cle{prepaid} (prépayée) et un abonnement \cle{postpaid} (postpayé), la lecture d'une \cle{bill} (facture) ou la compréhension d'une \cle{fair usage policy} sont des compétences langagières indispensables pour ne pas se faire piéger par les petites lignes du contrat.

\courssub{Les types d'offres : prépayé, postpayé et formules}

Commençons par observer une publicité typique d'un opérateur camerounais (MTN, Orange, Camtel, Nexttel) pour identifier les grands types d'offres.

\begin{textebox}{Extrait de publicité opérateur mobile — Cameroun}
\eng{“Looking for the best value? Choose your style!”}
\eng{“Option A — \textbf{Prepaid} (Pay-as-you-go): No contract, no monthly bill. Just buy a \textbf{recharge card} or \textbf{top-up} via Mobile Money (MoMo), dial the USSD code, and you are ready to call or surf. Ideal for students and light users.”}
\eng{“Option B — \textbf{Postpaid} (Monthly plan): Get a \textbf{monthly bill} at the end of the cycle. Enjoy \textbf{unlimited calls} to same network and a generous \textbf{data allowance}. Requires a 12-month \textbf{contract} and proof of income. Perfect for professionals.”}
\eng{“Option C — \textbf{Bundle} / \textbf{Package}: Combine \textbf{voice}, \textbf{SMS} and \textbf{data} in one \textbf{promotional package}. Subscribe to the 5\,GB monthly bundle and get \textbf{double data} for the first month!”}
\end{textebox}

\begin{center}
\begin{tabularx}{\linewidth}{|X|X|X|}
\hline
\rowcolor{popblueL} \textbf{English term} & \textbf{Traduction française} & \textbf{Contexte / Collocation} \\
\hline
\cle{prepaid} / \cle{pay-as-you-go} & Prépayé / Carte (recharge à l'usage) & \eng{a prepaid SIM card}, \eng{top up your credit} \\
\cle{postpaid} & Postpayé / Abonnement mensuel & \eng{a postpaid subscription}, \eng{monthly plan} \\
\cle{recharge card} / \cle{top-up voucher} & Carte de recharge / Code de recharge & \eng{buy a 1000 FCFA recharge card} \\
\cle{monthly plan} / \cle{monthly package} & Forfait mensuel / Plan mensuel & \eng{subscribe to a monthly plan} \\
\cle{bundle} / \cle{package} & Pack / Forfait groupé (voix + SMS + data) & \eng{a data bundle}, \eng{a combo package} \\
\cle{unlimited calls / SMS} & Appels / SMS illimités (souvent \og on-net\fg) & \eng{unlimited on-net calls} \\
\cle{allowance} / \cle{quota} & Plafond / Quota inclus & \eng{monthly data allowance}, \eng{exceed your quota} \\
\cle{contract} & Contrat (engagement) & \eng{a 12-month contract}, \cle{no-contract plan} \\
\hline
\end{tabularx}
\end{center}

\begin{popfigure}
\begin{tikzpicture}[
    node distance=0.8cm and 1.5cm,
    box/.style={rectangle, draw=popblue, thick, rounded corners, fill=popblueL, text width=5.5cm, align=center, minimum height=3.5cm},
    title/.style={rectangle, draw=popdark, thick, fill=popgold, text width=5.5cm, align=center, rounded corners=3pt, font=\bfseries},
    arrow/.style={-Stealth, thick, poporange}
]
% Prepaid column
\node[title] (pre-title) {PREPAID (Pay before)};
\node[box, below=of pre-title, align=center] (pre){
    \textbf{Key verbs:} \eng{top up, recharge, buy credit}\\
    \textbf{Bill:} \eng{No monthly bill}\\
    \textbf{Control:} \eng{Full cost control}\\
    \textbf{Commitment:} \eng{No contract}\\
    \textbf{Typical user:} Students, occasional users\\
    \textbf{Cameroon context:} \eng{MoMo top-up, *123\#}
};
% Postpaid column
\node[title, right=of pre-title] (post-title) {POSTPAID (Pay after)};
\node[box, below=of post-title, align=center] (post){
    \textbf{Key verbs:} \eng{subscribe, sign up, pay bill}\\
    \textbf{Bill:} \eng{Monthly invoice / statement}\\
    \textbf{Control:} \eng{Risk of \cle{extra fees} / \cle{out-of-bundle} rates}\\
    \textbf{Commitment:} \eng{Often 12--24 month contract}\\
    \textbf{Typical user:} Professionals, heavy users\\
    \textbf{Perks:} \eng{Unlimited calls, handset financing}
};
% Arrows
\draw[arrow] (pre-title) -- (post-title) node[midway, above, fill=white, inner sep=2pt] {CHOICE};
\end{tikzpicture}
\legende{Comparatif visuel : Prepaid vs Postpaid — Vocabulaire clé et différences structurelles}
\end{popfigure}

\begin{exemplebox}{Exemples traduits — Types d'offres}
\eng{I am on a \textbf{prepaid} plan; I just \textbf{top up} 500 FCFA when I need credit.} \\
« Je suis sur un forfait \textbf{prépayé} ; je recharge juste 500 FCFA quand j'ai besoin de crédit. »

\eng{My father prefers a \textbf{postpaid} \textbf{monthly plan} because he gets an \textbf{itemised bill} for his taxes.} \\
« Mon père préfère un \textbf{forfait mensuel postpayé} car il reçoit une \textbf{facture détaillée} pour ses impôts. »

\eng{The new \textbf{bundle} includes 10\,GB + 2h calls + 100 SMS for 30 days.} \\
« Le nouveau \textbf{pack} comprend 10\,Go + 2h d'appels + 100 SMS pour 30 jours. »

\eng{Be careful: \textbf{unlimited} often means \textbf{unlimited on-net} only (calls to the same network).} \\
« Attention : \textbf{illimité} signifie souvent \textbf{illimité vers le même réseau} seulement. »
\end{exemplebox}

\begin{attention}[Faux ami majeur : \eng{forfeit} $\neq$ \eng{forfait}]
En français, \emph{un forfait} = a fixed-price package.
En anglais, \eng{forfeit} (verbe ou nom) signifie \textbf{perdre un droit, être déclaré forfait (sport), confisquer}.
\eng{The team forfeited the match.} = L'équipe a déclaré forfait / a perdu le match par forfait.
\eng{He forfeited his deposit.} = Il a perdu son dépôt (confisqué).

\cle{Traductions correctes de \og forfait \fg :}
\begin{itemize}
    \item \cle{plan} (général) — \eng{a monthly plan}
    \item \cle{package} (commercial) — \eng{a promotional package}
    \item \cle{bundle} (groupé data/voix/SMS) — \eng{a data bundle}
    \item \cle{package deal} / \cle{combo} (marketing)
\end{itemize}
\eng{Never say: “I bought a new forfeit.” $\rightarrow$ Say: “I bought a new \textbf{bundle / plan}.”}
\end{attention}

\courssub{Formation des mots : familles lexicales (Word Families)}

Le vocabulaire des télécoms repose sur des racines productives. Maîtriser la dérivation (nom $\leftrightarrow$ verbe $\leftrightarrow$ adjectif) permet de décoder tout nouveau terme.

\begin{propriete}[Principales familles de mots du thème]
\begin{center}
\begin{tabular}{|l|l|l|l|}
\hline
\rowcolor{popblueL} \textbf{Verbe} & \textbf{Nom (action/état)} & \textbf{Nom (agent)} & \textbf{Adjectif / Participe} \\
\hline
\cle{subscribe} & \cle{subscription} & \cle{subscriber} & \cle{subscribed} \\
\cle{activate} & \cle{activation} & — & \cle{active} / \cle{activated} \\
\cle{deactivate} & \cle{deactivation} & — & \cle{inactive} / \cle{deactivated} \\
\cle{renew} & \cle{renewal} & — & \cle{renewed} \\
\cle{expire} & \cle{expiry} / \cle{expiration} & — & \cle{expired} \\
\cle{upgrade} / \cle{downgrade} & \cle{upgrade} / \cle{downgrade} & — & \cle{upgraded} / \cle{downgraded} \\
\cle{top up} & \cle{top-up} (n) / \cle{recharge} & — & \cle{topped up} \\
\cle{bill} & \cle{billing} / \cle{invoice} & \cle{biller} (rare) & \cle{billed} \\
\cle{pay} & \cle{payment} & \cle{payer} & \cle{payable} / \cle{paid} \\
\cle{charge} & \cle{charge} / \cle{surcharge} & — & \cle{chargeable} / \cle{charged} \\
\cle{roll over} (phr v) & \cle{rollover} (n/adj) & — & \cle{rolled over} \\
\hline
\end{tabular}
\end{center}
\textbf{Point d'attention :} \eng{expiry} (nom, fréquent en GB) vs \eng{expiration} (nom, plus formel/US). \eng{The expiry date} = \eng{The expiration date}. Le verbe \eng{expire} n'a \textbf{pas} de forme en \eng{-ation}.
\end{propriete}

\courssub{Souscrire et gérer son abonnement : verbes d'action}

La gestion d'une ligne mobile repose sur une série de verbes précis. Observons leurs emplois dans des phrases de la vie courante (dialogue entre deux collègues à Yaoundé).

\begin{textebox}{Dialogue : Gestion d'abonnement (Yaoundé, bureau)}
\eng{— “Did you \cle{subscribe to} the new 10\,GB \textbf{bundle}?”}
\eng{— “Yes, I \textbf{dialled *150\#} to \textbf{activate} it. It was instant.”}
\eng{— “Good. Does it \textbf{auto-renew}?”}
\eng{— “No, I have to \textbf{renew} it manually every 30 days. If I forget, it \textbf{expires} and I lose the remaining \textbf{balance}.”}
\eng{— “I \textbf{upgraded} my plan last month: from 5\,GB to 20\,GB. But if I travel to the village, I might \textbf{downgrade} to a cheaper \textbf{voice-only package}.”}
\eng{— “How do you \textbf{deactivate} a service you don't want?”}
\eng{— “Usually \eng{STOP} to the short code, or call \textbf{customer care}.”}
\end{textebox}

\begin{propriete}[Verbes clés de la gestion d'abonnement — Forme et Emploi]
\begin{center}
\begin{tabular}{|l|l|l|p{6cm}|}
\hline
\rowcolor{popblueL} \textbf{Verbe} & \textbf{Traduction} & \textbf{Construction} & \textbf{Exemple / Contexte} \\
\hline
\cle{subscribe to} & Souscrire à & \eng{subscribe to a plan / a bundle} & \eng{She subscribed to the unlimited package.} \\
\cle{unsubscribe from} & Se désabonner de & \eng{unsubscribe from a service} & \eng{Text STOP to unsubscribe from daily tips.} \\
\cle{sign up for} & S'inscrire à (plus familier) & \eng{sign up for a contract} & \eng{He signed up for a 24-month contract.} \\
\cle{activate} / \cle{deactivate} & Activer / Désactiver & \eng{activate a bundle, deactivate a service} & \eng{Dial *123\# to activate the night bonus.} \\
\cle{renew} & Renouveler & \eng{renew a bundle / a subscription} & \eng{Your plan renews automatically on the 1st.} \\
\cle{expire} & Expirer (intransitif) & \eng{The bundle expires in 2 days.} & \eng{Warning: your data expires at midnight.} \\
\cle{upgrade} & Monter en gamme / Augmenter & \eng{upgrade to a bigger plan} & \eng{I upgraded from 5GB to 20GB.} \\
\cle{downgrade} & Baisser en gamme / Réduire & \eng{downgrade to a cheaper plan} & \eng{If money is tight, downgrade your plan.} \\
\cle{top up} / \cle{recharge} & Recharger (crédit prépayé) & \eng{top up 1000 FCFA via MoMo} & \eng{I need to top up my credit.} \\
\hline
\end{tabular}
\end{center}
\end{propriete}

\begin{aretenir}[Collocations indispensables (mots qui vont ensemble)]
\begin{itemize}
    \item \eng{subscribe to a monthly plan} (pas *\eng{subscribe a plan})
    \item \eng{renew a bundle / a subscription}
    \item \eng{check one's balance} (solde crédit ou data)
    \item \eng{activate / deactivate a service}
    \item \eng{upgrade / downgrade one's plan}
    \item \eng{auto-renewal is on / off} (renouvellement automatique activé/désactivé)
    \item \eng{dial a USSD code} (composer un code USSD, ex. \eng{*123\#})
\end{itemize}
\end{aretenir}

\begin{exoresolu}[Mettre le verbe à la forme correcte (infinitif, présent, prétérit ou participe passé)]
\textbf{Énoncé :} Complétez ce courriel envoyé au service client.
\eng{Dear Customer Care,}
\eng{Last week, I \_\_\_\_\_\_\_ (subscribe) to the 10\,GB bundle, but I \_\_\_\_\_\_\_ (not / receive) the confirmation SMS. I tried to \_\_\_\_\_\_\_ (activate) it again via USSD, but the system said the service \_\_\_\_\_\_\_ (already / activate). Yesterday, I \_\_\_\_\_\_\_ (try) to \_\_\_\_\_\_\_ (renew) it early, but the \_\_\_\_\_\_\_ (expire) date is still in 5 days. I would like to \_\_\_\_\_\_\_ (upgrade) to the 20\,GB plan next month. Please \_\_\_\_\_\_\_ (deactivate) the auto-renewal on the current one.}

\textbf{Solution détaillée :}
\begin{enumerate}
    \item \eng{subscribed} (prétérit, action passée datée \eng{Last week}).
    \item \eng{did not receive} / \eng{didn't receive} (prétérit négatif).
    \item \eng{activate} (infinitif après \eng{tried to}).
    \item \eng{was already activated} / \eng{had already been activated} (passif : le service a été activé). \textbf{Piège :} \eng{activate} est transitif $\rightarrow$ passif nécessaire.
    \item \eng{tried} (prétérit).
    \item \eng{renew} (infinitif après \eng{tried to}).
    \item \eng{expiry} (nom ! \eng{the expiry date} = la date d'expiration. \eng{Expire} est verbe).
    \item \eng{upgrade} (infinitif après \eng{would like to}).
    \item \eng{deactivate} (impératif / infinitif sans \eng{to} après \eng{Please}).
\end{enumerate}
\end{exoresolu}

\courssub{La facturation et le suivi du compte : lexique de la \cle{bill}}

Que l'on soit \cle{prepaid} ou \cle{postpaid}, il faut savoir lire son compte. Voici les termes incontournables.

\begin{definition}[Rollover (Report de forfait)]
\eng{Rollover} (nom ou adjectif) : \textbf{the transfer of unused data, minutes or SMS to the next subscription period}.
Traduction : \textbf{Report du forfait non utilisé} (données, minutes reportées au mois suivant).
\eng{Unused data rolls over to the next month.} = Les données non utilisées sont reportées au mois suivant.
\eng{A rollover plan.} = Un forfait avec report.
\textbf{Attention :} Souvent conditionné par le renouvellement \emph{avant} la date d'expiration (\eng{renew before expiry to benefit from rollover}).
\end{definition}

\begin{center}
\begin{tabularx}{\linewidth}{|X|X|X|}
\hline
\rowcolor{popblueL} \textbf{Terme anglais} & \textbf{Français} & \textbf{Exemple contextuel} \\
\hline
\cle{bill} / \cle{invoice} & Facture (postpayé) & \eng{My monthly bill is 15\,000 FCFA.} \\
\cle{statement} & Relevé de compte & \eng{Check your online statement.} \\
\cle{payment} / \cle{due date} & Paiement / Date d'échéance & \eng{Payment is due by the 5th.} \\
\cle{balance} & Solde (crédit ou data restant) & \eng{Dial *123\# to check your balance.} \\
\cle{expiry date} / \cle{expiration date} & Date d'expiration & \eng{The bundle expiry date is 30/06.} \\
\cle{validity period} & Période de validité & \eng{Validity period: 30 days.} \\
\cle{rollover} & Report (forfait non utilisé) & \eng{Enjoy data rollover on this plan.} \\
\cle{out-of-bundle rate} / \cle{out-of-plan rate} & Tarif hors forfait (souvent très cher) & \eng{Out-of-bundle rate: 50 FCFA/MB.} \\
\cle{extra fees} / \cle{hidden charges} & Frais supplémentaires / cachés & \eng{Beware of hidden charges for roaming.} \\
\cle{airtime} & Crédit d'appel (temps d'antenne) & \eng{Convert airtime to data bundles.} \\
\hline
\end{tabularx}
\end{center}

\begin{exemplebox}{Exemples traduits — Facturation et solde}
\eng{\textbf{Dial *123\# to check your balance.}} \\
« Composez le *123\# pour consulter votre \textbf{solde}. »

\eng{\textbf{Your data bundle expires at midnight.}} \\
« Votre \textbf{forfait data expire à minuit}. »

\eng{\textbf{Unused minutes roll over to the next month if you renew on time.}} \\
« Les minutes non utilisées sont \textbf{reportées au mois suivant} si vous renouvelez à temps. »

\eng{\textbf{Avoid out-of-bundle charges: buy a data pack before you run out.}} \\
« Évitez les \textbf{frais hors forfait} : achetez un pack data avant d'épuiser le vôtre. »

\eng{\textbf{The invoice shows a payment due date of March 5th.}} \\
« La facture indique une \textbf{date d'échéance de paiement} au 5 mars. »
\end{exemplebox}

\courssub{Promotions, bonus et conditions générales : lire les petites lignes}

Les opérateurs rivalisent d'offres promotionnelles. Le vocabulaire marketing est codifié ; le comprendre évite les déceptions.

\begin{textebox}{SMS promotionnel / Notification App — Offre spéciale}
\eng{“🎉 \textbf{Special Offer!} Subscribe to the \textbf{5 GB Monthly Plan} (3\,000 FCFA) and get:”}
\eng{“✅ \textbf{Double Data} (10 GB total) for the 1st month!”}
\eng{“✅ \textbf{Free Night Calls} (10 PM -- 5 AM) to same network!”}
\eng{“✅ \textbf{Free SMS} weekend (Sat--Sun)!”}
\eng{“📅 \textbf{Validity:} 30 days. \textbf{Rollover} enabled: unused data rolls over.”}
\eng{“⚠ \textbf{Terms and Conditions apply.} \textbf{Fair Usage Policy:} Speed reduced after 10 GB/day. \textbf{No hidden charges.} Dial *150*1\# to subscribe.”}
\end{textebox}

\begin{center}
\begin{tabularx}{\linewidth}{|X|X|X|}
\hline
\rowcolor{poppurpleL} \textbf{Terme promotionnel} & \textbf{Sens réel} & \textbf{Traduction / Note} \\
\hline
\cle{special offer} / \cle{promotional package} & Offre promotionnelle / Pack promo & Souvent limité dans le temps (\eng{limited time offer}) \\
\cle{bonus} & Bonus (données, minutes offertes) & \eng{Get a 100\% bonus on recharge.} \\
\cle{double data} & Double data (volume doublé) & \eng{Enjoy double data for 3 months.} \\
\cle{free night calls} / \cle{off-peak calls} & Appels gratuits nuit / heures creuses & Vérifier l'horaire exact (\eng{10pm--5am}) \\
\cle{discount} / \cle{discounted rate} & Réduction / Tarif réduit & \eng{A 20\% discount on the monthly fee.} \\
\cle{terms and conditions (T\&Cs)} & Conditions générales (CGU) & \textbf{Toujours lire !} \eng{Subject to T\&Cs.} \\
\cle{fair usage policy (FUP)} & Politique d'usage raisonnable & Voir \textbf{Le saviez-vous ?} ci-dessous \\
\cle{hidden charges} & Frais cachés & \eng{No hidden charges = pas de frais surprises.} \\
\cle{extra fees} / \cle{surcharges} & Frais supplémentaires & \eng{Roaming incurs extra fees.} \\
\cle{eligibility} / \cle{eligible} & Éligibilité / Éligible & \eng{Offer eligible for new subscribers only.} \\
\hline
\end{tabularx}
\end{center}

\begin{savaistu}[Le piège du \eng{“Unlimited”} : \cle{Fair Usage Policy (FUP)}]
Dans les pays anglophones (UK, USA, Nigeria, Ghana, Afrique du Sud) et de plus en plus au Cameroun, les forfaits \eng{unlimited} (illimités) sont presque toujours soumis à une \textbf{Fair Usage Policy (FUP)}.
\eng{The FUP states that after 20 GB of high-speed data, your speed may be reduced to 256 kbps until the next billing cycle.}
Traduction : « La politique d'usage raisonnable stipule qu'après 20\,Go de data haut débit, votre vitesse peut être réduite à 256 kbps jusqu'au prochain cycle de facturation. »
\textbf{Conséquence :} \eng{Unlimited} $\neq$ \eng{unrestricted} (sans limite technique). C'est « illimité en volume mais pas en vitesse ». En français, on parle de \cle{débit réduit} ou \cle{bridage} après seuil (\cle{throttling} en anglais technique).
\end{savaistu}

\begin{attention}[Calques et erreurs fréquentes des francophones]
\begin{itemize}
    \item \textbf{\og Forfait \fg $\rightarrow$ \eng{forfeit}} : \textbf{FAUX}. (Voir encadré plus haut). Dire \eng{plan, bundle, package}.
    \item \textbf{\og Recharger \fg $\rightarrow$ \eng{recharge} (verbe)} : \textbf{FAUX}. \eng{Recharge} est un \textbf{nom} (\eng{a recharge card}). Verbe = \cle{top up} ou \cle{recharge} (rare, technique). Dire : \eng{I need to \textbf{top up} my credit.}
    \item \textbf{\og Abonnement \fg $\rightarrow$ \eng{abonnement}} : \textbf{FAUX}. Dire \cle{subscription}, \cle{plan}, \cle{contract} (postpaid).
    \item \textbf{\og Solde \fg $\rightarrow$ \eng{sold}} : \textbf{FAUX}. \eng{Sold} = participe passé de \eng{sell} (vendu). Dire \cle{balance} (\eng{check my balance}).
    \item \textbf{\og Valider \fg $\rightarrow$ \eng{validate}} : Possible mais \textbf{lourd}. Pour un code USSD ou un paiement : \cle{confirm}, \cle{enter}, \cle{submit}. \eng{Enter your PIN to confirm.}
    \item \textbf{\og Expirer \fg $\rightarrow$ \eng{expire} (transitif)} : \textbf{FAUX}. \eng{Expire} est \textbf{intransitif}. On ne dit pas \eng{*The system expires my bundle.} Mais \eng{My bundle expires.} / \eng{The system \textbf{terminates} my bundle.}
\end{itemize}
\end{attention}

\courssub{Le lexique en contexte : texte promotionnel complet et exploitation}

Le texte suivant synthétise l'essentiel du vocabulaire de cette section. Il imite une page « Offres » d'un site opérateur ou une notification App.

\begin{textebox}{Texte support : “Choose Your Perfect Plan” (Site opérateur fictif \eng{CamNet})}
\eng{“Welcome to CamNet! Whether you are a student in Buea, a trader in Douala market, or a civil servant in Yaoundé, we have a \textbf{plan} tailored for you.”}

\eng{“\textbf{1. Prepaid Freedom — No Strings Attached}”}
\eng{“Stay in control with our \textbf{Pay-As-You-Go} SIM. No \textbf{monthly bill}, no \textbf{contract}. Simply \textbf{top up} via Mobile Money (MoMo Express) or buy a \textbf{recharge card} at any kiosk. \textbf{Dial *100\#} to \textbf{check your balance} anytime.”}
\eng{“\textit{Perfect for:} Occasional users, visitors, students on a budget.”}

\eng{“\textbf{2. Postpaid Comfort — The Professional Choice}”}
\eng{“Enjoy peace of mind with a fixed \textbf{monthly invoice}. Our \textbf{Signature Plan} (15\,000 FCFA/month) offers \textbf{unlimited on-net calls}, 200 min off-net, 500 SMS and \textbf{25 GB high-speed data}. \textbf{Handset financing} available (pay for your smartphone in 12 \textbf{instalments}).”}
\eng{“\textit{Requires:} Proof of address, ID, 12-month \textbf{commitment}. \textbf{Early termination fees} apply.”}

\eng{“\textbf{3. Smart Bundles — Best Value for Money}”}
\eng{“Why choose? Get \textbf{Voice + SMS + Data} in one \textbf{bundle}.”}
\eng{“\textbf{-- Weekly Social Pack (500 FCFA):} 2 GB + Unlimited WhatsApp/Facebook + 50 SMS. \textbf{Validity:} 7 days.”}
\eng{“\textbf{-- Monthly Mega Pack (5\,000 FCFA):} 15 GB + \textbf{Unlimited calls} (CamNet-to-CamNet) + 200 SMS. \textbf{Validity:} 30 days. \textbf{Rollover:} Unused data rolls over if you \textbf{renew before expiry}.”}
\eng{“\textbf{-- Night Owl Pack (1\,000 FCFA):} 10 GB (10 PM -- 6 AM only). Ideal for downloads \& updates.”}

\eng{“⚡ \textbf{CURRENT PROMOTION: Double Data} on all Monthly Packs for NEW subscribers! \textbf{Free SIM} delivery in Yaoundé/Douala.”}
\eng{“⚠ \textbf{Important Reminders:}”}
\eng{“• \textbf{Fair Usage Policy} applies to “Unlimited” offers (speed reduced after threshold).”}
\eng{“• \textbf{Out-of-bundle rates} are high: 50 FCFA/MB. \textbf{Buy a data add-on} before you run out.”}
\eng{“• \textbf{Roaming} not included. \textbf{Activate roaming pack} before travel.”}
\eng{“• Read full \textbf{Terms \& Conditions} at camnet.cm/terms.”}

\eng{“Need help? \textbf{Call Customer Care} at 111 (free from CamNet) or chat live on the App.”}
\end{textebox}

\begin{center}
\begin{tabularx}{\linewidth}{|l|X|}
\hline
\rowcolor{popgreenL} \textbf{Mot difficile (POS)} & \textbf{Glossaire français} \\
\hline
\cle{tailored for} (adj) & Adapté à, taillé sur mesure pour \\
\cle{no strings attached} (idiom) & Sans engagement, sans condition cachée \\
\cle{kiosk} (n) & Kiosque, petit commerce de rue \\
\cle{invoice} (n) & Facture (synonyme formel de \eng{bill}) \\
\cle{handset financing} (n) & Financement de téléphone (paiement en plusieurs fois) \\
\cle{instalment} (n) & Mensualité, échéance de paiement \\
\cle{early termination fees} (n) & Frais de résiliation anticipée \\
\cle{add-on} (n) & Option supplémentaire (pack data extra) \\
\cle{threshold} (n) & Seuil (déclencheur de la FUP) \\
\cle{run out (of)} (phr v) & Épuiser, manquer de (\eng{run out of data}) \\
\hline
\end{tabularx}
\end{center}

\begin{exoresolu}[Compréhension et réemploi lexical — Texte “Choose Your Perfect Plan”]
\textbf{Énoncé :} Répondez en anglais complet en réemployant le vocabulaire de la leçon.
\begin{enumerate}
    \item What is the main difference between \eng{Prepaid Freedom} and \eng{Postpaid Comfort} regarding billing and contract?
    \item A student in Buea wants \eng{WhatsApp} and \eng{Facebook} cheap for one week. Which \eng{bundle} should she \eng{subscribe to}? How much? Validity?
    \item If a customer on the \eng{Monthly Mega Pack} uses 15 GB in 20 days and does not \eng{renew}, what happens to the unused data? (Use \eng{rollover} / \eng{expire}).
    \item Translate into English: « Le client doit lire les conditions générales car la politique d'usage raisonnable s'applique aux offres illimitées. »
    \item Translate into English: « Je dois recharger mon crédit avant la date d'expiration pour bénéficier du report. »
\end{enumerate}

\textbf{Solution détaillée :}
\begin{enumerate}
    \item \eng{Prepaid Freedom has no monthly bill and no contract (pay-as-you-go), whereas Postpaid Comfort has a fixed monthly invoice and requires a 12-month commitment.}
    \item \eng{She should subscribe to the Weekly Social Pack. It costs 500 FCFA and its validity is 7 days.}
    \item \eng{If he does not renew before the expiry date, the unused data will not roll over; it will expire / be lost. Rollover only works if you renew before expiry.}
    \item \eng{The customer must read the Terms and Conditions because the Fair Usage Policy applies to unlimited offers.}
    \item \eng{I must top up my credit before the expiry date to benefit from the rollover.}
\end{enumerate}
\end{exoresolu}

\begin{methode}[Comment parler de son forfait au téléphone (Speaking / Role-play)]
\textbf{Situation :} Vous appelez le service client (111) ou parlez à un vendeur en boutique.
\begin{enumerate}
    \item \textbf{Identifier sa ligne :} \eng{“Hello, I'm calling about my number \textbf{6XX XXX XXX}. I am on a \textbf{prepaid / postpaid} plan.”}
    \item \textbf{Décrire le problème / besoin :} \eng{“I want to \textbf{subscribe to} the 10 GB bundle. / My data \textbf{expired} yesterday. / I see \textbf{extra fees} on my \textbf{bill} I don't understand.”}
    \item \textbf{Demander une action précise :} \eng{“Can you \textbf{activate} the night bonus? / Please \textbf{deactivate} the auto-renewal. / I want to \textbf{upgrade / downgrade} my plan.”}
    \item \textbf{Vérifier :} \eng{“Can you send me a confirmation SMS? / What is the \textbf{validity period}? / Is there a \textbf{fair usage policy}?”}
    \item \textbf{Clôturer :} \eng{“Thank you. What is the \textbf{reference number} for this request?”}
\end{enumerate}
\end{methode}

\begin{experience}[Jeu de rôle : “Au service client CamNet”]
\textbf{Binômes.} L'un est le \textbf{Client}, l'autre l'\textbf{Agent}.
\textbf{Cartes scénarios (à tirer au sort) :}
\begin{enumerate}
    \item Client : Forfait 5 GB expiré hier. Veut renouveler + activer \eng{double data} promo. Agent : Vérifie éligibilité, explique procédure USSD (*150*1\#), confirme \eng{rollover} si renouvellement aujourd'hui.
    \item Client : Facture postpayée anormale (\eng{hidden charges} pour \eng{roaming} non demandé). Agent : Explique \eng{out-of-bundle} vs \eng{roaming}, propose \eng{downgrade} ou \eng{roaming pack}, note réclamation.
    \item Client : Veut \eng{upgrade} de \eng{Weekly} à \eng{Monthly Mega Pack}. Agent : Explique \eng{validity 30 days}, \eng{rollover}, \eng{FUP}, prix 5000 FCFA, paiement MoMo.
\end{enumerate}
\textbf{Contrainte :} Utiliser au moins 5 mots du tableau « Verbes clés » et 3 termes de « Facturation ».
\end{experience}

\begin{exercice}[Entraînement personnel — Vocabulaire \& Grammaire]
\begin{enumerate}
    \item \textbf{Traduisez en anglais :}
    \begin{itemize}
        \item[a)] Mon forfait prépayé a expiré hier, je dois le renouveler.
        \item[b)] Elle a souscrit à un abonnement postpayé avec engagement de 12 mois.
        \item[c)] Attention aux frais hors forfait si vous dépassez votre quota de données.
        \item[d)] Le report des données non utilisées n'est possible que si vous renouvelez avant la date d'échéance.
    \end{itemize}
    \item \textbf{Complétez avec le mot juste (choisissez dans la liste) :} \eng{allowance, bundle, invoice, rollover, top-up, subscription, fair usage policy, out-of-bundle, handset financing, expiry date.}
    \begin{itemize}
        \item[a)] Your monthly data \_\_\_\_\_\_\_ is 20 GB.
        \item[b)] The \_\_\_\_\_\_\_ for the postpaid plan arrives on the 1st of each month.
        \item[c)] Check the \_\_\_\_\_\_\_ of your bundle to avoid losing data.
        \item[d)] \_\_\_\_\_\_\_ allows you to pay for a phone in monthly instalments.
        \item[e)] Unlimited offers are subject to a \_\_\_\_\_\_\_.
        \item[f)] I need a \_\_\_\_\_\_\_ of 1000 FCFA via MoMo.
        \item[g)] The \_\_\_\_\_\_\_ rate is 50 FCFA per MB.
        \item[h)] Thanks to \_\_\_\_\_\_\_, my unused 2 GB were added to next month's plan.
        \item[i)] The \_\_\_\_\_\_\_ includes voice, SMS and data.
        \item[j)] You can cancel your \_\_\_\_\_\_\_ anytime with a prepaid plan.
    \end{itemize}
    \item \textbf{Corrigez les erreurs dans ces phrases (vocabulaire / grammaire) :}
    \begin{itemize}
        \item[a)] I bought a new forfeit yesterday.
        \item[b)] Please recharge my account with 500 FCFA.
        \item[c)] My balance is 0, I need to sold my credit.
        \item[d)] The system expired my bundle automatically.
        \item[e)] He subscribed a 24-month contract.
    \end{itemize}
\end{enumerate}
\end{exercice}

\begin{corrige}[Exercice Entraînement personnel]
\begin{enumerate}
    \item \textbf{Traduction :}
    \begin{itemize}
        \item[a)] \eng{My prepaid plan expired yesterday, I need to renew it.}
        \item[b)] \eng{She subscribed to a postpaid plan with a 12-month commitment.}
        \item[c)] \eng{Beware of out-of-bundle charges if you exceed your data allowance.}
        \item[d)] \eng{Data rollover is only possible if you renew before the expiry date.}
    \end{itemize}
    \item \textbf{Complétion :}
    \begin{itemize}
        \item[a)] \eng{allowance}
        \item[b)] \eng{invoice} (ou \eng{bill})
        \item[c)] \eng{expiry date}
        \item[d)] \eng{Handset financing}
        \item[e)] \eng{fair usage policy (FUP)}
        \item[f)] \eng{top-up}
        \item[g)] \eng{out-of-bundle}
        \item[h)] \eng{rollover}
        \item[i)] \eng{bundle} (ou \eng{package})
        \item[j)] \eng{subscription} (pas de \eng{contract} en prépayé)
    \end{itemize}
    \item \textbf{Correction :}
    \begin{itemize}
        \item[a)] \eng{I bought a new \textbf{bundle / plan} yesterday.} (\eng{forfeit} = faux ami)
        \item[b)] \eng{Please \textbf{top up} my account with 500 FCFA.} (\eng{recharge} = nom)
        \item[c)] \eng{My \textbf{balance} is 0, I need to \textbf{top up} my credit.} (\eng{sold} = vendu)
        \item[d)] \eng{My bundle \textbf{expired} automatically.} (\eng{expire} intransitif) ou \eng{The system \textbf{terminated} my bundle.}
        \item[e)] \eng{He \textbf{signed up for} a 24-month contract.} (\eng{subscribe to} ou \eng{sign up for} + préposition)
    \end{itemize}
\end{enumerate}
\end{corrige}

\coursec{Faux amis et calques du français}

Après avoir passé en revue le vocabulaire technique des forfaits, de la facturation et des abonnements (\emph{« Champ lexical 3 : Forfaits, abonnements et facturation »}), nous devons maintenant nous prémunir contre l'écueil numéro un des apprenants francophones : l'interférence de la langue maternelle. Le passage du français à l'anglais ne se fait pas mot à mot ; il repose sur des \cle{collocations} (associations de mots figées) et des \cle{faux amis} qui, s'ils sont ignorés, trahissent immédiatement un locuteur non natif et peuvent entraîner des malentendus coûteux — surtout quand il s'agit de recharger son crédit ou de contester une facture !

\courssub{Les faux amis majeurs du domaine télécoms}

Commençons par observer trois paires de mots qui piègent systématiquement les élèves de Terminale, même les meilleurs. Lisez les phrases suivantes et repérez l'erreur dans la traduction littérale.

\begin{textebox}{Erreurs typiques d'élèves francophones}
\textbf{Élève A :} \eng{I have consumed my forfait this month.} \\
\textbf{Élève B :} \eng{I need to buy a new credit for my phone.} \\
\textbf{Élève C :} \eng{The forfeit is too expensive for me.}
\end{textebox}

\begin{attention}[Erreur fréquente n°1 des candidats francophones : \eng{« I consumed my forfait »} — double faute !]
Ne dites \textbf{jamais} \eng{I consumed my forfait}. C'est une double erreur :
\begin{enumerate}
    \item \eng{Forfait} n'existe pas en anglais avec ce sens ; on dit \cle{bundle}, \cle{plan}, \cle{package} ou \cle{deal}.
    \item \eng{To consume} ne s'emploie pas pour le crédit ou les données téléphoniques (voir \eng{to use up} / \eng{to run out of} ci-dessous).
\end{enumerate}
\emph{Correction :} \eng{I have used up my bundle / I've run out of data.}
\end{attention}

\begin{definition}[Faux amis : \eng{Forfeit} vs \eng{Bundle/Plan} ; \eng{Consume} vs \eng{Use up} ; \eng{Credit} vs \eng{Airtime}]
\begin{itemize}
    \item \textbf{\eng{Forfeit} (nom/verbe) :} Pénalité, amende, perte d'un droit (sport, droit, jeu). Prononciation : « for-fit ». \eng{The team had to forfeit the match.} (L'équipe a dû déclarer forfait.) \textbf{Jamais} « forfait téléphonique ».
    \item \textbf{\eng{Bundle} (nom) :} Forfait, pack, offre groupée (data + appels + SMS). Prononciation : « ban-dle ». \eng{A monthly bundle of 10 GB.} (Un forfait mensuel de 10 Go.)
    \item \textbf{\eng{To consume} (verbe) :} Consommer (nourriture, énergie, ressources naturelles, biens). Prononciation : « kon-syoum ». \eng{This engine consumes too much fuel.} (Ce moteur consomme trop de carburant.) \textbf{Pas} pour le crédit/data.
    \item \textbf{\eng{To use up} (verbe à particule) :} Épuiser, utiliser entièrement (stock, crédit, data, temps). Prononciation : « youz ape ». \eng{I've used up all my data.} (J'ai épuisé toutes mes données.)
    \item \textbf{\eng{Credit} (nom) :} Crédit (bancaire, reconnaissance, note à l'école). Prononciation : « kre-dit ». En téléphonie, on comprendra mais le terme naturel est \cle{airtime} (temps d'appel) ou \cle{balance} (solde).
    \item \textbf{\eng{Airtime} (nom) :} Crédit d'appel, temps de communication acheté (recharge). Prononciation : « ère-taïm ». \eng{I need to buy airtime.} (Il me faut acheter du crédit.)
\end{itemize}
\end{definition}

\begin{center}
\begin{tabularx}{\linewidth}{|X|X|X|}
\hline
\rowcolor{popblueL} \textbf{Français (source de l'erreur)} & \textbf{Anglais incorrect (calque/faux ami)} & \textbf{Anglais correct (collocation native)} \\ \hline
Un forfait (téléphone) & \eng{A forfeit} & \eng{A bundle / A plan / A package / A deal} \\ \hline
Consommer (son crédit/data) & \eng{To consume} & \eng{To use up / To run out of} \\ \hline
Recharger / Acheter du crédit & \eng{To buy credit} & \eng{To top up / To buy airtime / To recharge} \\ \hline
Le solde / Le crédit restant & \eng{The credit} & \eng{The balance / My airtime} \\ \hline
\end{tabularx}
\end{center}

\courssub{Autres faux amis fréquents (anglais général)}

Au-delà du vocabulaire technique, certains faux amis « classiques » reviennent sans cesse dans les copies du Baccalauréat quand les élèves traitent de sujets de société ou de technologie. Maîtrisez-les une fois pour toutes.

\begin{exemplebox}[Trois faux amis « tueurs » au Bac]
\begin{itemize}
    \item \textbf{\eng{Actual} (adjectif) = \textbf{Réel, effectif, véritable} (pas « actuel »).} Prononciation : « ak-tchou-alle ». \\
    \eng{The actual cost is higher than the estimate.} (Le coût \textbf{réel} est plus élevé que l'estimation.) \\
    \rightarrow \textit{Pour « actuel » :} \cle{Current} / \cle{Present}. \eng{My current plan expires next week.} (Mon forfait \textbf{actuel} expire la semaine prochaine.)
    \item \textbf{\eng{Eventually} (adverbe) = \textbf{Finalement, à la fin, après un certain temps} (pas « éventuellement »).} Prononciation : « i-ven-tchou-el-man ». \\
    \eng{If you don't top up, you will eventually lose your number.} (Si tu ne recharges pas, tu \textbf{finiras par} / \textbf{finalement} perdre ton numéro.) \\
    \rightarrow \textit{Pour « éventuellement » :} \cle{Possibly} / \cle{Maybe} / \cle{If necessary}. \eng{We can possibly upgrade the plan.} (Nous pouvons \textbf{éventuellement} améliorer le forfait.)
    \item \textbf{\eng{A lecture} (nom) = \textbf{Une conférence, un cours magistral} (pas « une lecture »).} Prononciation : « lek-tcheur ». \\
    \eng{He gave a lecture on 5G technology.} (Il a fait une \textbf{conférence} sur la 5G.) \\
    \rightarrow \textit{Pour « une lecture » :} \cle{Reading} / \cle{A read}. \eng{It's a good read.} (C'est une bonne \textbf{lecture}.)
\end{itemize}
\end{exemplebox}

\courssub{Les calques syntaxicaux et lexicaux à corriger}

Un \cle{calque} est une traduction mot à mot de la structure française qui produit une phrase grammaticalement possible en anglais mais \textbf{non idiomatique} (ça sonne « faux », « traduit »). Voici les quatre calques les plus courants dans le thème des télécoms.

\begin{propriete}[Quatre calques à bannir et leurs corrections]
\begin{enumerate}
    \item \textbf{Calque :} \eng{*I want to subscribe a bundle.} (Structure FR : \textit{s'abonner à} qqch $\rightarrow$ transitive directe) \\
    \textbf{Règle :} \eng{Subscribe} est intransitif $\rightarrow$ \eng{Subscribe \textbf{TO} something}. \\
    \textbf{Correct :} \eng{I want to \textbf{subscribe to} a bundle.} / \eng{I want to \textbf{sign up for} a plan.}
    \item \textbf{Calque :} \eng{*The cost is expensive.} (Structure FR : \textit{Le coût est cher}) \\
    \textbf{Règle :} En anglais, un \eng{cost} / \eng{price} / \eng{fee} est \cle{high} / \cle{low} / \cle{reasonable}. Une \eng{plan} / \eng{phone} / \eng{service} est \cle{expensive} / \cle{cheap} / \cle{affordable}. \\
    \textbf{Correct :} \eng{The \textbf{price is high} / The \textbf{plan is expensive}.}
    \item \textbf{Calque :} \eng{*I have consumed my data.} (Structure FR : \textit{J'ai consommé mes données}) \\
    \textbf{Règle :} \eng{Consume} $\neq$ utiliser du data. Verbe correct : \eng{Use up} (accompli) ou \eng{Run out of} (état résultant). \\
    \textbf{Correct :} \eng{I have \textbf{used up} my data.} / \eng{I've \textbf{run out of} data.}
    \item \textbf{Calque :} \eng{*My phone has no money.} (Structure FR : \textit{Mon téléphone n'a pas d'argent}) \\
    \textbf{Règle :} Un téléphone ne « possède » pas d'argent. On parle de \eng{Credit}, \eng{Airtime}, \eng{Balance}. Verbe idiomatique : \eng{Run out of}. \\
    \textbf{Correct :} \eng{I've \textbf{run out of credit} / \textbf{airtime}.} / \eng{My \textbf{balance is zero}.}
\end{enumerate}
\end{propriete}

\begin{methode}[Stratégie pour éviter les calques : la « traduction mentale inversée »]
Avant d'écrire ou de dire une phrase en anglais :
\begin{enumerate}
    \item Écrivez votre phrase en français.
    \item Traduisez-la \textbf{littéralement} mot à mot en anglais (sur brouillon).
    \item Lisez la version littérale : \textbf{sonne-t-elle naturel ?} (Est-ce que \eng{« The cost is expensive »} sonne comme une phrase qu'un Anglais dirait ? Non, un Anglais dit \eng{« The price is high »}).
    \item Si ça sonne bizarre, \textbf{cherchez la collocation native} (le « bloc » de mots qui va ensemble) dans votre mémoire ou votre dictionnaire : \eng{High price, Heavy rain, Fast food, Run out of, Subscribe to}.
    \item Écrivez la version idiomatique.
\end{enumerate}
\emph{Exemple :} FR : « Je veux m'abonner à ce forfait. » $\rightarrow$ Littéral : \eng{*I want to subscribe this bundle.} $\rightarrow$ Bizarre ? Oui. $\rightarrow$ Collocation : \eng{Subscribe TO / Sign up FOR.} $\rightarrow$ Final : \eng{I want to sign up for this bundle.}
\end{methode}

\begin{popfigure}
\begin{tikzpicture}[
    node distance=0.5cm and 1cm,
    box/.style={rectangle, draw=popblue, thick, rounded corners=4pt, minimum height=1.2cm, text width=7.5cm, align=center, font=\small, fill=popblueL},
    arrow/.style={-Stealth, thick, poporange},
    title/.style={font=\bfseries\large, text=popdark, fill=popcream, draw=poporange, rounded corners=4pt, minimum height=1cm, text width=7.5cm, align=center}
]
\node[title] (t1) {\textbf{Ce que le francophone dit (FAUX / NON IDIOMATIQUE)}};
\node[title] (t2) [right=of t1] {\textbf{Ce qu'il faut dire (CORRECT / NATIF)}};

\node[box, align=center] (f1) [below=of t1]{\eng{*I want to subscribe a new bundle.}\\(*Je veux m'abonner un forfait.)};
\node[box, align=center] (c1) [below=of t2]{\eng{I want to \textbf{subscribe to} / \textbf{sign up for} a new bundle.}\\(Je veux m'abonner \textbf{à} un nouveau forfait.)};

\node[box, align=center] (f2) [below=of f1]{\eng{*The cost of this plan is expensive.}\\(*Le coût de ce forfait est cher.)};
\node[box, align=center] (c2) [below=of c1]{\eng{The \textbf{price} of this plan is \textbf{high}.} / \eng{This plan \textbf{is expensive}.}\\(Le prix de ce forfait est \textbf{élevé} / Ce forfait \textbf{est cher}.)};

\node[box, align=center] (f3) [below=of f2]{\eng{*I have consumed all my data.}\\(*J'ai consommé toutes mes données.)};
\node[box, align=center] (c3) [below=of c2]{\eng{I have \textbf{used up} all my data.} / \eng{I've \textbf{run out of} data.}\\(J'ai \textbf{épuisé} toutes mes données / Je n'ai \textbf{plus} de données.)};

\node[box, align=center] (f4) [below=of f3]{\eng{*My phone has no money.}\\(*Mon téléphone n'a pas d'argent.)};
\node[box, align=center] (c4) [below=of c3]{\eng{I've \textbf{run out of credit / airtime}.} / \eng{My \textbf{balance is zero}.}\\(J'ai \textbf{épuisé mon crédit} / Mon \textbf{solde est à zéro}.)};

\draw[arrow] (f1.east) -- (c1.west) node[midway, above, font=\tiny, text=popdark] {Correction};
\draw[arrow] (f2.east) -- (c2.west) node[midway, above, font=\tiny, text=popdark] {Correction};
\draw[arrow] (f3.east) -- (c3.west) node[midway, above, font=\tiny, text=popdark] {Correction};
\draw[arrow] (f4.east) -- (c4.west) node[midway, above, font=\tiny, text=popdark] {Correction};
\end{tikzpicture}
\legende{Tableau de correction des calques syntaxicaux et lexicaux les plus fréquents (Français $\rightarrow$ Anglais).}
\end{popfigure}

\courssub{Verbes à particule (Phrasal Verbs) indispensables}

Dans le domaine des télécoms, l'anglais privilégie massivement les \cle{verbes à particule} (phrasal verbs). La particule (\eng{up, off, out of, through}) change radicalement le sens du verbe de base. Apprenez-les comme des \textbf{blocs lexicaux insécables} (lexical chunks), non comme verbe + préposition.

\begin{propriete}[Les 5 phrasal verbs « télécoms » à connaître par cœur]
\begin{center}
\begin{tabular}{|l|l|l|l|}
\hline
\rowcolor{popblueL} \textbf{Phrasal Verb} & \textbf{Sens} & \textbf{Exemple (EN)} & \textbf{Traduction (FR)} \\ \hline
\eng{Top up} & Recharger (crédit, data) & \eng{I need to \textbf{top up} my phone.} & Il me faut \textbf{recharger} mon téléphone. \\ \hline
\eng{Run out of} & Ne plus avoir, épuiser (stock) & \eng{We've \textbf{run out of} airtime.} & Nous n'avons \textbf{plus de} crédit. \\ \hline
\eng{Use up} & Épuiser, utiliser entièrement & \eng{Don't \textbf{use up} all your data today.} & N'\textbf{épuise} pas toutes tes données aujourd'hui. \\ \hline
\eng{Cut off} & Couper (communication, électricité) & \eng{The call got \textbf{cut off}.} & L'appel a été \textbf{coupé}. \\ \hline
\eng{Get through (to)} & Réussir à joindre (qn au tel) & \eng{I can't \textbf{get through to} him.} & Je n'arrive pas à \textbf{le joindre}. \\ \hline
\end{tabular}
\end{center}
\emph{Note de grammaire :} \eng{Run out of} et \eng{Get through to} prennent \textbf{trois} mots (verbe + particule + préposition) et sont suivis d'un nom/pronom. \eng{Top up} et \eng{Use up} sont \textbf{séparables} : \eng{Top up my phone} = \eng{Top my phone up} (mais \eng{Top it up}, jamais \eng{Top up it}).
\end{propriete}

\courssub{Expressions idiomatiques utiles}

Ces expressions figées (« idioms ») permettent de parler de qualité de service, de prix et de problèmes de connexion avec une aisance native.

\begin{exemplebox}[Expressions idiomatiques : Qualité, Prix, Connexion]
\begin{itemize}
    \item \textbf{\eng{Value for money} (nom) :} Rapport qualité-prix. Prononciation : « va-lyou feu ma-ni ». \\
    \eng{This unlimited bundle is excellent \textbf{value for money}.} = Ce forfait illimité offre un \textbf{excellent rapport qualité-prix}.
    \item \textbf{\eng{To get through (to someone)} :} Réussir à joindre quelqu'un (au téléphone). \\
    \eng{I've been trying all morning but I can't \textbf{get through to} the customer service.} = J'essaie depuis ce matin mais je n'arrive pas à \textbf{joindre} le service client.
    \item \textbf{\eng{To be cut off} (passif) / \eng{To get cut off} :} Être coupé (communication). \\
    \eng{We were \textbf{cut off} in the middle of the conversation.} = Nous avons été \textbf{coupés} en pleine conversation.
    \item \textbf{\eng{To have no signal / reception} :} Ne pas avoir de réseau. \\
    \eng{I have \textbf{no signal} in this building.} = Je n'ai \textbf{pas de réseau} dans ce bâtiment.
    \item \textbf{\eng{To buffer} (verbe) :} Mettre en tampon (vidéo, streaming qui saccade). \\
    \eng{The video keeps \textbf{buffering} because the connection is slow.} = La vidéo \textbf{saccade / charge} car la connexion est lente.
\end{itemize}
\end{exemplebox}

\begin{savaistu}[Le saviez-vous ? \eng{To consume} existe bien, mais...]
Le verbe \eng{to consume} existe bel et bien en anglais, mais son \cle{champ sémantique} est restreint :
\begin{enumerate}
    \item \textbf{Nourriture/Boisson :} \eng{He consumed a large meal.} (Il a mangé un gros repas.)
    \item \textbf{Énergie/Ressources :} \eng{This car consumes a lot of petrol.} (Cette voiture consomme beaucoup d'essence.)
    \item \textbf{Médias/Contenus (formel) :} \eng{Young people consume media differently.} (Les jeunes consomment les médias différemment.)
\end{enumerate}
\textbf{JAMAIS} pour « consommer son forfait / son crédit / ses données ». C'est une erreur de \cle{collocation} (mauvais mariage de mots), pas une erreur de sens du verbe. L'anglais voit le forfait comme un \textbf{stock} qu'on \textbf{épuise} (\eng{use up, run out of}), pas comme une substance qu'on \textbf{consomme} (\eng{consume}).
\end{savaistu}

\courssub{Mots de la même famille et formation des mots}

Pour enrichir votre lexique et varier vos formulations à l'écrit comme à l'oral, observez les familles de mots dérivés des termes clés du thème. La dérivation (préfixes, suffixes) change la catégorie grammaticale.

\begin{propriete}[Familles de mots clés du thème « Service Packages »]
\begin{center}
\begin{tabularx}{\linewidth}{|X|X|X|X|}
\hline
\rowcolor{popblueL} \textbf{Verbe} & \textbf{Nom (action/état)} & \textbf{Nom (agent)} & \textbf{Adjectif / Participe} \\ \hline
\eng{Subscribe} & \eng{Subscription} & \eng{Subscriber} & \eng{Subscribed} \\ \hline
\eng{Consume} & \eng{Consumption} & \eng{Consumer} & \eng{Consumable} \\ \hline
\eng{Connect} & \eng{Connection / Connectivity} & \eng{—} & \eng{Connected / Connecting} \\ \hline
\eng{Recharge / Top up} & \eng{Recharge / Top-up} & \eng{—} & \eng{Rechargeable} \\ \hline
\eng{Bundle} & \eng{Bundle / Bundling} & \eng{—} & \eng{Bundled} \\ \hline
\eng{Expire} & \eng{Expiry / Expiration} & \eng{—} & \eng{Expired} \\ \hline
\eng{Bill / Charge} & \eng{Bill / Billing / Charge} & \eng{—} & \eng{Billed / Charged} \\ \hline
\end{tabularx}
\end{center}
\emph{Attention :} \eng{Subscription} (abonnement) $\neq$ \eng{Subscribing} (action en cours). \eng{Expiry date} = date d'expiration. \eng{Billing cycle} = cycle de facturation.
\end{propriete}

\courssub{Entraînement guidé : de l'erreur à la norme}

\begin{exoresolu}[Correction d'erreurs types « Spécial Télécoms »]
\textbf{Consigne :} Corrigez les phrases suivantes. Chaque phrase contient au moins une erreur de faux ami, de calque ou de phrasal verb. Réécrivez la phrase correcte et justifiez la correction en français.

\begin{enumerate}
    \item \eng{*I have consumed my forfait, I need to buy a new credit.}
    \item \eng{*The cost of this bundle is expensive, but the value for money is good.}
    \item \eng{*I cannot join my friend, his phone is cut.}
    \item \eng{*Subscribe this package if you want more data.}
    \item \eng{*Eventually, I will consume all my data if I watch videos.}
\end{enumerate}

\tcblower
\textbf{Solutions détaillées :}

\begin{enumerate}
    \item \textbf{Correct :} \eng{I have \textbf{used up} my \textbf{bundle} / \textbf{plan}, I need to \textbf{top up} / \textbf{buy airtime}.} \\
    \textit{Justification :} \eng{Forfait} $\rightarrow$ \eng{bundle/plan}. \eng{Consumed} $\rightarrow$ \eng{used up} (épuisé un stock). \eng{Buy a credit} $\rightarrow$ \eng{top up / buy airtime} (recharger).
    \item \textbf{Correct :} \eng{The \textbf{price} of this bundle \textbf{is high} / This bundle \textbf{is expensive}, but it offers good \textbf{value for money}.} \\
    \textit{Justification :} Calque \eng{*cost is expensive}. \eng{Price/Cost} = \eng{high/low}. \eng{Bundle/Plan} = \eng{expensive/cheap}. \eng{Value for money} est correct.
    \item \textbf{Correct :} \eng{I can't \textbf{get through to} my friend, his phone \textbf{is cut off} / \textbf{got cut off}.} \\
    \textit{Justification :} \eng{Join} = rejoindre (physiquement) / adhérer. \eng{Get through to} = joindre au téléphone. \eng{His phone is cut} n'existe pas. \eng{Be cut off} (passif) = être coupé (communication).
    \item \textbf{Correct :} \eng{\textbf{Subscribe to} / \textbf{Sign up for} this \textbf{bundle / plan} if you want more data.} \\
    \textit{Justification :} \eng{Subscribe} intransitif $\rightarrow$ \eng{Subscribe TO}. \eng{Package} possible mais \eng{bundle/plan} plus naturel pour télécoms. \eng{Forfait} $\neq$ \eng{forfeit}.
    \item \textbf{Correct :} \eng{\textbf{Eventually} / \textbf{Finally}, I will \textbf{use up} / \textbf{run out of} all my data if I watch videos.} \\
    \textit{Justification :} \eng{Eventually} = « finalement / à la fin » (OK ici si sens « au bout du compte »). Mais \eng{consume} $\rightarrow$ \eng{use up / run out of}. Attention : \eng{Eventually} $\neq$ \eng{éventuellement} (=\eng{possibly}).
\end{enumerate}
\end{exoresolu}

\begin{experience}[Jeu de rôle : Au service client d'un opérateur (MTN / Orange / Nexttel)]
\textbf{Contexte :} Vous appelez le service client pour résoudre un problème de forfait. Travaillez par paires (A = client, B = agent).
\textbf{Rôles :}
\begin{itemize}
    \item \textbf{Client (A) :} Vous avez \eng{run out of data} avant la fin du mois. Votre \eng{current plan} ne vous suffit plus. Vous voulez \eng{upgrade} vers un \eng{bundle} avec plus de data, mais vous trouvez le \eng{price high}. Négociez une \eng{value for money} offre.
    \item \textbf{Agent (B) :} Proposez deux \eng{packages} : un \eng{monthly bundle} 20 GB à 5000 FCFA et un \eng{unlimited plan} à 10000 FCFA. Expliquez la \eng{billing cycle} et comment \eng{top up} via \eng{USSD code} ou \eng{Mobile Money}.
\end{itemize}
\textbf{Contraintes :} Utilisez au moins 5 phrasal verbs / expressions vus dans la leçon (\eng{top up, run out of, get through to, cut off, value for money, subscribe to, sign up for, current plan, bundle, airtime, balance}).
\textbf{Débriefing :} Échangez les rôles. Notez les erreurs de calques entendues.
\end{experience}

\begin{exercice}[À vous de jouer : Traduction et reformulation]
\textbf{Traduisez en anglais naturel (évitez les calques !) :}
\begin{enumerate}
    \item Je dois recharger mon téléphone, je n'ai plus de crédit.
    \item Ce forfait est cher, mais le rapport qualité-prix est excellent.
    \item Je n'arrive pas à joindre le service client, la ligne est coupée.
    \item Mon forfait actuel expire demain, je veux m'abonner à un nouveau pack.
    \item Si tu regardes trop de vidéos, tu vas épuiser tes données très vite.
\end{enumerate}
\end{exercice}

\begin{corrige}[Exercice : Faux amis et calques]
\begin{enumerate}
    \item \eng{I need to \textbf{top up} my phone, I've \textbf{run out of credit / airtime}.} (\textit{Pas : *buy credit, *my phone has no money})
    \item \eng{This \textbf{bundle / plan is expensive}, but it offers excellent \textbf{value for money}.} (\textit{Pas : *The cost is expensive})
    \item \eng{I can't \textbf{get through to} customer service, the line \textbf{keeps getting cut off} / \textbf{got cut off}.} (\textit{Pas : *join, *the line is cut})
    \item \eng{My \textbf{current plan} expires tomorrow, I want to \textbf{subscribe to} / \textbf{sign up for} a new \textbf{bundle}.} (\textit{Pas : *actual plan, *subscribe a package, *forfait})
    \item \eng{If you watch too many videos, you'll \textbf{use up} / \textbf{run out of} your data very quickly.} (\textit{Pas : *consume your data})
\end{enumerate}
\end{corrige}

\coursec{Mots de la même famille et formation des mots}

Après avoir identifié les \cle{faux amis} et les \cle{calques} qui piègent l'apprenant francophone (section précédente), nous allons maintenant voir comment \textbf{multiplier votre vocabulaire} de manière systématique. En anglais, un mot racine génère souvent une famille entière : verbe, nom, adjectif, parfois adverbe. Maîtriser ces mécanismes de \cle{dérivation} et de \cle{composition} vous permet de deviner le sens de mots inconnus et d'exprimer des nuances précises sans apprendre chaque mot isolément. C'est la clé pour passer d'un vocabulaire « de survie » à un vocabulaire « expert » sur les services télécoms.

\courssub{Observation dans un texte authentique : la plainte d'un abonné}

Lisez ce message posté sur le forum d'assistance de \eng{MTN Cameroon} par un utilisateur frustré. Repérez les mots en \textbf{gras} et essayez de deviner leur nature grammaticale (verbe, nom, adjectif) et leur lien de parenté.

\begin{textebox}[Forum MTN Cameroon — Customer Complaint]
\textbf{Subject:} Unfair deduction on my \textbf{subscription} — \textbf{validity} issue

Hi \textbf{Moderators},

I \textbf{subscribed} to the \textbf{Weekly Data Bundle} (5 GB) three days ago. I received the confirmation \textbf{SMS} stating the \textbf{validity} was 7 days. However, this morning I got a notification that my bundle has \textbf{expired}. I checked my \textbf{account balance} and the data is gone.

I tried to \textbf{recharge} my line with a scratch card, but the USSD code failed due to poor \textbf{connectivity}. I am a loyal \textbf{subscriber} since 2015 and I expect better \textbf{communication} from your \textbf{operators}. 

Please \textbf{validate} my complaint and \textbf{renew} my bundle immediately. I am on a \textbf{prepaid} plan, not \textbf{postpaid}, so I manage my \textbf{payments} carefully. This \textbf{disconnection} is unacceptable.

\textbf{Regards,}
\textbf{Franck M.} (Buea)
\end{textebox}

\begin{definition}[Glossaire du texte]
\begin{itemize}
    \item \eng{Moderator} (n) : modérateur (forum).
    \item \eng{Deduction} (n) : déduction, prélèvement.
    \item \eng{Bundle} (n) : forfait, pack (data/voix).
    \item \eng{Scratch card} (n) : carte à gratter (recharge) ; synonyme : \eng{recharge card}, \eng{voucher}.
    \item \eng{USSD code} (n) : code USSD (ex. *123\#).
    \item \eng{Loyal} (adj) : fidèle.
    \item \eng{Regards} (n pl) : salutations (formule de politesse).
    \item \eng{Expired} (v, past part.) : expiré, arrivé à échéance.
\end{itemize}
\end{definition}

\courssub{Analyse des familles lexicales : règles de formation}

Observez le tableau ci-dessous. Chaque ligne part d'un \cle{radical} (souvent le verbe) et ajoute des \cle{suffixes} pour changer la catégorie grammaticale. C'est la \cle{dérivation suffixale}.

\begin{propriete}[Règle générale : les suffixes nominaux et adjectivaux fréquents]
\begin{itemize}
    \item \textbf{-tion / -sion} : transforme un verbe en nom d'action ou de résultat (\eng{communicate} $\to$ \eng{communication}, \eng{subscribe} $\to$ \eng{subscription}). \textbf{Attention :} l'orthographe change souvent (disparition du \textit{e} final, \textit{ss} $\to$ \textit{ssion}, \textit{de} $\to$ \textit{ssion}).
    \item \textbf{-ity} : transforme un adjectif en nom abstrait (\eng{valid} $\to$ \eng{validity}, \eng{connective} $\to$ \eng{connectivity}).
    \item \textbf{-ive} : transforme un verbe/nom en adjectif (\eng{communicate} $\to$ \eng{communicative}).
    \item \textbf{-er / -or} : transforme un verbe en nom d'agent (celui qui fait l'action) : \eng{subscribe} $\to$ \eng{subscriber}, \eng{operate} $\to$ \eng{operator}, \eng{provide} $\to$ \eng{provider}.
    \item \textbf{-ment} : transforme un verbe en nom d'action/résultat (\eng{pay} $\to$ \eng{payment}, \eng{register} $\to$ \eng{registration}).
\end{itemize}
\textbf{Comparaison FR/EN :} Le français utilise souvent \textit{-tion}, \textit{-ité}, \textit{-if}, \textit{-eur}, \textit{-ment} de façon très parallèle. C'est un point d'appui majeur pour les francophones.
\end{propriete}

\begin{center}
\begin{tabularx}{\linewidth}{|X|X|X|X|}
\hline
\rowcolor{popblueL}
\textbf{Radical / Verbe} & \textbf{Nom (action/état)} & \textbf{Nom (agent)} & \textbf{Adjectif(s)} \\
\hline
\eng{to subscribe} (s'abonner) & \eng{subscription} (abonnement) & \eng{subscriber} (abonné) & \eng{subscribed} (abonné, part. passé) \\
\hline
\eng{to connect} (connecter) & \eng{connection} (connexion) & — & \eng{connected} / \eng{disconnected} (connecté / déconnecté) \\
 & \eng{connectivity} (connectivité) & & \\
\hline
\eng{to communicate} (communiquer) & \eng{communication} (communication) & \eng{communicator} (communicant) & \eng{communicative} (communicatif) \\
 & \eng{telecommunications} (télécoms) & & \\
\hline
\eng{to pay} (payer) & \eng{payment} (paiement) & \eng{payer} (payeur) & \eng{prepaid} (prépayé) / \eng{postpaid} (postpayé) \\
 & \eng{billing} (facturation) & & \eng{pay-as-you-go} (paiement à l'usage) \\
\hline
\eng{to validate} (valider) & \eng{validation} (validation) & \eng{validator} (validateur) & \eng{valid} (valide) / \eng{invalid} (invalide) \\
\eng{—} (base adj. \eng{valid}) & \eng{validity} (validité) & — & \\
\hline
\end{tabularx}
\end{center}

\courssub{Zoom sur la famille de CONNECT : arbre de dérivation}

La famille de \eng{connect} est particulièrement riche pour le domaine télécoms. Voici sa structure visuelle.

\begin{popfigure}
\begin{center}\tcbox[colback=popblue,colframe=popblue,coltext=white,arc=9pt,boxsep=4pt,left=6pt,right=6pt]{\bfseries\small connect (verb)}\end{center}
\vspace{-2pt}
\begin{tcbraster}[raster columns=2,raster equal height=rows,raster column skip=6pt,raster row skip=6pt,enhanced,blankest]
\begin{tcolorbox}[enhanced,colback=white,colframe=popblue,boxrule=0.9pt,arc=3pt,title={connection (noun)},fonttitle=\bfseries\scriptsize,coltitle=white,colbacktitle=popblue,left=4pt,right=4pt,top=2pt,bottom=2pt,boxsep=1pt,toptitle=1.5pt,bottomtitle=1.5pt,halign=left]
\begin{itemize}[nosep,leftmargin=*,itemsep=1pt,font=\scriptsize]
\item reconnection
\item disconnection
\end{itemize}
\end{tcolorbox}
\begin{tcolorbox}[enhanced,colback=white,colframe=poporange,boxrule=0.9pt,arc=3pt,title={connectivity (noun)},fonttitle=\bfseries\scriptsize,coltitle=white,colbacktitle=poporange,left=4pt,right=4pt,top=2pt,bottom=2pt,boxsep=1pt,toptitle=1.5pt,bottomtitle=1.5pt,halign=left]
\begin{itemize}[nosep,leftmargin=*,itemsep=1pt,font=\scriptsize]
\item high-speed
\item mobile
\end{itemize}
\end{tcolorbox}
\begin{tcolorbox}[enhanced,colback=white,colframe=popgreen,boxrule=0.9pt,arc=3pt,title={connected (adj)},fonttitle=\bfseries\scriptsize,coltitle=white,colbacktitle=popgreen,left=4pt,right=4pt,top=2pt,bottom=2pt,boxsep=1pt,toptitle=1.5pt,bottomtitle=1.5pt,halign=left]
\begin{itemize}[nosep,leftmargin=*,itemsep=1pt,font=\scriptsize]
\item well-connected
\item interconnected
\end{itemize}
\end{tcolorbox}
\begin{tcolorbox}[enhanced,colback=white,colframe=poppurple,boxrule=0.9pt,arc=3pt,title={disconnected (adj)},fonttitle=\bfseries\scriptsize,coltitle=white,colbacktitle=poppurple,left=4pt,right=4pt,top=2pt,bottom=2pt,boxsep=1pt,toptitle=1.5pt,bottomtitle=1.5pt,halign=left]
\begin{itemize}[nosep,leftmargin=*,itemsep=1pt,font=\scriptsize]
\item offline
\end{itemize}
\end{tcolorbox}
\end{tcbraster}

\legende{Arbre de dérivation de \eng{connect} : le radical central génère noms, adjectifs et composés par préfixation (re-, dis-, inter-) et suffixation (-ity, -ed).}
\end{popfigure}

\courssub{Les préfixes productifs : \texorpdfstring{\eng{re-}}{re-}, \texorpdfstring{\eng{un-}}{un-}, \texorpdfstring{\eng{dis-}}{dis-}, \texorpdfstring{\eng{pre-}}{pre-}, \texorpdfstring{\eng{post-}}{post-}}

En téléphonie, plusieurs préfixes reviennent sans cesse pour exprimer la gestion du forfait : la répétition, la négation, l'antériorité et la postériorité.

\begin{propriete}[Règles des préfixes courants dans le domaine télécoms]
\textbf{1. \eng{re-} (répétition, retour en arrière)} \\
S'ajoute aux verbes et noms. Sens : « refaire », « faire à nouveau », « restaurer ».
\begin{itemize}
    \item \eng{recharge} (v/n) : recharger / recharge (crédit). \eng{« I need to \textbf{recharge} my line. »} (Synonyme UK : \eng{top up}).
    \item \eng{renew} (v) : renouveler. \eng{« The bundle \textbf{renews} automatically every Sunday. »}
    \item \eng{reconnect} (v) : se reconnecter. \eng{« The app tries to \textbf{reconnect} after a drop. »}
    \item \eng{re-subscribe} / \eng{resubscribe} (v) : se réabonner.
\end{itemize}

\textbf{2. \eng{un-} (négation, privation, retour d'état)} \\
S'ajoute aux adjectifs, participes passés (adjectivaux) et certains verbes.
\begin{itemize}
    \item \eng{unlimited} (adj) : illimité. \eng{« \textbf{Unlimited} calls to MTN numbers. »}
    \item \eng{unsubscribe} (v) : se désabonner. \eng{« Dial *123\# to \textbf{unsubscribe} from the service. »}
    \item \eng{unlimited} vs \eng{limited} (limité).
    \item \eng{unpaid} (adj) : impayé. \eng{« \textbf{Unpaid} bills lead to disconnection. »}
    \item \eng{unregistered} (adj) : non enregistré.
\end{itemize}

\textbf{3. \eng{dis-} (négation d'action, séparation, inversion)} \\
Principalement sur verbes et noms d'action.
\begin{itemize}
    \item \eng{disconnect} (v) : déconnecter, couper. \eng{« The call was \textbf{disconnected}. »} (Pas \eng{*unconnect}).
    \item \eng{disable} (v) : désactiver. \eng{« \textbf{Disable} data roaming to save money. »}
\end{itemize}

\textbf{4. \eng{pre-} / \eng{post-} (antériorité / postériorité temporelle ou logique)} \\
\begin{itemize}
    \item \eng{prepaid} (adj) : prépayé (on paie avant de consommer).
    \item \eng{postpaid} (adj) : postpayé (on consomme, puis on reçoit une facture).
    \item \eng{pre-register} (v) : pré-enregistrer.
\end{itemize}
\textbf{Note :} \eng{un-} ne s'utilise pas devant tous les verbes (on dit \eng{disconnect}, pas \eng{unconnect} ; \eng{disable}, pas \eng{unable} comme verbe).
\end{propriete}

\begin{exemplebox}[Exemples traduits — Préfixes en contexte]
\begin{itemize}
    \item \eng{Subscribers can \textbf{renew} their plans automatically.} \\
    « Les abonnés peuvent \textbf{renouveler} leurs forfaits automatiquement. »
    \item \eng{The \textbf{validity} of this bundle is seven days.} \\
    « La \textbf{validité} de ce forfait est de sept jours. »
    \item \eng{He chose a \textbf{prepaid} plan because he wants to control his spending.} \\
    « Il a choisi un forfait \textbf{prépayé} car il veut contrôler ses dépenses. »
    \item \eng{To \textbf{unsubscribe}, send STOP to 8080.} \\
    « Pour \textbf{se désabonner}, envoyez STOP au 8080. »
    \item \eng{The network offers \textbf{unlimited} social media access on this package.} \\
    « Le réseau offre un accès \textbf{illimité} aux réseaux sociaux sur ce pack. »
    \item \eng{Please \textbf{disable} mobile data when abroad to avoid roaming charges.} \\
    « Veuillez \textbf{désactiver} les données mobiles à l'étranger pour éviter les frais d'itinérance. »
\end{itemize}
\end{exemplebox}

\courssub{Prononciation et accent tonique : le piège du déplacement}

C'est l'erreur n°1 des francophones : lire l'anglais comme du français. En anglais, l'accent tonique (\cle{word stress}) détermine le rythme de la phrase. Quand on ajoute un suffixe comme \eng{-tion} ou \eng{-ity}, l'accent \textbf{se déplace} souvent vers la syllabe précédant le suffixe.

\begin{attention}[Pièges de prononciation — Déplacement de l'accent tonique]
\begin{center}
\begin{tabular}{|l|l|l|l|}
\hline
\rowcolor{poporangeL}
\textbf{Mot de base} & \textbf{Prononciation approx. (FR)} & \textbf{Dérivé} & \textbf{Prononciation approx. (FR)} \\
\hline
\eng{subSCRIBE} & « sèb-\textbf{SKRAÏB} » & \eng{subSCRIPtion} & « sèb-\textbf{SKRIP}-cheune » \\
\hline
\eng{commuNIcate} & « ke-miou-\textbf{NI}-kéït » & \eng{communiCAtion} & « ke-miou-ni-\textbf{KÉI}-cheune » \\
\hline
\eng{VA-lid} & « \textbf{VÉ}-lide » & \eng{vaLI-dity} & « va-\textbf{LI}-di-ti » \\
\hline
\eng{PAY} & « \textbf{PÉÏ} » & \eng{PAY-ment} & « \textbf{PÉÏ}-meunte » (accent stable) \\
\hline
\eng{conNECT} & « ke-\textbf{NÈKTE} » & \eng{connecTIVity} & « ke-neck-\textbf{TI}-vi-ti » \\
\hline
\eng{reCHARGE} & « ri-\textbf{TCHARDJ} » & \eng{REcharge} (n) & « \textbf{RI}-tchardj » (accent sur préfixe nom) \\
\hline
\end{tabular}
\end{center}
\textbf{Règle mnémotechnique :} \eng{-tion} attire l'accent sur la syllabe d'avant (\textit{-ta-tion} $\to$ accent sur \textit{ta}). \eng{-ity} attire l'accent sur la syllabe d'avant (\textit{-vi-ty} $\to$ accent sur \textit{vi}). \textbf{Entraînez-vous à taper du pied sur la syllabe forte !}
\end{attention}

\courssub{Méthode : enrichir son lexique par les familles de mots}

\begin{methode}[Comment tripler son vocabulaire sans effort]
\textbf{Objectif :} À partir d'un mot nouveau, reconstruire sa famille complète pour le fixer durablement.
\begin{enumerate}
    \item \textbf{Identifiez le radical} (souvent le verbe à l'infinitif sans \textit{to}).
    \item \textbf{Appliquez les suffixes standards} : \eng{-tion/-sion} (nom action), \eng{-ment} (nom action), \eng{-er/-or} (nom agent), \eng{-ing} (nom action/gérondif), \eng{-ed/-ing} (adjectifs participes), \eng{-ive} (adjectif), \eng{-ity} (nom qualité), \eng{-al} (adjectif \eng{operational}).
    \item \textbf{Testez les préfixes courants} : \eng{re-}, \eng{un-}, \eng{dis-}, \eng{pre-}, \eng{post-}, \eng{inter-}, \eng{multi-}, \eng{non-}.
    \item \textbf{Vérifiez l'orthographe} (changements : \textit{y $\to$ i} \eng{apply $\to$ application}, \textit{e} muet tombe \eng{subscribe $\to$ subscription}, doublement consonne \eng{run $\to$ running}).
    \item \textbf{Notez la prononciation} de chaque forme (accent tonique !).
    \item \textbf{Écrivez une phrase personnelle} pour chaque forme (contexte camerounais : MTN, Orange, Camtel, Nexttel, forfait, data, crédit, MoMo).
\end{enumerate}
\textbf{Exemple appliqué à \eng{operate} :} \eng{operator, operation, operational, operationalize, pre-operational, co-operate/cooperate, operating system}...
\end{methode}

\courssub{Le lexique en contexte : dialogue authentique (Boutique Orange Douala)}

Ce dialogue met en scène un client (\eng{Customer}) et un agent (\eng{Agent}) dans une boutique Orange à Douala. Il recycle \textbf{toutes} les familles étudiées.

\begin{textebox}[Dialogue : At the Orange Shop — Douala Akwa]
\textbf{Agent:} Good morning, sir. Welcome to Orange. How can I help you?

\textbf{Customer:} Good morning. I want to \textbf{subscribe} to a new data \textbf{bundle}. My current \textbf{subscription} expires today.

\textbf{Agent:} Alright. Are you on a \textbf{prepaid} or \textbf{postpaid} plan?

\textbf{Customer:} \textbf{Prepaid}. I prefer to \textbf{recharge} (ou \textbf{top up}) when I need to. I don't want a monthly \textbf{commitment}.

\textbf{Agent:} Understood. We have the \textbf{« Mega Data »} package. 20 GB for 30 days. The \textbf{validity} is one month.

\textbf{Customer:} Is the \textbf{connection} speed good in my area? I live near \textbf{Bonapriso}.

\textbf{Agent:} Yes, sir. 4G \textbf{connectivity} is excellent there. You will get high-speed \textbf{access}.

\textbf{Customer:} Okay. And if I finish the data before the \textbf{validity} period ends?

\textbf{Agent:} You can \textbf{renew} the bundle early via USSD *141\# or the \textbf{MyOrange} app. The \textbf{payment} is deducted from your main \textbf{account balance}.

\textbf{Customer:} Perfect. What if I want to \textbf{unsubscribe} later?

\textbf{Agent:} Just dial *141*0\#. It's instant. No \textbf{penalty} fee.

\textbf{Customer:} Great. Let's do it. Here is my ID for \textbf{registration}.

\textbf{Agent:} Thank you. Your \textbf{subscription} is now \textbf{active}. You will receive a confirmation \textbf{SMS}. Enjoy your \textbf{communication}!

\textbf{Customer:} Thank you. Have a good day.
\end{textebox}

\begin{definition}[Glossaire clé du dialogue]
\begin{itemize}
    \item \eng{Bundle} (n) : forfait, pack (synonyme de \eng{package}, \eng{plan}, \eng{data plan}).
    \item \eng{Commitment} (n) : engagement (souvent 12/24 mois pour postpaid).
    \item \eng{Area} (n) : zone, quartier.
    \item \eng{High-speed access} (n) : accès haut débit.
    \item \eng{Account balance} (n) : solde du compte (crédit principal).
    \item \eng{Deduct} (v) : déduire, prélever.
    \item \eng{Penalty} (n) : pénalité, frais de résiliation anticipée.
    \item \eng{Registration} (n) : enregistrement (obligatoire au Cameroun : \eng{SIM registration}).
    \item \eng{Active} (adj) : actif, activé.
    \item \eng{Top up} (v/n) : recharger / recharge (synonyme de \eng{recharge}, très courant au UK et en Afrique).
\end{itemize}
\end{definition}

\courssub{Exercice résolu : transformation de phrases}

\begin{exoresolu}[Transformation : forme nominale $\leftrightarrow$ verbale $\leftrightarrow$ adjectivale]
\textbf{Consigne :} Réécrivez chaque phrase en utilisant le mot entre parenthèses (changez la forme si nécessaire) sans changer le sens.

\begin{enumerate}
    \item \eng{He \textbf{pays} his bill every month.} (\eng{payment}) \\
    \textbf{Solution :} \eng{He makes his monthly \textbf{payment} regularly.} \\
    \textit{Analyse :} Verbe \eng{pay} $\to$ Nom \eng{payment}. Ajout de \eng{make a} (collocation). \eng{Monthly} (adj) dérivé de \eng{month}.
    
    \item \eng{The \textbf{validity} of the voucher is 30 days.} (\eng{valid}) \\
    \textbf{Solution :} \eng{The voucher \textbf{is valid} for 30 days.} \\
    \textit{Analyse :} Nom \eng{validity} $\to$ Adjectif \eng{valid} + verbe \eng{be}. Structure \eng{be valid for + durée}.
    
    \item \eng{Please \textbf{connect} to the Wi-Fi network.} (\eng{connection}) \\
    \textbf{Solution :} \eng{Please establish a \textbf{connection} to the Wi-Fi network.} \\
    \textit{Analyse :} Verbe \eng{connect} $\to$ Nom \eng{connection}. Verbe support \eng{establish} (établir) ou \eng{make}.
    
    \item \eng{She is a loyal \textbf{subscriber} of MTN.} (\eng{subscribe}) \\
    \textbf{Solution :} \eng{She \textbf{subscribed} to MTN years ago and stayed loyal.} \\
    \textit{Analyse :} Nom agent \eng{subscriber} $\to$ Verbe \eng{subscribe} (passé \eng{subscribed}). Attention : \eng{subscribe TO} (pas \eng{at}).
    
    \item \eng{The \textbf{communication} was cut off.} (\eng{disconnect}) \\
    \textbf{Solution :} \eng{They \textbf{got disconnected} / The call \textbf{was disconnected}.} \\
    \textit{Analyse :} Nom \eng{communication} $\to$ Verbe \eng{disconnect} (passif \eng{be disconnected}). \eng{Cut off} = \eng{disconnect} (couper la ligne).
\end{enumerate}
\end{exoresolu}

\courssub{Le saviez-vous ? : \texorpdfstring{\eng{Telecoms}}{Telecoms}, le mot passe-partout du Bac}

\begin{savaistu}[Telecoms : l'abréviation star des sujets d'examen]
Dans la presse anglophone (\eng{The Guardian}, \eng{BBC}, \eng{Financial Times}, \eng{The Post} au Cameroun) et dans les sujets de \textbf{compréhension écrite du Baccalauréat}, le mot \eng{telecommunications} est presque toujours abrégé en \textbf{\eng{telecoms}} (toujours avec un \textbf{s} final, même au singulier : \eng{the telecoms sector}).

\begin{itemize}
    \item \eng{The \textbf{telecoms} regulator in Cameroon is the ART (Agence de Régulation des Télécommunications).}
    \item \eng{\textbf{Telecoms} giants like MTN and Orange dominate the African market.}
    \item \eng{A career in \textbf{telecoms} engineering is promising.}
\end{itemize}
\textbf{Attention :} On ne dit pas \eng{*a telecom} (sauf adjectif : \eng{telecom engineer}). Le nom du secteur est \eng{telecoms} (pluriel en forme, singulier en sens souvent). C'est un \cle{mot-outil} pour comprendre les titres d'articles : \eng{« Telecoms Shake-up in Nigeria »}, \eng{« Cameroon Telecoms Revenue Hits Record High »}.
\end{savaistu}

\courssub{Exercices d'entraînement}

\begin{exercice}[Entraînement : Familles de mots et préfixes]
\textbf{A. Complétez le tableau suivant. Utilisez un dictionnaire si besoin.}
\begin{center}
\begin{tabular}{|l|l|l|l|}
\hline
\rowcolor{popblueL}
\textbf{Verbe} & \textbf{Nom (action/état)} & \textbf{Nom (agent)} & \textbf{Adjectif(s)} \\
\hline
\eng{to provide} & & \eng{provider} & \\
\hline
& \eng{operation} & \eng{operator} & \\
\hline
\eng{to register} & & & \eng{registered} \\
\hline
& \eng{activation} & & \eng{active} / \eng{inactive} \\
\hline
\eng{to bill} & & & \eng{billed} / \eng{unbilled} \\
\hline
\end{tabular}
\end{center}

\textbf{B. Choisissez la forme correcte (mot entre parenthèses) pour compléter les phrases.}
\begin{enumerate}
    \item My data \eng{\_\_\_\_\_\_\_\_\_\_} (subscribe / subscription / subscriber) expires tomorrow. I must renew it.
    \item The network \eng{\_\_\_\_\_\_\_\_\_\_} (connect / connectivity / connection) in my village is very slow; only 2G works.
    \item You can \eng{\_\_\_\_\_\_\_\_\_\_} (pay / payment / prepaid) your bill via Mobile Money (MoMo).
    \item This SIM card is not \eng{\_\_\_\_\_\_\_\_\_\_} (valid / validity / validate) anymore; it has been \eng{\_\_\_\_\_\_\_\_\_\_} (disconnect / disconnected / disconnection) for three months.
    \item To stop the service, dial the \eng{\_\_\_\_\_\_\_\_\_\_} (unsubscribe / unsubscription / unsubscribed) code.
    \item The \eng{\_\_\_\_\_\_\_\_\_\_} (operate / operator / operation) launched a new 5G \eng{\_\_\_\_\_\_\_\_\_\_} (package / pack / packaging) yesterday.
\end{enumerate}

\textbf{C. Traduisez en anglais en utilisant le vocabulaire de la leçon.}
\begin{enumerate}
    \item Je dois recharger mon crédit pour renouveler mon forfait data.
    \item Les abonnés postpayés reçoivent une facture à la fin du mois.
    \item La validité de cette carte de recharge est de 90 jours.
    \item Notre connexion internet a été coupée (disconnected) hier soir.
    \item Ce forfait offre des appels illimités vers tous les réseaux.
\end{enumerate}
\end{exercice}

\begin{corrige}[Exercice : Familles de mots et préfixes]
\textbf{A. Tableau complété}
\begin{center}
\begin{tabular}{|l|l|l|l|}
\hline
\rowcolor{popgreenL}
\textbf{Verbe} & \textbf{Nom (action/état)} & \textbf{Nom (agent)} & \textbf{Adjectif(s)} \\
\hline
\eng{to provide} & \eng{provision} / \eng{provisioning} & \eng{provider} & \eng{provided} / \eng{provisional} \\
\hline
\eng{to operate} & \eng{operation} & \eng{operator} & \eng{operational} / \eng{operative} \\
\hline
\eng{to register} & \eng{registration} & \eng{registrar} & \eng{registered} / \eng{unregistered} \\
\hline
\eng{to activate} & \eng{activation} & \eng{activator} & \eng{active} / \eng{inactive} \\
\hline
\eng{to bill} & \eng{billing} & \eng{biller} & \eng{billed} / \eng{unbilled} \\
\hline
\end{tabular}
\end{center}

\textbf{B. Formes correctes}
\begin{enumerate}
    \item \eng{subscription} (nom, sujet de \eng{expires}).
    \item \eng{connectivity} (nom qualité du réseau, synonyme de \eng{connection quality} ; \eng{connection} possible mais \eng{connectivity} rend mieux la qualité du signal).
    \item \eng{pay} (verbe après modal \eng{can}). \textit{Note :} \eng{make a payment} possible mais \eng{pay} plus direct.
    \item \eng{valid} (adjectif après \eng{is not}) / \eng{disconnected} (participe passé adjectival après \eng{has been}).
    \item \eng{unsubscription} (nom attributif modifiant \eng{code} : « code de désabonnement ») ; \eng{unsubscribe code} est fréquent dans l'UI mais \eng{unsubscription code} est grammaticalement plus standard pour un nom composé.
    \item \eng{operator} (nom agent, sujet de \eng{launched}) / \eng{package} (nom, synonyme de \eng{bundle/plan} ; \eng{pack} possible mais moins formel, \eng{packaging} = emballage).
\end{enumerate}

\textbf{C. Traductions proposées}
\begin{enumerate}
    \item \eng{I need to \textbf{recharge} (ou \textbf{top up}) my credit to \textbf{renew} my data \textbf{bundle} (ou \textbf{subscription}).}
    \item \eng{\textbf{Postpaid subscribers} (ou \textbf{customers}) receive a \textbf{bill} at the end of the month.}
    \item \eng{The \textbf{validity} of this \textbf{recharge card} (ou \textbf{voucher} / \textbf{scratch card}) is 90 days.}
    \item \eng{Our internet \textbf{connection} \textbf{was disconnected} (ou \textbf{got disconnected}) last night.}
    \item \eng{This \textbf{bundle} (ou \textbf{plan} / \textbf{package}) offers \textbf{unlimited} calls to all \textbf{networks}.}
\end{enumerate}
\end{corrige}

\coursec{Le lexique en contexte : dialogue et texte authentiques}

Dans la section précédente, nous avons appris à construire des familles de mots (\emph{to subscribe, subscriber, subscription} ; \emph{to connect, connection} ; \emph{valid, validity}) et à les employer correctement. Mais un mot n'existe vraiment que lorsqu'il est utilisé dans une situation réelle. Cette dernière section te montre donc le lexique des forfaits téléphoniques et internet \textbf{en action} : d'abord dans un dialogue authentique à l'agence télécom, puis dans un article de presse sur l'internet mobile au Cameroun. Tu y apprendras aussi la méthode de lecture qui te fera réussir les questions de compréhension au Baccalauréat.

\courssub{Dialogue modèle : choisir un forfait à l'agence télécom}

\begin{textebox}[At the telecom agency — Choosing a bundle]
\textbf{Agent:} Good morning! How can I help you?\\
\textbf{Customer:} Good morning. I'd like to subscribe to an internet bundle. What plans do you have?\\
\textbf{Agent:} We have a 2 GB weekly plan for 1,000 FCFA and a 5 GB monthly plan for 3,000 FCFA, with free night calls.\\
\textbf{Customer:} How long is the weekly plan valid?\\
\textbf{Agent:} Seven days, from the day of activation.\\
\textbf{Customer:} And what does the monthly plan include exactly?\\
\textbf{Agent:} It includes 5 GB of data, unlimited SMS and free calls from midnight to 5 a.m.\\
\textbf{Customer:} I see. I'll take the monthly plan. Does unused data roll over?\\
\textbf{Agent:} Yes, it does. If you renew before the expiry date, the remaining megabytes are carried over to the next month.\\
\textbf{Customer:} Perfect. How do I activate it?\\
\textbf{Agent:} Please dial *150\# and follow the instructions. You will receive a confirmation SMS.\\
\textbf{Customer:} Thank you very much. Have a nice day!\\
\textbf{Agent:} You're welcome. Goodbye!
\end{textebox}

\textbf{Glossaire du dialogue}

\begin{center}
\begin{tabularx}{\linewidth}{|l|X|l|}\hline
\rowcolor{popblueL} \textbf{Mot anglais} & \textbf{Nature} & \textbf{Traduction} \\ \hline
bundle & nom & forfait (de données) \\ \hline
plan & nom & forfait, offre \\ \hline
weekly / monthly & adjectifs & hebdomadaire / mensuel \\ \hline
valid & adjectif & valable \\ \hline
to include & verbe & comprendre, inclure \\ \hline
unlimited & adjectif & illimité \\ \hline
to roll over & verbe & être reporté (données non utilisées) \\ \hline
to renew & verbe & renouveler \\ \hline
expiry date & nom composé & date d'expiration \\ \hline
to dial & verbe & composer (un numéro) \\ \hline
\end{tabularx}
\end{center}

\courssub{Formules fonctionnelles : les phrases clés du dialogue}

Observe les questions du client : chacune a une fonction précise. Voici le tableau des formules à connaître par cœur pour tout dialogue commercial sur les forfaits.

\begin{center}
\begin{tabularx}{\linewidth}{|X|X|X|}\hline
\rowcolor{popblueL} \textbf{English} & \textbf{Français} & \textbf{Fonction} \\ \hline
How can I help you? & Comment puis-je vous aider ? & accueillir le client \\ \hline
I'd like to subscribe to an internet bundle. & Je voudrais souscrire à un forfait internet. & exprimer un souhait poli \\ \hline
What plans do you have? & Quels forfaits avez-vous ? & demander l'offre \\ \hline
How much is the 2 GB bundle? & Combien coûte le forfait de 2 Go ? & demander le prix \\ \hline
What does this plan include? & Que comprend ce forfait ? & demander le contenu \\ \hline
How long is it valid? & Combien de temps est-il valable ? & demander la durée \\ \hline
Does unused data roll over? & Les données non utilisées sont-elles reportées ? & vérifier une condition \\ \hline
I'll take the monthly plan. & Je prends le forfait mensuel. & décider, choisir \\ \hline
Please dial *150\# to activate it. & Veuillez composer le *150\# pour l'activer. & donner une instruction \\ \hline
\end{tabularx}
\end{center}

\begin{exemplebox}[Exemples traduits]
\eng{“What does this plan include?”} « Que comprend ce forfait ? »\\
\eng{“It includes unlimited SMS and 5 GB of data.”} « Il comprend des SMS illimités et 5 Go de données. »\\
\eng{“How much is the weekly bundle?”} « Combien coûte le forfait hebdomadaire ? »\\
\eng{“It costs 1,000 FCFA and it is valid for seven days.”} « Il coûte 1 000 FCFA et il est valable sept jours. »
\end{exemplebox}

\begin{attention}[How much ou How long ?]
Dans les dialogues commerciaux, \eng{How much is...?} demande le \textbf{prix} (\emph{It costs 3,000 FCFA}) ; \eng{How long is it valid?} demande la \textbf{durée} (\emph{It is valid for one month}). Ne confonds pas les deux questions : répondre \eng{“It is 3,000 FCFA”} à \eng{“How long is it valid?”} est une erreur classique au Bac. Autre piège : on dit \eng{I'd like to subscribe to a plan} (avec \emph{to}), jamais \eng{“subscribe a plan”} sur le modèle français « souscrire un forfait ».
\end{attention}

\courssub{Texte de lecture : la croissance de l'internet mobile au Cameroun}

\begin{textebox}[Mobile internet is changing Cameroon]
Ten years ago, going online in Cameroon meant finding a cybercafé in Yaoundé or Douala and paying for a slow connection by the hour. Today, millions of Cameroonians carry the internet in their pockets. With a cheap smartphone and a small data bundle, a student in Bamenda can watch a lesson on video, a cocoa farmer near Buea can check market prices, and a shopkeeper in Garoua can receive mobile money payments instantly.

Telecom companies have understood this new demand. They now compete fiercely to attract customers with attractive service packages: daily, weekly and monthly bundles, free night calls, unlimited SMS, and bonus data for subscribers who renew their plans before the expiry date. Prices have fallen: a gigabyte of data that cost a small fortune a few years ago is now affordable for many families.

However, challenges remain. In rural areas, network coverage is still weak, and many users complain that their data disappears too quickly. Consumer associations advise customers to compare plans carefully before subscribing, to check how long each bundle is valid, and to ask whether unused data rolls over to the next period.

Despite these problems, the growth of mobile internet is transforming education, business and even politics. For young Cameroonians, mastering this connected world — and the English vocabulary that describes it — has become an essential skill.
\end{textebox}

\textbf{Glossaire du texte}

\begin{center}
\begin{tabularx}{\linewidth}{|l|X|l|}\hline
\rowcolor{popblueL} \textbf{Mot anglais} & \textbf{Nature} & \textbf{Traduction} \\ \hline
to go online & expression & se connecter à internet \\ \hline
data bundle & nom composé & forfait de données \\ \hline
shopkeeper & nom & commerçant \\ \hline
to compete & verbe & être en concurrence \\ \hline
affordable & adjectif & abordable \\ \hline
coverage & nom & couverture (réseau) \\ \hline
to complain & verbe & se plaindre \\ \hline
skill & nom & compétence \\ \hline
\end{tabularx}
\end{center}

\courssub{Méthode de lecture : attaquer un texte au Bac}

\begin{methode}[Lire un texte en cinq étapes]
\begin{enumerate}
\item \textbf{Lis le titre} et anticipe le sujet : quels mots du champ lexical vont apparaître ?
\item \textbf{Repère les mots du thème} dans le texte (\emph{bundle, plan, data, subscribe, coverage...}) : ils te donnent la structure du texte.
\item \textbf{Devine le sens des mots inconnus par le contexte} : par exemple, si \emph{“prices have fallen”} et qu'un gigaoctet est \emph{affordable}, tu devines que ce mot signifie « abordable ».
\item \textbf{Vérifie au glossaire} (ou au dictionnaire en classe) pour confirmer tes hypothèses.
\item \textbf{Réponds aux questions} en localisant d'abord la phrase du texte qui contient la réponse.
\end{enumerate}
\end{methode}

\begin{popfigure}
\begin{center}
\begin{tikzpicture}[node distance=0.55cm,
box/.style={draw=popblue, fill=popblueL, rounded corners, text width=6.5cm, align=center, font=\small, thick},
fleche/.style={-{Stealth[length=3mm]}, thick, poporange}]
\node[box] (e1) {\textbf{1.} Lire le titre et anticiper le sujet};
\node[box, below=of e1] (e2) {\textbf{2.} Repérer les mots du champ lexical};
\node[box, below=of e2] (e3) {\textbf{3.} Deviner les mots inconnus par le contexte};
\node[box, below=of e3] (e4) {\textbf{4.} Vérifier au glossaire};
\node[box, below=of e4] (e5) {\textbf{5.} Localiser la réponse et répondre aux questions};
\draw[fleche] (e1) -- (e2);
\draw[fleche] (e2) -- (e3);
\draw[fleche] (e3) -- (e4);
\draw[fleche] (e4) -- (e5);
\end{tikzpicture}
\end{center}
\legende{Organigramme de la méthode de lecture en cinq étapes}
\end{popfigure}

\begin{methode}[Réussir les questions de compréhension au Bac]
\begin{enumerate}
\item Lis la question \textbf{avant} de relire le texte : tu sais ce que tu cherches.
\item \textbf{Souligne dans le texte} la phrase qui contient la réponse avant d'écrire quoi que ce soit.
\item \textbf{Reformule} avec tes propres mots : ne recopie jamais plus de quelques mots du texte (le correcteur pénalise la copie intégrale).
\item Réponds par une \textbf{phrase complète}, avec sujet et verbe, pas par un mot isolé.
\item Vérifie que ta réponse répond bien à la question posée (\emph{How much} = prix, \emph{How long} = durée, \emph{Why} = raison).
\end{enumerate}
\end{methode}

\begin{exoresolu}[Question de compréhension résolue pas à pas]
\textbf{Question:} \eng{“What two pieces of advice do consumer associations give to customers before subscribing?”} (Answer in your own words.)
\tcblower
\textbf{Étape 1 — Je localise.} La réponse se trouve dans le troisième paragraphe : \emph{“Consumer associations advise customers to compare plans carefully before subscribing, to check how long each bundle is valid, and to ask whether unused data rolls over.”}\\
\textbf{Étape 2 — Je souligne} les deux conseils : \emph{compare plans} et \emph{check validity / ask about rollover}.\\
\textbf{Étape 3 — Je reformule} sans copier la phrase entière :\\
\eng{Consumer associations recommend that customers should compare the different plans and check the validity period of each bundle before they subscribe.}\\
« Les associations de consommateurs recommandent aux clients de comparer les différents forfaits et de vérifier la durée de validité de chacun avant de souscrire. »\\
\textbf{Remarque :} j'ai remplacé \emph{advise} par \emph{recommend} et \emph{how long it is valid} par \emph{the validity period} : c'est cela, reformuler !
\end{exoresolu}

\begin{savaistu}[Un argument choc pour tes compositions]
Le Cameroun compte plus d'abonnés au téléphone mobile que d'habitants équipés en électricité : beaucoup de familles rechargent leur téléphone avant même d'avoir l'éclairage à la maison ! C'est un argument percutant pour tes compositions sur les médias et la communication : \eng{“In Cameroon, there are more mobile phone subscribers than households with electricity.”} « Au Cameroun, il y a plus d'abonnés au téléphone mobile que de foyers disposant d'électricité. »
\end{savaistu}

\courssub{Production guidée : à toi d'écrire}

\begin{experience}[Rédige ton propre dialogue]
Par groupes de deux, rédigez un court dialogue (10 à 12 répliques) chez un vendeur de forfaits à Douala. Vous devez \textbf{obligatoirement} employer ces 8 mots du lexique : \emph{bundle, to subscribe, valid, to include, unlimited, to roll over, to dial, expiry date}.\\
\textbf{Plan suggéré :} salutation → demande des offres disponibles → comparaison de deux prix (\emph{How much is...?}) → question sur la durée (\emph{How long is it valid?}) → question sur le contenu (\emph{What does it include?}) → décision (\emph{I'll take...}) → instruction d'activation (\emph{Please dial...}) → au revoir.\\
Puis jouez le dialogue devant la classe : le public note les 8 mots imposés à mesure qu'il les entend.
\end{experience}

\begin{exercice}[Paragraphe de conseil]
Rédige un paragraphe de 6 à 8 lignes intitulé \eng{Advice for choosing a good internet bundle}. Utilise au moins 6 mots du lexique de la leçon (\emph{to compare, plan, data, valid, affordable, to renew, coverage, to activate}) et deux formules interrogatives du dialogue.
\end{exercice}

\begin{corrige}[Exercice — modèle de paragraphe]
\eng{Before you subscribe to an internet bundle, you should compare the different plans carefully. First, ask: “How much is the bundle?” and “How long is it valid?” Then check what the plan includes: some bundles offer unlimited SMS, others give more data. Choose an affordable plan that matches your needs, and make sure the network coverage is good in your area. Finally, renew your bundle before the expiry date so that your unused data can roll over.}\\
« Avant de souscrire à un forfait internet, compare soigneusement les différentes offres. D'abord, demande : “Combien coûte le forfait ?” et “Combien de temps est-il valable ?”. Ensuite, vérifie ce que le forfait comprend : certains offrent des SMS illimités, d'autres plus de données. Choisis un forfait abordable adapté à tes besoins et assure-toi que la couverture réseau est bonne dans ta région. Enfin, renouvelle ton forfait avant la date d'expiration pour que tes données non utilisées soient reportées. »
\end{corrige}

\begin{aretenir}[L'essentiel de la section]
\begin{itemize}
\item Un dialogue commercial suit un schéma fixe : salutation → demande → comparaison → décision → activation.
\item \eng{How much is...?} = prix ; \eng{How long is it valid?} = durée ; \eng{What does it include?} = contenu.
\item Méthode de lecture : titre → mots du thème → devinette par le contexte → glossaire → réponse localisée et reformulée.
\item Au Bac, souligne la phrase-réponse dans le texte et reformule : ne recopie jamais le texte mot à mot.
\end{itemize}
\end{aretenir}

\newpage

\coursec{Activité d'intégration}

\coursec{Lesson 51 — Vocabulary: Exploiting Service Packages (Telephone/Internet Services)}

\courssub{Objectifs de l'activité}
\begin{itemize}
    \item Réinvestir le vocabulaire des forfaits téléphoniques et internet (\cle{data plan}, \cle{subscription}, \cle{rollover}, \cle{validity period}, \cle{bonus}, \cle{top-up}, \cle{unlimited calls}, \cle{fair usage policy}, \cle{roaming}, \cle{bundle}) dans une interaction orale authentique.
    \item Maîtriser les collocations clés (\cle{subscribe to a plan}, \cle{exceed the data limit}, \cle{renew a subscription}, \cle{check the balance}, \cle{activate a bundle}) et les formules fonctionnelles de conseil, de comparaison et de négociation.
    \item Éviter les faux amis et calques du français (\eng{subscription} ≠ \eng{subcription}, \eng{recharge} vs \cle{top-up}, \eng{forfait} → \cle{plan/package/bundle}, \eng{valider} → \cle{confirm/activate}, pas \eng{validate}).
    \item Produire un support écrit persuasif (affiche publicitaire) respectant les conventions du genre : slogan accrocheur, informations clés, mentions légales.
    \item Développer l'esprit critique en évaluant les offres (rapport qualité/prix, conditions cachées) et en votant pour la plus convaincante.
\end{itemize}

\courssub{Situation}
\begin{textebox}[Scenario: At the MTN / Orange / Camtel Service Centre — Yaoundé]
You are at a busy telecom service centre in Yaoundé (Carrefour Warda, Bastos or Mvog-Ada). Three roles are assigned per group:

\textbf{Role A — Sales Agent:} You work for \textit{ConnectPlus}, a new mobile virtual network operator (MVNO) targeting students and young professionals. You have two flagship offers: \textbf{“Student Boost”} and \textbf{“Pro Connect”}. You must present them clearly, using the functional language for comparing plans, explaining terms and closing the sale.

\textbf{Role B — Customer:} You are a Terminale A student at Lycée Leclerc or a young intern at a startup in Akwa. You have a budget of 5,000–10,000 FCFA/month. You need reliable internet for research, WhatsApp/Zoom classes, and social media, plus enough call minutes for family. You are hesitant: you have been disappointed by “unlimited” offers that throttle speed after 2 GB. You will ask at least four precise questions.

\textbf{Role C — Observer:} You hold the \textit{Lexical Checklist}. You tick every target word/collocation used correctly, note any French calque or false friend, and give feedback at the end.

\textbf{Offer Sheets (distributed to Sales Agents only):}

\textbf{Plan 1 — Student Boost} \\
Price: 5,500 FCFA / 30 days \\
Data: 15 GB (4G/4G+) — \textit{Fair Usage Policy: speed reduced to 256 kbps after 15 GB} \\
Calls: Unlimited on-net (ConnectPlus), 200 min off-net \\
SMS: Unlimited national \\
Bonus: Free access to educational platforms (Khan Academy, Coursera, Wikipedia Zero) — no data deducted \\
Rollover: Unused data rolls over once if renewed before expiry \\
Activation: Dial *123*1\# or \textit{MyConnect} app

\textbf{Plan 2 — Pro Connect} \\
Price: 9,500 FCFA / 30 days \\
Data: 40 GB (4G/4G+), then unlimited at 1 Mbps \\
Calls: Unlimited national (all networks) \\
SMS: Unlimited national \\
Bonus: 5 GB \textit{roaming} in CEMAC zone (Cameroon, Chad, CAR, Gabon, Equatorial Guinea, Congo) \\
Rollover: Full data rollover (no expiry if renewed monthly) \\
Priority customer care via WhatsApp \\
Activation: Dial *123*2\# or \textit{MyConnect} app

\textbf{General Terms (apply to both):} \\
\textit{Terms and conditions apply.} Fair usage policy applies to “unlimited” data. Roaming subject to partner network coverage. Prices include taxes. Offer valid until 31 December 2025.
\end{textebox}

\begin{popfigure}
\begin{tikzpicture}[
    node distance=1.2cm and 1.5cm,
    box/.style={rectangle, rounded corners, draw=popblue, thick, fill=popblueL, text width=3.2cm, align=center, font=\small},
    arrow/.style={-Stealth, thick, popdark}
]
\node[box, align=center] (start){Client arrives \\ states needs};
\node[box, right=of start, align=center] (present){Agent presents \\ 2 plans \\ (collocations!)};
\node[box, below=of present, align=center] (questions){Client asks \\ $\ge$ 4 questions \\ (price, validity, \\ data, conditions)};
\node[box, right=of questions, align=center] (decide){Client chooses \\ \& subscribes \\ (functional lang.)};
\node[box, below=of decide, align=center] (poster){Group writes \\ promo poster \\ (slogan + lexique)};
\node[box, left=of poster, align=center] (vote){Class votes \\ best offer};
\draw[arrow] (start) -- (present);
\draw[arrow] (present) -- (questions);
\draw[arrow] (questions) -- (decide);
\draw[arrow] (decide) -- (poster);
\draw[arrow] (poster) -- (vote);
\end{tikzpicture}
\legende{Flux de l'activité : du jeu de rôle à la production écrite}
\end{popfigure}

\courssub{Consignes}
\begin{enumerate}
    \item \textbf{Préparation (5 min) :} L'enseignant distribue les fiches \textit{Offer Sheets} aux agents. Les clients lisent la situation et préparent leurs questions. Les observateurs prennent la \textit{Lexical Checklist} (voir Ressources).
    \item \textbf{Jeu de rôle — Phase 1 : Présentation (3–4 min).} L'agent accueille le client, identifie ses besoins (\eng{“What are you looking for in a plan?”}) et présente les deux offres en comparant : data, validité, appels, bonus, rollover. Il utilise au moins 5 collocations cibles.
    \item \textbf{Jeu de rôle — Phase 2 : Questions \& Négociation (4–5 min).} Le client pose \textbf{au moins 4 questions} couvrant : prix exact (taxes comprises?), durée de validité, volume de données réel (throttling?), conditions du bonus/rollover, frais de résiliation, couverture réseau. L'agent répond avec précision. Le client décide et souscrit (\eng{“I’ll go for the Student Boost. How do I activate it?”}).
    \item \textbf{Débriefing oral (3 min).} L'observateur donne son feedback : nombre de mots du champ lexical employés ($\ge$ 10), collocations correctes ($\ge$ 3), calques/faux amis relevés. L'enseignant corrige collectivement les erreurs fréquentes.
    \item \textbf{Production écrite — Affiche publicitaire (15 min).} En groupe, rédigez une \textbf{mini-affiche (6–8 lignes)} en anglais pour le forfait choisi. Contraintes :
    \begin{itemize}
        \item Un \textbf{slogan} court et percutant (ex. \eng{“Learn More, Pay Less”}).
        \item Les infos clés : prix, data, validité, bonus principal.
        \item Au moins \textbf{3 collocations} étudiées (ex. \eng{unlimited high-speed data}, \eng{roll over unused data}, \eng{subscribe today}).
        \item La mention légale \eng{“Terms and conditions apply.”} en bas.
        \item Zéro calque du français (pas \eng{“forfait”}, \eng{“valider”}, \eng{“abonnement”} pour \eng{subscription}).
    \end{itemize}
    \item \textbf{Présentation \& Vote (5 min).} Chaque groupe affiche son poster au tableau. La classe lit, pose une question si besoin, puis vote pour l'offre la plus claire, honnête et attrayante. L'enseignant commente les points forts linguistiques.
\end{enumerate}

\courssub{Ressources à mobiliser}
\begin{definition}[Lexical Checklist — Observer’s Grid (extrait)]
\begin{center}
\begin{tabularx}{\linewidth}{|X|X|X|}
\hline
\rowcolor{popblueL} \textbf{Target Word / Collocation} & \textbf{Used correctly?} & \textbf{Notes / Calque detected} \\
\hline
data plan / subscription plan & \square\ Yes\ \square\ No & \\
\hline
unlimited calls / on-net / off-net & \square\ Yes\ \square\ No & \\
\hline
fair usage policy (FUP) & \square\ Yes\ \square\ No & \\
\hline
validity period / expiry date & \square\ Yes\ \square\ No & \\
\hline
roll over (unused data) & \square\ Yes\ \square\ No & \\
\hline
top-up / recharge (v.) & \square\ Yes\ \square\ No & \\
\hline
subscribe to a plan & \square\ Yes\ \square\ No & \\
\hline
activate a bundle / bonus & \square\ Yes\ \square\ No & \\
\hline
check (my) balance / data usage & \square\ Yes\ \square\ No & \\
\hline
roaming (CEMAC zone) & \square\ Yes\ \square\ No & \\
\hline
throttle / speed cap & \square\ Yes\ \square\ No & \\
\hline
terms and conditions apply & \square\ Yes\ \square\ No & \\
\hline
\end{tabularx}
\end{center}
\end{definition}

\begin{propriete}[Functional Language — Comparing \& Advising]
\begin{itemize}
    \item \eng{“If you mainly use WhatsApp and Zoom, the Student Boost gives you 15 GB plus free educational sites.”}
    \item \eng{“The Pro Connect costs more, but it includes roaming in the CEMAC zone and full data rollover.”}
    \item \eng{“Unlike other operators, we don’t cut you off after the limit; we just reduce the speed.”}
    \item \eng{“You can activate it by dialling *123*1\# or via the MyConnect app.”}
    \item \eng{“Would you like me to walk you through the activation steps?”}
\end{itemize}
\end{propriete}

\begin{attention}[Faux amis \& Calques à bannir]
\begin{center}
\begin{tabularx}{\linewidth}{|X|X|X|}
\hline
\rowcolor{poporangeL} \textbf{Français (calque)} & \textbf{Anglais correct} & \textbf{Remarque} \\
\hline
un forfait & a plan / a package / a bundle & \eng{Forfeit} = perte/déclaration de forfait (sport/juridique) \\
\hline
un abonnement & a subscription (service récurrent) & \eng{Subscription} ≠ \eng{subcription} (orthographe) \\
\hline
valider (un code, une option) & activate / confirm / redeem & \eng{Validate} = valider scientifiquement/légalement \\
\hline
recharger (son crédit) & top up / recharge (v.) & \eng{Recharge} existe mais \cle{top up} plus courant (UK/Afrique) \\
\hline
la validité & validity period / expiry date & pas \eng{validity} seul (trop abstrait) \\
\hline
le report (de data) & rollover (n.) / roll over (v.) & \eng{Report} = rapport / signaler \\
\hline
illimité (data) & unlimited (with FUP) & toujours préciser \eng{fair usage policy} \\
\hline
bonus & bonus / perk / extra & \eng{Bonus} ok, mais \cle{perk} plus marketing \\
\hline
\end{tabularx}
\end{center}
\end{attention}

\begin{methode}[Rédiger une affiche publicitaire efficace (6–8 lignes)]
\begin{enumerate}
    \item \textbf{Accroche (Slogan) :} 3–5 mots, impératif ou promesse. Ex. \eng{“Stay Connected, Stay Ahead.”}
    \item \textbf{Offre clé (Hook) :} Prix + Data + Validité en une phrase. Ex. \eng{“15 GB + Unlimited Calls for 5,500 FCFA / 30 days.”}
    \item \textbf{Bonus distinctif (USP) :} Ce qui différencie. Ex. \eng{“Free educational access — zero data cost.”}
    \item \textbf{Appel à l'action (CTA) :} Verbe d'action + canal. Ex. \eng{“Dial *123*1\# or download the MyConnect app.”}
    \item \textbf{Mention légale :} \eng{“Terms and conditions apply.”} (obligatoire, petit caractère).
    \item \textbf{Relecture :} Vérifier collocations, orthographe (\eng{subscription}, \eng{rollover}), majuscules (noms propres, I), guillemets anglais “ ”.
\end{enumerate}
\end{methode}

\courssub{Proposition de production et corrigé commenté}
\begin{exoresolu}[Modèle d'affiche — Groupe « Student Boost »]
“Learn More, Spend Less” \\
15 GB High-Speed Data + Unlimited On-Net Calls \\
Only 5,500 FCFA — 30 Days Validity \\
Free Access to Educational Platforms (No Data Deducted) \\
Roll Over Unused Data If You Renew On Time \\
Dial *123*1\# or Get the MyConnect App Today \\
Terms and conditions apply.
\tcblower
\textbf{Commentaire linguistique détaillé :}

\begin{enumerate}
    \item \textbf{Slogan :} \eng{“Learn More, Spend Less”} — parallélisme (deux impératifs), vocabulaire simple, promesse claire pour étudiants. Pas de calque.
    \item \textbf{Ligne 2 :} \eng{High-Speed Data} (collocation standard), \eng{Unlimited On-Net Calls} (terme technique précis, \cle{on-net} vs \cle{off-net}). \eng{Unlimited} utilisé honnêtement car le FUP est implicite dans les \eng{Terms}.
    \item \textbf{Ligne 3 :} \eng{Only 5,500 FCFA} — \cle{Only} bien placé avant le prix (focus). \eng{30 Days Validity} (nominalisation correcte, plutôt que \eng{validity period of 30 days} mais acceptable).
    \item \textbf{Ligne 4 :} \eng{Free Access to Educational Platforms} — \cle{Free access to} + nom. Parenthèse \eng{(No Data Deducted)} : \cle{deduct} (v.) technique précis pour “déduire du crédit data”. Évite \eng{“not counted”} (trop vague).
    \item \textbf{Ligne 5 :} \eng{Roll Over Unused Data If You Renew On Time} — \cle{roll over} (verbe phrasal, intransitif ici, sujet = data). \cle{Unused data} (collocation). \cle{Renew on time} (collocation temporelle). Structure conditionnelle réelle (If + present, imperative) adaptée à l'affiche.
    \item \textbf{Ligne 6 :} \eng{Dial *123*1\# or Get the MyConnect App Today} — deux impératifs parallèles. \cle{Get the app} (familier/standard marketing). \cle{Today} (urgence marketing).
    \item \textbf{Ligne 7 :} \eng{Terms and conditions apply.} — mention légale standard, en minuscules sauf \eng{Terms} (début de phrase). Point final.
\end{enumerate}

\textbf{Points de vigilance corrigés par rapport aux erreurs fréquentes :}
\begin{itemize}
    \item \textbf{Orthographe :} \eng{subscription} (pas \eng{subcription}), \eng{rollover} (un mot, nom), \eng{unlimited} (pas \eng{unlimitted}).
    \item \textbf{Faux amis évités :} Pas de \eng{forfait}, \eng{valider}, \eng{abonnement} (sauf si contexte bancaire), \eng{report} pour rollover.
    \item \textbf{Collocations cibles présentes :} \eng{high-speed data}, \eng{unlimited calls}, \eng{free access to}, \eng{roll over unused data}, \eng{renew on time}, \eng{dial a code}, \eng{terms and conditions apply} (7 collocations \textgreater\ 3 exigées).
    \item \textbf{Registre :} Marketing direct, positif, concis. Pas de phrases complètes avec sujet (style télégraphique d'affiche).
\end{itemize}
\end{exoresolu}

\begin{exoresolu}[Modèle d'affiche — Groupe « Pro Connect » (variante)]
“Work Anywhere, Stay Connected” \\
40 GB + Unlimited National Calls — 9,500 FCFA / Month \\
5 GB Roaming Data in CEMAC Zone Included \\
Full Data Rollover — No Expiry On Renewal \\
Priority Support via WhatsApp \\
Activate Now: *123*2\# or MyConnect App \\
Terms and conditions apply.
\tcblower
\textbf{Commentaire :} Slogan ciblé professionnels (\eng{Work Anywhere}). \eng{Roaming Data in CEMAC Zone} (précision géographique). \eng{Full Data Rollover — No Expiry On Renewal} (nominalisé, percutant). \eng{Priority Support} (collocation service client). CTA clair. Respecte toutes les contraintes.
\end{exoresolu}

\courssub{Bilan de l'activité}
\begin{aretenir}[Ce qu'il faut retenir de cette intégration]
\begin{itemize}
    \item Le vocabulaire des télécoms (\cle{plan}, \cle{bundle}, \cle{rollover}, \cle{FUP}, \cle{roaming}, \cle{top-up}) s'emploie dans des \textbf{collocations fixes} : \eng{subscribe to a plan}, \eng{activate a bundle}, \eng{exceed the limit}, \eng{check the balance}, \eng{renew on time}.
    \item Les \textbf{faux amis} sont nombreux : \eng{forfait $\neq$ forfeit}, \eng{valider $\neq$ validate}, \eng{report $\neq$ rollover}, \eng{abonnement $\rightarrow$ subscription}. Le verbe anglais \eng{recharge} existe mais \cle{top up} est la collocation standard en Afrique anglophone et au Royaume-Uni.
    \item À l'oral, les \textbf{formules fonctionnelles} de conseil (\eng{“If you need… I’d recommend…”}), de comparaison (\eng{“Unlike X, Y offers…”}) et de clôture (\eng{“Shall I activate it for you?”}) structurent l'échange professionnel.
    \item À l'écrit (affiche), le \textbf{style télégraphique} (noms, impératifs, phrases nominales) prime. La mention \eng{“Terms and conditions apply.”} est incontournable pour la conformité légale.
    \item Le \textbf{contexte camerounais} (FCFA, CEMAC, MTN/Orange/Camtel, WhatsApp, bendskin, Khan Academy) ancre l'apprentissage dans la réalité vécue par les élèves de Terminale A.
\end{itemize}
\end{aretenir}

\begin{experience}[Prolongement possible — Projet « My Ideal Plan »]
En devoir ou en séance suivante : chaque élève imagine \textbf{son forfait idéal} (prix, data, appels, bonus, validité) pour un étudiant camerounais. Il rédige une \textbf{présentation orale de 2 min} (type pitch startup) et une \textbf{fiche technique en anglais} (tableau + 3 phrases de vente). Les meilleures sont exposées au CDI ou publiées sur le journal du lycée. Critères : lexique cible (15 items), collocations (5), zéro calque, créativité réaliste.
\end{experience}

\begin{savaistu}[Le saviez-vous ? — Telecoms au Cameroun]
Le Cameroun compte plus de \textbf{25 millions d'abonnés mobiles} (données récentes de l'ARTP) pour une population d'environ 28 millions d'habitants. Les deux opérateurs historiques, \textbf{MTN Cameroon} et \textbf{Orange Cameroun}, se partagent ~90\% du marché. \textbf{Camtel} (public) déploie la 4G et la fibre. Le \textbf{Mobile Money} (MoMo, Orange Money) transforme le téléphone en portefeuille bancaire : paiement factures, épargne, crédit. Le régulateur (ARTP) impose l'affichage clair du \textit{Fair Usage Policy} depuis 2021. Le vocabulaire anglais s'impose dans les apps (\textit{MyMTN}, \textit{MyOrange}, \textit{MyConnect}) et les USSD (*123\#, *150\#) — d'où l'importance de le maîtriser pour ne pas se tromper de menu !
\end{savaistu}

\newpage

\coursec{Méthodes et exercices résolus}

\coursec{Lesson 51 — Vocabulary: Words and expressions related to exploiting service packages (telephone/internet services)}

\courssub{Savoir-faire 1 : Identifier et employer le vocabulaire du thème}

\begin{methode}[Identifier et employer le vocabulaire des forfaits téléphone/internet]
Pour maîtriser le lexique des \cle{service packages} (forfaits) au Bac, suis cette démarche :
\begin{enumerate}
\item \textbf{Repère le champ lexical du texte} : souligne tous les mots liés au téléphone (\eng{call, top up, airtime, coverage, network, SIM card}), à Internet (\eng{data, browse, download, Wi-Fi, bandwidth, connection}) et à la facturation (\eng{subscription, bundle, bill, recharge, expiry date}).
\item \textbf{Classe les mots par famille} : verbe (\eng{to subscribe}), nom (\eng{a subscription}), adjectif (\eng{unlimited}), pour enrichir ta réponse.
\item \textbf{Apprends les collocations} : on dit \eng{to top up one's phone} (recharger), \eng{to run out of credit} (être à court de crédit), \eng{to subscribe to a bundle} (souscrire à un forfait). Une collocation exacte vaut plus qu'un mot isolé.
\item \textbf{Vérifie le registre} : \eng{airtime} et \eng{credit} sont courants au Cameroun ; \eng{prepaid} s'oppose à \eng{postpaid}.
\item \textbf{Réemploie le mot dans ta propre phrase} : c'est la meilleure preuve que tu l'as compris. Exemple : \eng{My data bundle expires at midnight.} « Mon forfait internet expire à minuit. »
\end{enumerate}
\end{methode}

\begin{exoresolu}[Vocabulary in context — type Bac]
\textbf{Énoncé.} Read the text and answer the questions.

\eng{“Last week, Amina subscribed to a monthly internet bundle with her mobile operator. The package includes 5 gigabytes of data and unlimited night calls. Unfortunately, the network coverage in her village is poor, so she often runs out of patience before she runs out of data. She plans to switch to another operator that offers cheaper prepaid plans.”}

\begin{enumerate}
\item Find in the text: (a) a word meaning “forfait”, (b) a word meaning “couverture réseau”, (c) a verb meaning “changer d'opérateur”.
\item Is Amina's plan prepaid or postpaid? Justify.
\item Use \eng{to run out of} in your own sentence.
\end{enumerate}

\tcblower
\textbf{Solution détaillée.}

\begin{enumerate}
\item (a) \eng{bundle} ou \eng{package} — les deux signifient « forfait » ; (b) \eng{coverage} (\eng{network coverage} = « couverture réseau ») ; (c) \eng{to switch (to)} = « changer de / passer à ». Attention : \eng{switch} exige la préposition \eng{to} : \eng{switch to another operator}.
\item Her plan is \textbf{postpaid} (or a monthly subscription): the text says she \eng{subscribed to a monthly internet bundle}, which means she pays for a package every month, unlike prepaid credit that you top up when it is finished.
\item Exemple de production : \eng{I always run out of credit before the end of the week because I call my grandmother in Bafoussam every day.} « Je suis toujours à court de crédit avant la fin de la semaine car j'appelle ma grand-mère à Bafoussam tous les jours. » Structure à retenir : \eng{to run out of + nom} (credit, data, airtime, patience).
\end{enumerate}

\textbf{Conclusion.} Identifier le mot juste dans le texte, puis le réemployer dans une phrase personnelle grammaticalement correcte : voilà ce que l'examinateur attend.
\end{exoresolu}

\courssub{Savoir-faire 2 : Réemployer le lexique à l'oral et à l'écrit}

\begin{methode}[Réemployer le lexique dans un dialogue ou une rédaction]
\begin{enumerate}
\item \textbf{Prépare un “coffre à phrases”} pour chaque situation : à la boutique de l'opérateur (\eng{I would like to top up my phone, please.}), au service client (\eng{I am calling to complain about my bill.}), entre amis (\eng{My data is finished, can you share your connection?}).
\item \textbf{Utilise les formules de politesse} : \eng{Could you...?}, \eng{I would like to...}, \eng{How much does the monthly bundle cost?}
\item \textbf{Enchaîne avec des connecteurs} : \eng{first, then, however, finally} pour structurer un récit (\eng{First, I checked my balance; then I bought a data bundle.}).
\item \textbf{Varie les verbes} : ne répète pas \eng{buy} ; utilise \eng{purchase, subscribe to, activate, renew, top up}.
\item \textbf{À l'écrit, vérifie l'accord sujet-verbe} : \eng{The bundle costs 2,000 francs.} (et non \emph{cost}).
\end{enumerate}
\end{methode}

\begin{exoresolu}[Guided dialogue — expression orale/écrite]
\textbf{Énoncé.} Complete the dialogue between a customer and an agent at a mobile phone shop in Douala with words from the box: \emph{balance, top up, bundle, expires, coverage}.

\eng{Agent: Good morning, how can I help you?}
\eng{Customer: I would like to (1) \dots\ my phone with 1,000 francs, please.}
\eng{Agent: Would you like airtime or a data (2) \dots ?}
\eng{Customer: A data one. How long is it valid?}
\eng{Agent: It (3) \dots\ after seven days.}
\eng{Customer: And how can I check my (4) \dots ?}
\eng{Agent: Just dial the short code. By the way, our network (5) \dots\ is excellent in Douala.}

\tcblower
\textbf{Solution détaillée.}

\begin{enumerate}
\item \eng{top up} — c'est la collocation exacte pour « recharger » un téléphone ; \eng{I would like to + base verbale}.
\item \eng{bundle} — \eng{a data bundle} = « un forfait internet » ; l'article \eng{a} impose un nom singulier.
\item \eng{expires} — sujet \eng{it} à la 3\textsuperscript{e} personne du singulier au présent simple, d'où le \textbf{-s} obligatoire.
\item \eng{balance} — \eng{to check one's balance} = « consulter son solde » ; possessif \eng{my} + nom.
\item \eng{coverage} — \eng{network coverage} = « couverture réseau » ; invariable ici car nom non comptable dans ce sens.
\end{enumerate}

Dialogue complet : \eng{I would like to top up my phone... a data bundle... It expires after seven days... check my balance... our network coverage is excellent.}

\textbf{Conclusion.} Chaque case se justifie par une collocation apprise : c'est la logique du lexique, pas le hasard.
\end{exoresolu}

\courssub{Savoir-faire 3 : Éviter les faux amis et les calques du français}

\begin{methode}[Déjouer les faux amis du domaine téléphone/internet]
\begin{enumerate}
\item \textbf{Méfie-toi des mots “transparents”} : \eng{actually} ne veut pas dire « actuellement » (dis \eng{currently}) ; \eng{a delay} = « un retard », pas « un délai » au sens de « date limite » (dis \eng{a deadline}).
\item \textbf{Ne calque pas les structures françaises} : « j'ai consommé mes données » ne se traduit pas mot à mot ; dis \eng{I have used up my data}.
\item \textbf{Attention aux prépositions} : \eng{to subscribe \textbf{to} a package}, \eng{to pay \textbf{for} a service}, \eng{to complain \textbf{about} a bill}, \eng{to connect \textbf{to} the internet}.
\item \textbf{« Recharger » n'est pas} \emph{to recharge} au sens courant du crédit téléphonique : la forme la plus naturelle est \eng{to top up} ; \eng{to recharge} convient pour une batterie (\eng{recharge the battery}).
\item \textbf{Vérifie toujours le mot dans une phrase complète} avant de l'écrire dans ta copie.
\end{enumerate}
\end{methode}

\begin{exoresolu}[Faux amis et calques — correction de phrases]
\textbf{Énoncé.} Each sentence contains a mistake (faux ami, calque or wrong preposition). Correct it.

\begin{enumerate}
\item I want to subscribe at the new internet package.
\item Actually, I don't have any credit on my phone. (sens : « en ce moment »)
\item She complained for her bill at the agency.
\item My phone is discharged, I must top up the battery.
\end{enumerate}

\tcblower
\textbf{Solution détaillée.}

\begin{enumerate}
\item \eng{I want to subscribe \textbf{to} the new internet package.} — Le verbe \eng{subscribe} est toujours suivi de la préposition \eng{to} ; « s'abonner à » se traduit par \eng{subscribe to}, jamais \emph{subscribe at}.
\item \eng{\textbf{Currently}, I don't have any credit on my phone.} — Faux ami classique : \eng{actually} signifie « en fait, en réalité » ; « actuellement » se dit \eng{currently} ou \eng{at the moment}.
\item \eng{She complained \textbf{about} her bill at the agency.} — \eng{to complain about something} = « se plaindre de quelque chose » ; la préposition \eng{for} est un calque du français « se plaindre pour ».
\item \eng{My phone is \textbf{dead}, I must \textbf{recharge} the battery.} — On dit \eng{my phone is dead} ou \eng{the battery is flat} pour « déchargé » ; \eng{to top up} se réserve au crédit téléphonique, \eng{to recharge} à la batterie.
\end{enumerate}

\textbf{Conclusion.} Chaque erreur vient d'un calque du français ; la solution repose sur la collocation anglaise authentique.
\end{exoresolu}

\courssub{Savoir-faire 4 : Exploiter les familles de mots et la formation des mots}

\begin{methode}[Construire et reconnaître les mots de la même famille]
\begin{enumerate}
\item \textbf{Identifie la racine} : \eng{connect} donne \eng{connection} (nom), \eng{connected} (adjectif/participe), \eng{disconnect} (verbe avec préfixe négatif \eng{dis-}), \eng{reconnect} (préfixe \eng{re-} = « de nouveau »).
\item \textbf{Retiens les suffixes fréquents} : \eng{-tion} (nom d'action : \eng{subscription, connection}), \eng{-er/-or} (personne ou appareil : \eng{operator, caller}), \eng{-less} (privation : \eng{wireless} = « sans fil »), \eng{-able} (possibilité : \eng{affordable} = « abordable »).
\item \textbf{Utilise les préfixes} : \eng{un-} (\eng{unlimited} = « illimité »), \eng{pre-} (\eng{prepaid} = « prépayé »), \eng{post-} (\eng{postpaid} = « postpayé »), \eng{re-} (\eng{renew} = « renouveler »).
\item \textbf{Pour les exercices de transformation} : repère la nature grammaticale exigée par la phrase (verbe après un sujet, nom après un article, adjectif devant un nom), puis dérive la racine correctement.
\item \textbf{Orthographe} : double consonne parfois (\eng{calling, dropped}), \eng{-y} devient \eng{-i} (\eng{connectivity}).
\end{enumerate}
\end{methode}

\begin{exoresolu}[Word formation — type Bac]
\textbf{Énoncé.} Complete each sentence with the correct form of the word in brackets.

\begin{enumerate}
\item The \dots\ of this bundle is very easy. (activate)
\item This operator offers \dots\ calls at night. (limit)
\item The \dots\ in rural areas is often slow. (connect)
\item A \dots\ router allows you to browse without cables. (wire)
\item These data plans are not \dots\ for students. (afford)
\end{enumerate}

\tcblower
\textbf{Solution détaillée.}

\begin{enumerate}
\item \eng{activation} — après l'article \eng{the}, il faut un nom ; suffixe \eng{-tion} : \eng{activate} $\rightarrow$ \eng{activation}. « L'activation de ce forfait est très facile. »
\item \eng{unlimited} — devant le nom \eng{calls}, il faut un adjectif ; le contexte (« at night ») suggère le sens « illimité », d'où le préfixe \eng{un-} : \eng{unlimited calls}.
\item \eng{connection} — après \eng{the}, un nom ; \eng{connect} $\rightarrow$ \eng{connection}. « La connexion dans les zones rurales est souvent lente. »
\item \eng{wireless} — adjectif devant \eng{router} ; suffixe de privation \eng{-less} : \eng{a wireless router} = « un routeur sans fil ».
\item \eng{affordable} — après \eng{are not}, un adjectif ; suffixe \eng{-able} : \eng{affordable} = « abordable ». « Ces forfaits ne sont pas abordables pour les élèves. »
\end{enumerate}

\textbf{Conclusion.} La méthode gagnante : déterminer d'abord la nature du mot attendu, puis appliquer le suffixe ou le préfixe adapté.
\end{exoresolu}

\begin{attention}[Les 5 erreurs qui coûtent des points au Bac]
\begin{enumerate}
\item \textbf{Faux ami \eng{actually}} : écrire \emph{Actually, I have no data} pour « actuellement » est une faute. Dis \eng{Currently} ou \eng{At the moment}. \eng{Actually} = « en fait ».
\item \textbf{Préposition oubliée après \eng{subscribe}} : \emph{I subscribed the bundle} est incorrect. La seule forme correcte est \eng{I subscribed \textbf{to} the bundle.}
\item \textbf{Calque de « recharger »} : \emph{I want to recharge my phone with 500 francs} est un calque ; pour le crédit, dis \eng{to top up my phone}. Garde \eng{recharge} pour la batterie.
\item \textbf{Oubli du -s à la 3\textsuperscript{e} personne} : \emph{The bundle expire after seven days} est fautif ; écris \eng{The bundle expire\textbf{s} after seven days.} Le présent simple exige le \textbf{-s} avec \eng{he/she/it}.
\item \textbf{Orthographe des mots dérivés} : \emph{conection}, \emph{subscribtion}, \emph{unlimitted} sont des fautes fréquentes. Écris \eng{connection}, \eng{subscription}, \eng{unlimited} (un seul \textbf{t}).
\end{enumerate}
\end{attention}

\newpage

\coursec{Exercices}

\coursec{Lesson 51 — Vocabulary: Words and expressions related to exploiting service packages (telephone/internet services)}

\begin{textebox}[Situation de communication : À l'agence MTN de Buea]
“Good morning, madam. I’d like to \cle{subscribe} to a new \cle{data bundle}. My current \cle{plan} \cle{expires} tomorrow.”
“Certainly. We have a \cle{monthly package} with 20 GB, \cle{valid} for 30 days. It includes \cle{rollover} for unused data.”
“Does it have a \cle{fair usage policy}? I hate \cle{throttling}.”
“Yes, speeds reduce after 5 GB daily. You can \cle{top up} via \cle{Mobile Money} anytime.”
“Perfect. I’ll \cle{activate} it now. Here is my \cle{reference number}.”
\end{textebox}

\courssub{Champ lexical 1 : Téléphone et appels (Phone \& Calls)}

\begin{definition}[Vocabulaire clé : Voix et crédit]
Le tableau ci-dessous regroupe les termes essentiels pour parler des appels, du crédit et de la recharge au Cameroun.
\begin{center}
\begin{tabularx}{\linewidth}{|X|X|X|}
\hline
\rowcolor{popblueL} \textbf{English} & \textbf{Français} & \textbf{Contexte / Collocation} \\
\hline
\cle{airtime} & crédit d'appel / temps d'antenne & \eng{buy airtime}, \eng{airtime balance} \\
\cle{credit} & crédit (monétaire) & \eng{credit balance}, \eng{out of credit} \\
\cle{top up} (v.) / \cle{top-up} (n.) & recharger / recharge & \eng{top up my phone}, \eng{a top-up card} \\
\cle{recharge} (v./n.) & recharger / recharge & \eng{recharge my account}, \eng{recharge voucher} \\
\cle{balance} & solde & \eng{check my balance}, \eng{low balance} \\
\cle{USSD code} & code USSD & \eng{dial *157\#} \\
\cle{voice bundle} / \cle{call package} & forfait appels & \eng{subscribe to a voice bundle} \\
\cle{outgoing calls} & appels sortants & \eng{make outgoing calls} \\
\cle{incoming calls} & appels entrants & \eng{receive incoming calls} \\
\hline
\end{tabularx}
\end{center}
\end{definition}

\begin{propriete}[Top up vs Recharge]
\begin{itemize}
    \item \cle{Top up} (souvent séparable : \eng{top my phone up}) est le terme le plus courant au Royaume-Uni et en Afrique anglophone pour ajouter du crédit prépayé.
    \item \cle{Recharge} est plus formel, utilisé dans les interfaces bancaires ou applications (ex. \eng{recharge via Mobile Money}).
    \item \cle{Airtime} désigne spécifiquement le crédit pour les appels vocaux (par opposition à \cle{data} pour internet).
\end{itemize}
\end{propriete}

\courssub{Champ lexical 2 : Internet et données (Internet \& Data)}

\begin{definition}[Vocabulaire clé : Connexion et volumes]
\begin{center}
\begin{tabularx}{\linewidth}{|X|X|X|}
\hline
\rowcolor{popblueL} \textbf{English} & \textbf{Français} & \textbf{Contexte / Collocation} \\
\hline
\cle{data} / \cle{mobile data} & données mobiles / internet mobile & \cle{data consumption}, \cle{data usage} \\
\cle{data bundle} / \cle{data pack} & forfait data / paquet data & \eng{buy a 5 GB data bundle} \\
\cle{megabyte (MB)} / \cle{gigabyte (GB)} & mégaoctet / gigaoctet & Unités de volume \\
\cle{high-speed} / \cle{4G} / \cle{5G} & haut débit / 4G / 5G & \eng{high-speed internet} \\
\cle{throttling} & bridage / réduction de débit & \eng{speed throttling after 2 GB} \\
\cle{fair usage policy (FUP)} & politique d'utilisation équitable & \eng{FUP limit reached} \\
\cle{network congestion} & congestion du réseau & \eng{during peak hours} \\
\cle{browse} / \cle{surf} (v.) & naviguer & \eng{browse the web} \\
\cle{download} / \cle{upload} (v./n.) & télécharger / envoyer (fichiers) & \eng{download speed} \\
\cle{streaming} & lecture en continu & \eng{video streaming} \\
\hline
\end{tabularx}
\end{center}
\end{definition}

\begin{attention}[Faux ami : \cle{Consume} vs \cle{Use up} / \cle{Run out of}]
En français, on dit “j'ai \textbf{consommé} mon forfait”. En anglais, \eng{consume} est technique (consommation énergétique). Pour un forfait, on utilise :
\begin{itemize}
    \item \eng{use up} (phrasal verb) : \eng{I used up all my data.}
    \item \eng{run out of} (phrasal verb) : \eng{I ran out of data yesterday.}
    \item \eng{exhaust} (formel) : \eng{He exhausted his monthly allowance.}
\end{itemize}
Ne dites pas : \eng{*I consumed my bundle.}
\end{attention}

\courssub{Champ lexical 3 : Forfaits, abonnements et facturation (Plans, Subscriptions \& Billing)}

\begin{definition}[Vocabulaire clé : Cycle de vie d'un forfait]
\begin{center}
\begin{tabularx}{\linewidth}{|X|X|X|}
\hline
\rowcolor{popblueL} \textbf{English} & \textbf{Français} & \textbf{Contexte / Collocation} \\
\hline
\cle{plan} / \cle{package} / \cle{bundle} / \cle{pack} & forfait / offre / pack & \eng{a monthly plan}, \eng{a Social Pack} \\
\cle{subscribe (to)} & s'abonner (à) & \eng{subscribe to a plan} \\
\cle{subscription} & abonnement & \eng{monthly subscription}, \eng{cancel subscription} \\
\cle{subscriber} & abonné & \eng{a loyal subscriber} \\
\cle{validity} / \cle{validity period} & validité / durée de validité & \eng{valid for 30 days}, \eng{expiry date} \\
\cle{expire} (v.) / \cle{expiry} (n.) & expirer / expiration & \eng{The plan expires tomorrow.} \\
\cle{rollover} (n./v.) & report (données non utilisées) & \eng{data rollover}, \eng{unused data rolls over} \\
\cle{auto-renewal} & renouvellement automatique & \eng{disable auto-renewal} \\
\cle{terms and conditions} / \cle{T\&Cs} & conditions générales & \eng{read the T\&Cs} \\
\cle{hidden charges} & frais cachés & \eng{avoid hidden charges} \\
\cle{bill shock} & choc de la facture / surprise facture & \eng{prevent bill shock} \\
\cle{per-megabyte cost} / \cle{unit price} & coût au mégaoctet / prix unitaire & \eng{compare per-MB costs} \\
\cle{regulator} (ART) & régulateur (ART) & \eng{the regulator enforces rules} \\
\hline
\end{tabularx}
\end{center}
\end{definition}

\begin{savaistu}[Le marché télécoms au Cameroun]
Le Cameroun compte deux opérateurs majeurs : \textbf{MTN Cameroon} et \textbf{Orange Cameroon}. Le régulateur est l'\textbf{ART} (Agence de Régulation des Télécommunications). L'usage du \textbf{Mobile Money} (MTN MoMo, Orange Money) pour l'achat de \cle{airtime} et de \cle{data bundles} via codes \cle{USSD} (ex. \eng{*157\#}, \eng{*141\#}) est la norme. Les offres “\cle{Night Plan}” (souvent 23h--5h) et les “\cle{Social Packs}” (WhatsApp, Facebook illimités) sont très populaires chez les étudiants.
\end{savaistu}

\courssub{Faux amis et calques du français : Tableau récapitulatif}

\begin{attention}[Erreurs fréquentes des francophones -- À éviter absolument]
\begin{center}
\begin{tabularx}{\linewidth}{|X|X|X|}
\hline
\rowcolor{popblueL} \textbf{Français (Calque)} & \textbf{Anglais INCORRECT} & \textbf{Anglais CORRECT} \\
\hline
Mon forfait est fini & \eng{My forfait is finished} & \eng{My \cle{bundle/plan} has \cle{expired} / \cle{run out}} \\
J'ai consommé ma data & \eng{I consumed my data} & \eng{I \cle{used up} / \cle{ran out of} my data} \\
Je veux faire un abonnement & \eng{I want to make an abonnement} & \eng{I want to \cle{subscribe to} / \cle{take out} a \cle{plan}} \\
Le tarif est cher & \eng{The tarif is expensive} & \eng{The \cle{rate} / \cle{price} / \cle{fee} is high} \\
Recharger avec 1000 francs & \eng{Recharge with 1000 francs} & \eng{\cle{Top up} \cle{credit/airtime} with 1\,000 FCFA} \\
Le réseau ralentit & \eng{The network is slowing} & \eng{The network \cle{is slow} / \cle{is throttling}} \\
\hline
\end{tabularx}
\end{center}
\end{attention}

\courssub{Mots de la même famille et formation des mots (Word Formation)}

\begin{propriete}[Dérivation : Verbe $\rightarrow$ Noms / Adjectifs]
La maîtrise des familles lexicales permet de varier l'expression (ex. \eng{subscribe} $\rightarrow$ \eng{subscription}, \eng{subscriber}, \eng{subscribed}).
\begin{center}
\begin{tabularx}{\linewidth}{|X|X|X|X|}
\hline
\rowcolor{popblueL} \textbf{Verbe} & \textbf{Nom (chose/concept)} & \textbf{Nom (agent)} & \textbf{Adjectif / Participe} \\
\hline
subscribe & subscription & subscriber & subscribed \\
connect & connection & connector / user & connected \\
validate & validation & validator & valid / validated \\
activate & activation & activator & active / activated \\
consume & consumption & consumer & consuming / consumable \\
expire & expiry / expiration & -- & expired \\
register & registration & registrant & registered \\
\hline
\end{tabularx}
\end{center}
\end{propriete}

\begin{exemplebox}[Exemples en contexte]
\eng{The \textbf{activation} of your \textbf{subscription} is instant.} (L'activation de votre abonnement est instantanée.) \\
\eng{A \textbf{registered} \textbf{subscriber} enjoys \textbf{rollover} benefits.} (Un abonné enregistré bénéficie du report de données.) \\
\eng{Check the \textbf{validity} period before you \textbf{validate} the transaction.} (Vérifiez la validité avant de valider.)
\end{exemplebox}

\courssub{Le lexique en contexte : Dialogue authentique et écoute}

\begin{textebox}[Script d'écoute : Choosing a Student Plan at Orange Shop -- Douala]
\textbf{Agent:} “Welcome to Orange. How can I help you today?”
\textbf{Student:} “Hi. I need a \cle{data bundle} for my studies. I mostly do research and download PDFs.”
\textbf{Agent:} “We have the \textbf{'Orange Edu Bundle'}: 15 GB for 30 days at 3\,500 FCFA. The \cle{validity} is one month.”
\textbf{Student:} “What happens if I don't finish the 15 GB?”
\textbf{Agent:} “It has \cle{rollover}. If you \cle{renew} before the \cle{expiry} date, unused GBs move to the next cycle.”
\textbf{Student:} “And the speed? Is there \cle{throttling}?”
\textbf{Agent:} “Standard speed up to 4G+. \cle{Fair usage policy} applies only after 3 GB daily — then speed drops to 256 kbps.”
\textbf{Student:} “Okay. I'll \cle{subscribe}. I'll pay via \cle{Orange Money}.”
\textbf{Agent:} “Dial \eng{*141\#}, select ‘Buy Bundle’, then ‘Edu Pack’. You'll get a confirmation SMS with your new \cle{balance}.”
\end{textebox}

\begin{experience}[Jeu de rôle : À l'agence de télécoms]
\textbf{Consigne :} En binôme. L'élève A est le client (étudiant à Yaoundé), l'élève B est l'agent (MTN ou Orange).
\textbf{Client :} Vous voulez un forfait 30 jours avec rollover, pour réviser le Bac (recherche, vidéos YouTube éducatives). Budget : 5\,000 FCFA. Posez des questions sur : validité, rollover, FUP/throttling, recharge Mobile Money.
\textbf{Agent :} Présentez deux offres imaginaires (ex. \eng{MTN Student Pro} vs \eng{MTN Night Owl}). Expliquez les différences (prix, volume, validité, rollover, FUP). Incitez à lire les T\&Cs.
\textbf{Durée :} 3 min. Changez de rôle.
\end{experience}

\courssub{Je vérifie mes connaissances}

\begin{exercice}[QCM : Vocabulaire des services télécoms]
Choisissez la bonne réponse parmi les quatre propositions. Une seule réponse est correcte.
\begin{enumerate}
    \item I need to \cle{top up} my phone; I have no \cle{airtime} left.
    \begin{enumerate}
        \item consume
        \item \textbf{top up}
        \item forfeit
        \item subscribe
    \end{enumerate}
    \item The \cle{data bundle} I bought yesterday \cle{expires} at midnight tonight.
    \begin{enumerate}
        \item validates
        \item \textbf{expires}
        \item connects
        \item rolls over
    \end{enumerate}
    \item Before you \cle{subscribe} to a new plan, always \cle{check} the \cle{terms and conditions}.
    \begin{enumerate}
        \item sign / read
        \item \textbf{subscribe / check}
        \item buy / ignore
        \item cancel / verify
    \end{enumerate}
    \item My \cle{balance} is low; I cannot make any outgoing calls until I \cle{recharge}.
    \begin{enumerate}
        \item credit / consume
        \item \textbf{balance / recharge}
        \item bonus / forfeit
        \item package / browse
    \end{enumerate}
    \item The \cle{rollover} feature allows unused data to be carried over to the next billing \cle{cycle}.
    \begin{enumerate}
        \item \textbf{rollover / cycle}
        \item bonus / month
        \item refund / period
        \item transfer / year
    \end{enumerate}
\end{enumerate}
\end{exercice}

\begin{exercice}[Vrai ou Faux justifié : Offres internet au Cameroun]
Lisez le texte ci-dessous et dites si les affirmations sont vraies (V) ou fausses (F). Justifiez chaque réponse par une phrase extraite du texte.

\begin{textebox}[Article : Mobile Data Trends in Cameroon]
Mobile internet penetration has grown rapidly in Yaoundé and Douala over the past five years. Major operators like MTN Cameroon and Orange Cameroon compete aggressively with affordable data bundles. Customers can now purchase daily, weekly, or monthly packages via USSD codes or Mobile Money apps. A popular feature is the “Night Plan”, which offers high-speed data at a reduced rate between 11 PM and 5 AM. However, many users complain about \cle{throttling} speeds once the \cle{fair usage policy} limit is reached. To avoid bill shock, subscribers are advised to monitor their \cle{data consumption} regularly using the operator's self-care app.
\end{textebox}

\begin{enumerate}
    \item Mobile internet growth is limited to rural areas only.
    \item Operators offer packages with different validity periods (daily, weekly, monthly).
    \item The “Night Plan” is more expensive than standard daytime rates.
    \item \cle{Throttling} occurs when a user exceeds the fair usage policy limit.
    \item Subscribers cannot check their data usage on their phones.
\end{enumerate}
\end{exercice}

\begin{exercice}[Vocabulaire : Compléter le dialogue à l'agence]
Complétez le dialogue suivant en utilisant les mots de la liste ci-dessous. Chaque mot s'utilise une seule fois.
\medskip
\noindent \textbf{Liste :} \cle{bundle} -- \cle{valid} -- \cle{balance} -- \cle{subscribe} -- \cle{rollover} -- \cle{expires} -- \cle{airtime} -- \cle{recharge} -- \cle{package} -- \cle{data}

\begin{textebox}[Dialogue at the Telecom Agency]
\textbf{Agent:} “Good morning! How can I help you?”
\textbf{Customer:} “Hello. I want to \_\_\_\_\_\_\_\_\_\_ to a new internet \_\_\_\_\_\_\_\_\_\_ for my smartphone.”
\textbf{Agent:} “We have a 10 GB \_\_\_\_\_\_\_\_\_\_ that is \_\_\_\_\_\_\_\_\_\_ for 30 days. It costs 5\,000 FCFA.”
\textbf{Customer:} “Does it have the \_\_\_\_\_\_\_\_\_\_ feature? I hate losing unused megabytes at the end of the month.”
\textbf{Agent:} “Yes, sir. Unused data rolls over to the next cycle if you \_\_\_\_\_\_\_\_\_\_ before the current plan \_\_\_\_\_\_\_\_\_\_.”
\textbf{Customer:} “Perfect. I'll take it. How do I check my \_\_\_\_\_\_\_\_\_\_ later?”
\textbf{Agent:} “Dial *157\# or use the MyMTN app. You can also buy \_\_\_\_\_\_\_\_\_\_ for voice calls there.”
\end{textebox}
\end{exercice}

\begin{exercice}[Formation des mots : Familles lexicales]
Complétez le tableau suivant en trouvant la forme manquante (nom, verbe, adjectif ou adverbe). La première ligne est donnée en exemple.

\begin{center}
\begin{tabularx}{\linewidth}{|X|X|X|X|}
\hline
\rowcolor{popblueL} \textbf{Verbe} & \textbf{Nom (chose/concept)} & \textbf{Nom (agent)} & \textbf{Adjectif / Adverbe} \\
\hline
subscribe & subscription & subscriber & \textbf{subscribed} \\
\hline
connect & \textbf{connection} & \textbf{connector / user} & connected \\
\hline
validate & validation & \textbf{validator} & \textbf{valid / validated} \\
\hline
activate & \textbf{activation} & activator & \textbf{active / activated} \\
\hline
consume & consumption & consumer & \textbf{consuming / consumable} \\
\hline
expire & \textbf{expiry / expiration} & -- & \textbf{expired} \\
\hline
register & registration & \textbf{registrant} & registered \\
\hline
\end{tabularx}
\end{center}
\end{exercice}

\courssub{Je m'entraîne}

\begin{exercice}[Traduction : Du français vers l'anglais]
Traduisez les phrases suivantes en anglais en réutilisant le vocabulaire de la leçon.
\begin{enumerate}
    \item J'ai épuisé mon forfait data hier soir ; je dois le recharger.
    \item Ce forfait n'est plus valable depuis le 1er janvier.
    \item As-tu activé l'option de report des données (rollover) sur ta ligne ?
    \item Vérifie ton solde avant d'appeler, tu n'as plus de crédit.
    \item Les opérateurs proposent des offres promotionnelles le week-end.
    \item Je me suis abonné à un forfait illimité pour les réseaux sociaux.
\end{enumerate}
\end{exercice}

\begin{exercice}[Transformation de phrases : Calques et corrections]
Les phrases ci-dessous contiennent des erreurs typiques de francophones (calques, faux amis). Réécrivez-les correctement en anglais.
\begin{enumerate}
    \item I \textbf{consumed} my \textbf{forfait} yesterday. (Indice : \cle{use up} / \cle{run out of} / \cle{bundle})
    \item My internet \textbf{finished} at midnight. (Indice : \cle{expire} / \cle{run out})
    \item I want to \textbf{make} a new \textbf{abonnement}. (Indice : \cle{take out} / \cle{subscribe to} / \cle{plan})
    \item The \textbf{tarif} is too expensive for me. (Indice : \cle{rate} / \cle{price} / \cle{fee})
    \item He \textbf{recharged} his phone with 1000 francs. (Indice : \cle{top up} / \cle{credit} / \cle{airtime})
    \item The network is \textbf{slowing} since morning. (Indice : \cle{throttling} / \cle{slow} / \cle{since this morning})
\end{enumerate}
\end{exercice}

\begin{exercice}[Texte à trous : Choisir le bon mot]
Complétez le texte avec les mots appropriés parmi les propositions entre parenthèses.

\begin{textebox}[Managing Your Mobile Account]
Managing a mobile account in Cameroon has become easier with USSD codes. To check your \_\_\_\_\_\_\_\_\_\_ (balance / validity / bundle), dial *157\#. If you want to \_\_\_\_\_\_\_\_\_\_ (subscribe / register / activate) to a weekly data plan, select “Data Bundles” then “Weekly”. Make sure you have enough \_\_\_\_\_\_\_\_\_\_ (airtime / credit / bonus) on your account before confirming. Once the transaction is successful, you will receive an SMS confirming the \_\_\_\_\_\_\_\_\_\_ (activation / expiration / consumption). Remember that the \_\_\_\_\_\_\_\_\_\_ (validity / cycle / rollover) period starts immediately. If you don't use all your data, check if your plan allows \_\_\_\_\_\_\_\_\_\_ (rollover / throttling / forfeit) to the next month. To avoid \_\_\_\_\_\_\_\_\_\_ (bill shock / fair usage / subscription), monitor your usage daily.
\end{textebox}
\end{exercice}

\begin{exercice}[Appariement : Collocations télécoms]
Associez chaque verbe de la colonne A à l'élément de la colonne B pour former une collocation correcte. Écrivez les paires (ex. 1-f).

\begin{center}
\begin{tabularx}{\linewidth}{|X|X|}
\hline
\rowcolor{popblueL} \textbf{Colonne A (Verbes)} & \textbf{Colonne B (Noms/Compléments)} \\
\hline
1. top up & a. the internet \\
2. run out of & b. my balance \\
3. check & c. a service \\
4. activate & d. credit / airtime \\
5. browse & e. data / credit \\
\hline
\end{tabularx}
\end{center}

\textbf{Écrivez ensuite une phrase originale pour chaque collocation.}
\end{exercice}

\courssub{Je me prépare au Bac}

\begin{exercice}[Compréhension de texte et questions de langue -- Type Bac]
Lisez attentivement le texte suivant et répondez aux questions qui suivent.

\begin{textebox}[Mobile Money and Data Bundles: Changing Communication in Cameroon]
In the bustling streets of Douala, street vendors like Pierre no longer need to carry stacks of physical recharge cards. A simple transaction on his smartphone via Mobile Money — MTN MoMo or Orange Money — instantly credits his line with airtime or a data bundle. This digital shift has revolutionized how Cameroonians \cle{exploit} service packages.

Operators now design \cle{tailored} offers: “Social Packs” for WhatsApp and Facebook addicts, “Night Owl” bundles for students downloading heavy files, and “Family Share” plans allowing data \cle{pooling} across multiple numbers. The \cle{validity} periods range from 24 hours to 30 days, with \cle{rollover} options rewarding loyal customers who renew before \cle{expiry}.

However, challenges persist. \cle{Network congestion} during peak hours leads to \cle{throttling}, frustrating users who have paid for “high speed”. The \cle{fair usage policy} (FUP), often hidden in fine print, caps high-speed data at 2GB or 5GB per day on so-called “unlimited” plans. Consumer associations urge the \cle{regulator} (ART) to enforce transparency on \cle{hidden charges} and \cle{auto-renewal} traps.

For the savvy user, the strategy is clear: compare \cle{per-megabyte} costs, disable auto-renewal via USSD, use Wi-Fi for heavy updates, and always read the \cle{terms and conditions} before clicking “Subscribe”.
\end{textebox}

\textbf{Questions de compréhension (Réponses courtes attendues) :}
\begin{enumerate}
    \item How do street vendors like Pierre buy airtime today?
    \item Name two types of tailored data offers mentioned in the text.
    \item What happens to internet speed during peak hours according to the text?
    \item What does the Fair Usage Policy (FUP) do on “unlimited” plans?
    \item What advice does the text give to avoid auto-renewal traps?
\end{enumerate}

\textbf{Questions de langue :}
\begin{enumerate}
    \setcounter{enumi}{5}
    \item Find in the text a synonym for “customised” (paragraph 2).
    \item Find in the text the opposite of “visible / obvious” (paragraph 3).
    \item Rewrite the sentence “Consumers urge the regulator to enforce transparency” in the passive voice.
    \item Put the verbs in brackets in the correct tense: “If Pierre \_\_\_\_\_\_\_\_\_\_ (not / disable) auto-renewal, he \_\_\_\_\_\_\_\_\_\_ (lose) credit automatically next week.”
    \item Explain the meaning of the phrasal verb “\cle{run out of}” in this context: “He ran out of data yesterday.”
\end{enumerate}
\end{exercice}

\begin{exercice}[Traduction : Thème -- Type Bac]
Traduisez le passage suivant en anglais. Utilisez le vocabulaire précis de la leçon (forfaits, data, recharge, validité, report, abonnement, solde).

\begin{quote}
« Au Cameroun, choisir un forfait internet n'est pas facile. Hier, j'ai voulu m'abonner à un forfait mensuel de 20 Go, mais mon solde était insuffisant. J'ai dû recharger mon compte via Mobile Money. Le vendeur m'a expliqué que ce forfait avait une validité de 30 jours et que l'option de report des données non utilisées était activée automatiquement. Cependant, il m'a prévenu que la vitesse serait réduite après 5 Go par jour à cause de la politique d'utilisation équitable. Je vais comparer les offres des deux opérateurs avant de prendre une décision finale. »
\end{quote}
\end{exercice}

\begin{exercice}[Expression écrite : Composition -- Type Bac]
\textbf{Sujet :} Your friend wants to choose an internet plan for his/her studies. Write a paragraph advising him/her, comparing two specific offers (real or imaginary names like “MTN Student Pack” vs “Orange Edu Bundle”) and using \textbf{at least eight words or expressions} from the lesson vocabulary list (e.g., \cle{bundle}, \cle{validity}, \cle{rollover}, \cle{throttling}, \cle{subscribe}, \cle{airtime}, \cle{fair usage policy}, \cle{top up}, \cle{balance}, \cle{auto-renewal}).

\textbf{Contraintes :}
\begin{itemize}
    \item Longueur : 80--100 mots.
    \item Structure : Introduction (contexte), Comparaison (prix, volume, validité, report), Conseil final.
    \item Vocabulaire : Soulignez les 8 mots imposés utilisés.
    \item Registre : Semi-formel (conseil entre amis).
\end{itemize}
\end{exercice}

\begin{exemplebox}[Modèle de réponse guidée (environ 90 mots)]
“Hey! For your studies, compare the \textbf{\cle{MTN Student Pack}} (15 GB, 30-day \textbf{\cle{validity}}, 4\,500 FCFA) and the \textbf{\cle{Orange Edu Bundle}} (10 GB + unlimited Wikipedia, 30 days, 3\,500 FCFA). Both offer \textbf{\cle{rollover}} if you \textbf{\cle{renew}} before \textbf{\cle{expiry}}. Watch out: MTN applies a strict \textbf{\cle{fair usage policy}} (throttling after 3 GB/day), while Orange \textbf{\cle{throttling}} starts at 2 GB. \textbf{\cle{Subscribe}} to Orange if budget is tight, but \textbf{\cle{top up}} your \textbf{\cle{balance}} via Mobile Money. Disable \textbf{\cle{auto-renewal}} to control costs. Read \textbf{\cle{T\&Cs}} first!”
\end{exemplebox}

\courssub{Corrigés des exercices d'entraînement et de révision}
% Les corrigés sont intégrés après chaque exercice ci-dessus pour un feedback immédiat.
% Cette section peut rester vide ou servir de sommaire.
\begin{aretenir}[Check-list vocabulaire Lesson 51]
\begin{itemize}
    \item \textbf{Verbes clés :} \eng{subscribe to}, \eng{top up}, \eng{recharge}, \eng{activate}, \eng{renew}, \eng{expire}, \eng{run out of}, \eng{use up}, \eng{throttle}.
    \item \textbf{Noms clés :} \eng{bundle/pack/plan}, \eng{airtime/credit}, \eng{balance}, \eng{validity/expiry}, \eng{rollover}, \eng{FUP/fair usage policy}, \eng{auto-renewal}, \eng{terms and conditions}, \eng{bill shock}.
    \item \textbf{Adjectifs :} \eng{valid}, \eng{active}, \eng{unlimited}, \eng{tailored}, \eng{hidden}.
    \item \textbf{Prononciation :} \eng{validity} [va-li-di-ti], \eng{expiry} [eks-pai-ri], \eng{throttling} [trot-ling], \eng{subscription} [sab-skrip-chon].
\end{itemize}
\end{aretenir}

\newpage

\coursec{Corrigés des exercices}

\coursec{Lesson 51 — Vocabulary: Words and expressions related to exploiting service packages (telephone/internet services) — Corrigés des exercices}

\begin{corrige}[Exercice 1 — QCM : Vocabulaire des services télécoms]
\begin{enumerate}
    \item \textbf{top up} — \eng{I need to top up my phone; I have no airtime left.} « Je dois recharger mon téléphone ; je n'ai plus de crédit. » \eng{Consume} est un faux ami ici ; \eng{forfeit} signifie « perdre (un droit) » ; \eng{subscribe} exige la préposition \eng{to} et un complément (\eng{subscribe to a plan}).
    \item \textbf{expires} — \eng{The data bundle I bought yesterday expires at midnight tonight.} « Le forfait data que j'ai acheté hier expire ce soir à minuit. » \eng{Validate} = valider ; \eng{roll over} = être reporté (ne s'applique pas à une date de fin).
    \item \textbf{subscribe / check} — \eng{Before you subscribe to a new plan, always check the terms and conditions.} Collocations standard : \eng{subscribe to a plan}, \eng{check the T\&Cs}.
    \item \textbf{balance / recharge} — \eng{My balance is low; I cannot make any outgoing calls until I recharge.} \eng{Balance} est le terme affiché à l'écran ; en anglais standard \eng{recharge} est transitif (\eng{recharge my account}), mais l'intransitif \eng{I need to recharge} est fréquent en anglais camerounais.
    \item \textbf{rollover / cycle} — \eng{The rollover feature allows unused data to be carried over to the next billing cycle.} Terminologie technique exacte : \eng{billing cycle} = cycle de facturation.
\end{enumerate}
\textbf{Point de vigilance :} ne jamais traduire « consommer son forfait » par \eng{*consume one's bundle} ; préférer \eng{use up} ou \eng{run out of}.
\end{corrige}

\begin{corrige}[Exercice 2 — Vrai ou Faux justifié]
\begin{enumerate}
    \item \textbf{False.} Justification : \eng{“Mobile internet penetration has grown rapidly in Yaoundé and Douala.”} Le texte parle de grandes villes (zones urbaines), pas de zones rurales.
    \item \textbf{True.} Justification : \eng{“Customers can now purchase daily, weekly, or monthly packages via USSD codes or Mobile Money apps.”}
    \item \textbf{False.} Justification : \eng{“The ‘Night Plan’... offers high-speed data at a reduced rate between 11 PM and 5 AM.”} Il est \emph{moins} cher, pas plus cher.
    \item \textbf{True.} Justification : \eng{“Many users complain about throttling speeds once the fair usage policy limit is reached.”}
    \item \textbf{False.} Justification : \eng{“Subscribers are advised to monitor their data consumption regularly using the operator's self-care app.”} Ils peuvent donc vérifier leur consommation sur leur téléphone.
\end{enumerate}
\textbf{Méthode :} recopier la phrase exacte du texte qui justifie la réponse ; ne jamais répondre seulement V ou F au Bac.
\end{corrige}

\begin{corrige}[Exercice 3 — Dialogue complété]
\textbf{Customer :} \eng{“Hello. I want to \textbf{subscribe} to a new internet \textbf{package} for my smartphone.”} \\
\textbf{Agent :} \eng{“We have a 10 GB \textbf{bundle} that is \textbf{valid} for 30 days. It costs 5\,000 FCFA.”} \\
\textbf{Customer :} \eng{“Does it have the \textbf{rollover} feature? I hate losing unused megabytes at the end of the month.”} \\
\textbf{Agent :} \eng{“Yes, sir. Unused data rolls over to the next cycle if you \textbf{recharge} before the current plan \textbf{expires}.”} \\
\textbf{Customer :} \eng{“Perfect. I'll take it. How do I check my \textbf{balance} later?”} \\
\textbf{Agent :} \eng{“Dial *157\# or use the MyMTN app. You can also buy \textbf{airtime} for voice calls there.”}
\medskip

\textbf{Justifications :} \eng{subscribe to} (verbe + préposition obligatoire) ; \eng{valid for 30 days} (adjectif + \eng{for}) ; \eng{expires} (3\textsuperscript{e} personne du singulier, sujet \eng{the current plan}) ; \eng{airtime} (crédit pour appels vocaux, par opposition à \eng{data}).
\end{corrige}

\begin{corrige}[Exercice 4 — Formation des mots]
\begin{center}
\begin{tabularx}{\linewidth}{|X|X|X|X|}
\hline
\rowcolor{popblueL} \textbf{Verbe} & \textbf{Nom (chose)} & \textbf{Nom (agent)} & \textbf{Adjectif} \\
\hline
subscribe & subscription & subscriber & subscribed \\
\hline
connect & \textbf{connection} & \textbf{connector / user} & connected \\
\hline
validate & validation & \textbf{validator} & \textbf{valid / validated} \\
\hline
activate & \textbf{activation} & activator & \textbf{active / activated} \\
\hline
consume & consumption & consumer & \textbf{consumable} \\
\hline
expire & \textbf{expiry / expiration} & \textemdash & \textbf{expired} \\
\hline
register & registration & \textbf{registrant} & registered \\
\hline
\end{tabularx}
\end{center}
\textbf{Remarques :} \eng{expiry} est la forme britannique courante (\eng{expiry date}), \eng{expiration} la forme américaine ou formelle ; \eng{registrant} désigne la personne qui s'enregistre ; attention à l'orthographe : \eng{connection} (UK) / \eng{connexion} est une variante rare — ne pas écrire \eng{*connecttion} ; l'adjectif usuel pour \eng{consume} est \eng{consumable} (\eng{consuming} est un participe présent).
\end{corrige}

\begin{corrige}[Exercice 5 — Traduction (français $\rightarrow$ anglais)]
\begin{enumerate}
    \item \eng{I used up my data bundle last night; I need to top it up.} (Variante : \eng{I ran out of data last night; I have to recharge.}) — \eng{it} est placé entre \eng{top} et \eng{up} (phrasal verb séparable).
    \item \eng{This plan has not been valid since January 1st.} (Variante : \eng{This bundle expired on January 1st.}) — présent perfect avec \eng{since}.
    \item \eng{Have you activated the rollover option on your line?} — présent perfect pour une action au présent d'actualité.
    \item \eng{Check your balance before calling; you have no credit left.} (Variante : \eng{...you are out of airtime.}) — impératif + \eng{before} + V-ing.
    \item \eng{Operators offer promotional deals at weekends.} (Variante US : \eng{on weekends.})
    \item \eng{I subscribed to an unlimited plan for social media.} — prétérit simple (action ponctuelle passée) ; \eng{an} devant voyelle (\eng{unlimited}).
\end{enumerate}
\textbf{Points de vigilance :} pas de \eng{*forfait}, pas de \eng{*abonnement}, pas de \eng{*I consumed my bundle}.
\end{corrige}

\begin{corrige}[Exercice 6 — Transformation : correction des calques]
\begin{enumerate}
    \item \eng{I used up my bundle yesterday.} / \eng{I ran out of data yesterday.} — \eng{consume} est un faux ami ; \eng{forfait} n'existe pas en anglais.
    \item \eng{My data bundle expired at midnight.} / \eng{My plan ran out at midnight.} — un forfait « expire », il ne « finit » pas (\eng{*finished}).
    \item \eng{I want to subscribe to a new plan.} / \eng{I want to take out a new plan.} — on ne « fait » pas un abonnement ; \eng{*abonnement} est un calque.
    \item \eng{The rate is too high for me.} / \eng{The price is too high.} — un prix est \eng{high/low}, pas \eng{*expensive} (c'est le produit qui est \eng{expensive}) ; \eng{tariff} = tarif douanier.
    \item \eng{He topped up his credit with 1\,000 FCFA.} / \eng{He topped up his phone with 1\,000 FCFA worth of airtime.}
    \item \eng{The network has been slow since this morning.} / \eng{The network has been throttling speeds since this morning.} — présent perfect continu avec \eng{since} ; \eng{*since morning} est un calque.
\end{enumerate}
\end{corrige}

\begin{corrige}[Exercice 7 — Texte à trous]
Réponses dans l'ordre : \textbf{balance} -- \textbf{subscribe} -- \textbf{credit} (ou \textbf{airtime}) -- \textbf{activation} -- \textbf{validity} -- \textbf{rollover} -- \textbf{bill shock}.
\medskip

Texte complet : \eng{“To check your \textbf{balance}, dial *157\#. If you want to \textbf{subscribe} to a weekly data plan, select ‘Data Bundles’ then ‘Weekly’. Make sure you have enough \textbf{credit} on your account before confirming. Once the transaction is successful, you will receive an SMS confirming the \textbf{activation}. Remember that the \textbf{validity} period starts immediately. If you don't use all your data, check if your plan allows \textbf{rollover} to the next month. To avoid \textbf{bill shock}, monitor your usage daily.”}
\medskip

\textbf{Justifications :} \eng{check one's balance} (collocation) ; \eng{subscribe to} ; \eng{activation} (nom de l'action confirmée par SMS, pas \eng{expiration}) ; \eng{validity period} (collocation) ; \eng{rollover} (seul mot compatible avec \eng{to the next month}) ; \eng{bill shock} (surprise de facture que l'on « évite »).
\end{corrige}

\begin{corrige}[Exercice 8 — Appariement : collocations]
\textbf{Correspondances :} 1-d ; 2-e ; 3-b ; 4-c ; 5-a.
\medskip

\textbf{Phrases originales possibles :}
\begin{enumerate}
    \item \eng{I need to top up my credit before I call my uncle in Bamenda.} « Je dois recharger mon crédit avant d'appeler mon oncle à Bamenda. »
    \item \eng{She ran out of data while downloading her lecture notes.} « Elle a épuisé ses données en téléchargeant ses notes de cours. »
    \item \eng{Please check your balance before subscribing to any bundle.} « Vérifie ton solde avant de t'abonner à un forfait. »
    \item \eng{You must activate the service by dialling *157\#.} « Tu dois activer le service en composant le *157\#. »
    \item \eng{My brother browses the internet every evening after school.} « Mon frère navigue sur internet chaque soir après l'école. »
\end{enumerate}
\textbf{Point de vigilance :} \eng{run out of} est suivi directement du nom sans article possessif obligatoire (\eng{run out of data}, \eng{run out of credit}).
\end{corrige}

\begin{corrige}[Exercice 9 — Compréhension et langue (type Bac)]
\textbf{Barème indicatif (10 pts) :} compréhension 5 $\times$ 1 pt ; langue 5 $\times$ 1 pt.
\medskip

\textbf{A. Questions de compréhension :}
\begin{enumerate}
    \item \eng{He buys airtime via Mobile Money (MTN MoMo or Orange Money) on his smartphone.} « Il achète du crédit via Mobile Money sur son smartphone. »
    \item \eng{Two tailored offers are “Social Packs” and “Night Owl” bundles.} (Accepter aussi \eng{Family Share plans}.)
    \item \eng{During peak hours, network congestion leads to throttling, so the speed drops.} « Pendant les heures de pointe, la congestion du réseau provoque un bridage du débit. »
    \item \eng{On so-called “unlimited” plans, the FUP caps high-speed data at 2 GB or 5 GB per day.} « La FUP plafonne le volume de données haut débit à 2 ou 5 Go par jour. »
    \item \eng{The text advises users to disable auto-renewal via USSD and to read the terms and conditions before subscribing.}
\end{enumerate}

\textbf{B. Questions de langue :}
\begin{enumerate}
    \setcounter{enumi}{5}
    \item Synonyme de \eng{customised} : \textbf{tailored} (paragraphe 2).
    \item Contraire de \eng{visible / obvious} : \textbf{hidden} (paragraphe 3 : \eng{hidden in fine print}).
    \item Voix passive : \eng{The regulator is urged to enforce transparency (by consumers).} — auxiliaire \eng{be} au présent + participe passé \eng{urged} ; le complément d'agent \eng{by consumers} est facultatif.
    \item \eng{If Pierre \textbf{does not disable} auto-renewal, he \textbf{will lose} credit automatically next week.} — conditionnelle de type 1 : \eng{if} + présent simple, futur avec \eng{will} dans la principale. Ne jamais mettre \eng{will} dans la subordonnée (\eng{*if he will not disable}).
    \item \eng{Run out of} signifie « épuiser sa réserve de, ne plus avoir de » : \eng{He ran out of data yesterday} = « Il a épuisé ses données hier. » C'est un phrasal verb inséparable suivi d'un complément.
\end{enumerate}

\textbf{Points de vigilance :} en compréhension, répondre par une phrase complète en s'appuyant sur le texte ; à la voix passive, vérifier l'accord du participe passé et le temps de l'auxiliaire.
\end{corrige}

\begin{corrige}[Exercice 10 — Traduction : thème (type Bac)]
\textbf{Barème indicatif (5 pts) :} 1 pt par segment correct (5 segments), tolérance zéro sur les calques lexicaux (\eng{*forfait}, \eng{*abonnement}, \eng{*consommé}).
\medskip

\textbf{Traduction modèle :}

\eng{“In Cameroon, choosing an internet plan is not easy. Yesterday, I wanted to subscribe to a monthly 20 GB bundle, but my balance was insufficient. I had to top up my account via Mobile Money. The seller explained to me that this plan had a validity of 30 days and that the rollover option for unused data was activated automatically. However, he warned me that the speed would be reduced after 5 GB per day because of the fair usage policy. I am going to compare the offers of the two operators before making a final decision.”}

« Au Cameroun, choisir un forfait internet n'est pas facile. Hier, j'ai voulu m'abonner à un forfait mensuel de 20 Go, mais mon solde était insuffisant. J'ai dû recharger mon compte via Mobile Money. Le vendeur m'a expliqué que ce forfait avait une validité de 30 jours et que l'option de report des données non utilisées était activée automatiquement. Cependant, il m'a prévenu que la vitesse serait réduite après 5 Go par jour à cause de la politique d'utilisation équitable. Je vais comparer les offres des deux opérateurs avant de prendre une décision finale. »
\medskip

\textbf{Points de vigilance :}
\begin{itemize}
    \item « s'abonner à » $\rightarrow$ \eng{subscribe to} (préposition obligatoire) ;
    \item « recharger » $\rightarrow$ \eng{top up} ou \eng{recharge} (jamais \eng{*charge}) ;
    \item « la vitesse serait réduite » $\rightarrow$ conditionnel \eng{would be reduced} (discours rapporté au passé) ;
    \item « je vais comparer » $\rightarrow$ \eng{I am going to compare} (projet) ou \eng{I will compare} ;
    \item « 20 Go » $\rightarrow$ \eng{20 GB} (gigabytes) ; « FCFA » reste en anglais camerounais.
\end{itemize}
\end{corrige}

\begin{corrige}[Exercice 11 — Expression écrite : composition (type Bac)]
\textbf{Barème indicatif (10 pts) :} respect des contraintes (80--100 mots, 8 mots du lexique soulignés) : 2 pts ; richesse et exactitude du vocabulaire de la leçon : 3 pts ; correction grammaticale : 3 pts ; organisation (introduction, comparaison, conseil) : 2 pts.
\medskip

\textbf{Proposition de rédaction modèle (environ 95 mots) :}

\eng{“Hi Sandra! For your studies, I advise you to compare two offers before you \textbf{subscribe}. The MTN Student Pack gives you 15 GB with a 30-day \textbf{validity} for 4\,500 FCFA, while the Orange Edu Bundle offers 10 GB plus free Wikipedia for 3\,500 FCFA. Both include \textbf{rollover}, so unused data is not lost if you renew before \textbf{expiry}. Be careful: MTN applies a strict \textbf{fair usage policy} and \textbf{throttling} starts after 3 GB per day. \textbf{Top up} your \textbf{balance} via Mobile Money, disable \textbf{auto-renewal}, and always read the terms and conditions first!”}
\medskip

\textbf{Points de vigilance :}
\begin{itemize}
    \item Souligner au moins 8 mots du lexique (ici : \eng{subscribe}, \eng{validity}, \eng{rollover}, \eng{expiry}, \eng{fair usage policy}, \eng{throttling}, \eng{top up}, \eng{balance}, \eng{auto-renewal}) ;
    \item éviter les calques : pas de \eng{*forfait}, \eng{*consommer la data}, \eng{*faire un abonnement} ;
    \item registre semi-formel adapté à un conseil entre amis (\eng{I advise you to...}, \eng{Be careful:...}) ;
    \item vérifier la 3\textsuperscript{e} personne du singulier (\eng{the pack gives}, \eng{throttling starts}) et les comparatifs (\eng{cheaper than}, \eng{more data than}).
\end{itemize}
\end{corrige}

\newpage

\coursec{Fiche bilan}

\coursec{Lesson 51 — Vocabulary: Words and expressions related to exploiting service packages (telephone/internet services)}

\begin{popfigure}
\begin{center}\tcbox[colback=poporange,colframe=poporange,coltext=white,arc=9pt,boxsep=4pt,left=6pt,right=6pt]{\bfseries\small Exploiting Service Packages (Phone / Internet)}\end{center}
\vspace{-2pt}
\begin{tcbraster}[raster columns=3,raster equal height=rows,raster column skip=6pt,raster row skip=6pt,enhanced,blankest]
\begin{tcolorbox}[enhanced,colback=poporange,colframe=poporange,coltext=white,boxrule=0.9pt,arc=3pt,left=4pt,right=4pt,top=3pt,bottom=3pt,boxsep=1pt,halign=left]\bfseries\scriptsize Phone \& Calls \textit{call, ring, pick up, hang up, voicemail, roaming, prepaid, credit}\end{tcolorbox}
\begin{tcolorbox}[enhanced,colback=popgreen,colframe=popgreen,coltext=white,boxrule=0.9pt,arc=3pt,left=4pt,right=4pt,top=3pt,bottom=3pt,boxsep=1pt,halign=left]\bfseries\scriptsize Internet \& Data \textit{broadband, bandwidth, Wi-Fi, hotspot, data cap, throttle, stream, download, upload}\end{tcolorbox}
\begin{tcolorbox}[enhanced,colback=poppurple,colframe=poppurple,coltext=white,boxrule=0.9pt,arc=3pt,left=4pt,right=4pt,top=3pt,bottom=3pt,boxsep=1pt,halign=left]\bfseries\scriptsize Plans \& Billing \textit{subscription, bundle, plan, package, bill, invoice, top-up, recharge, auto-renew, term}\end{tcolorbox}
\begin{tcolorbox}[enhanced,colback=poppink,colframe=poppink,coltext=white,boxrule=0.9pt,arc=3pt,left=4pt,right=4pt,top=3pt,bottom=3pt,boxsep=1pt,halign=left]\bfseries\scriptsize False Friends \textit{forfeit $\neq$ forfait, abonnement $\neq$ subscription, factory $\neq$ facture, credit $\neq$ mark}\end{tcolorbox}
\begin{tcolorbox}[enhanced,colback=popgold,colframe=popgold,coltext=white,boxrule=0.9pt,arc=3pt,left=4pt,right=4pt,top=3pt,bottom=3pt,boxsep=1pt,halign=left]\bfseries\scriptsize Word Families \textit{subscribe $\to$ subscription/subscriber, connect $\to$ connection/connectivity, bill $\to$ billing, pay $\to$ payment/payer}\end{tcolorbox}
\begin{tcolorbox}[enhanced,colback=popblue,colframe=popblue,coltext=white,boxrule=0.9pt,arc=3pt,left=4pt,right=4pt,top=3pt,bottom=3pt,boxsep=1pt,halign=left]\bfseries\scriptsize In Context \textit{Dialogue: Top-up / Support call; Text: Choosing a bundle; Collocations: run out of credit, exceed limit}\end{tcolorbox}
\end{tcbraster}

\legende{Carte mentale : Vocabulaire des forfaits téléphone / internet (Tle A — Lesson 51)}
\end{popfigure}

\courssub{Lexique clé — Anglais / Français}

\textbf{1. Téléphone et appels (Phone \& Calls)}
\begin{center}
\begin{tabularx}{\linewidth}{|X|X|X|}\hline
\rowcolor{poporangeL} \textbf{English} & \textbf{Français} & \textbf{Collocation / Exemple} \\ \hline
\cle{landline} / \cle{fixed line} & ligne fixe & a \eng{landline number} — un numéro de fixe \\ \hline
\cle{mobile phone} / \cle{cell phone (US)} & téléphone portable & \eng{My mobile is off.} — Mon portable est éteint. \\ \hline
\cle{smartphone} & smartphone & \eng{smartphone user} — utilisateur de smartphone \\ \hline
\cle{prepaid} / \cle{pay-as-you-go} & prépayé / carte & a \eng{prepaid SIM card} — une carte SIM prépayée \\ \hline
\cle{credit} / \cle{airtime} & crédit / temps d'appel & \eng{run out of credit} — ne plus avoir de crédit \\ \hline
\cle{top-up} (n/v) / \cle{recharge} & recharger / recharge & \eng{I need to top up my phone.} — Je dois recharger. \\ \hline
\cle{roaming} & itinérance (étranger) & \cle{data roaming} — données en itinérance \\ \hline
\cle{voicemail} / \cle{answerphone (UK)} & messagerie vocale & \eng{leave a voicemail} — laisser un message \\ \hline
\cle{missed call} & appel manqué & \eng{I have three missed calls.} — J'ai trois appels manqués. \\ \hline
\cle{dial} / \cle{call} & composer / appeler & \eng{dial a number} — composer un numéro \\ \hline
\cle{pick up} / \cle{answer} & décrocher & \eng{He didn't pick up.} — Il n'a pas décroché. \\ \hline
\cle{hang up} & raccrocher & \eng{Don't hang up!} — Ne raccroche pas ! \\ \hline
\cle{ringtone} & sonnerie & \eng{change my ringtone} — changer de sonnerie \\ \hline
\cle{contact list} / \cle{phonebook} & répertoire / contacts & \eng{save to contacts} — enregistrer dans contacts \\ \hline
\end{tabularx}
\end{center}

\textbf{2. Internet et données (Internet \& Data)}
\begin{center}
\begin{tabularx}{\linewidth}{|X|X|X|}\hline
\rowcolor{popgreenL} \textbf{English} & \textbf{Français} & \textbf{Collocation / Exemple} \\ \hline
\cle{broadband} & haut débit / large bande & \eng{broadband connection} — connexion haut débit \\ \hline
\cle{fibre (UK)} / \cle{fiber (US)} / \cle{fibre optics} & fibre optique & \eng{fibre broadband} — fibre jusqu'au domicile (FTTH) \\ \hline
\cle{Wi-Fi} / \cle{wireless} & Wi-Fi / sans fil & \eng{connect to Wi-Fi} — se connecter au Wi-Fi \\ \hline
\cle{hotspot} & point d'accès / borne & \eng{mobile hotspot} — partage de connexion \\ \hline
\cle{bandwidth} & bande passante & \eng{high bandwidth} — grande bande passante \\ \hline
\cle{data allowance} / \cle{data cap} & quota de données / plafond & \eng{exceed my data cap} — dépasser mon forfait data \\ \hline
\cle{throttle} (v) / \cle{throttling} (n) & brider / bridage & \eng{The ISP throttles speed after 10 GB.} — Le FAI bride la vitesse après 10 Go. \\ \hline
\cle{stream} (v) / \cle{streaming} (n) & streamer / streaming & \eng{stream videos in HD} — regarder des vidéos en HD \\ \hline
\cle{download} / \cle{upload} & télécharger / envoyer (vers serveur) & \eng{download speed} / \eng{upload speed} — débit descendant / montant \\ \hline
\cle{buffering} & mise en tampon (chargement) & \eng{constant buffering} — chargement constant \\ \hline
\cle{lag} (n/v) & latence / ralentir & \eng{The game lags.} — Le jeu a du lag / ralentit. \\ \hline
\cle{router} / \cle{modem} & routeur / box (modem) & \eng{restart the router} — redémarrer la box \\ \hline
\cle{ISP} (Internet Service Provider) & FAI (Fournisseur d'Accès Internet) & \eng{MTN, Orange, Camtel are ISPs.} — MTN, Orange, Camtel sont des FAI. \\ \hline
\end{tabularx}
\end{center}

\textbf{3. Forfaits, abonnements et facturation (Plans, Subscriptions \& Billing)}
\begin{center}
\begin{tabularx}{\linewidth}{|X|X|X|}\hline
\rowcolor{poppurpleL} \textbf{English} & \textbf{Français} & \textbf{Collocation / Exemple} \\ \hline
\cle{plan} / \cle{package} / \cle{bundle} & forfait / pack / offre groupée & a \eng{monthly plan} — un forfait mensuel \\ \hline
\cle{subscription} & abonnement & \eng{monthly subscription} — abonnement mensuel \\ \hline
\cle{subscribe to} (v) & s'abonner à & \eng{subscribe to a 4G plan} — s'abonner à un forfait 4G \\ \hline
\cle{subscriber} & abonné & \eng{a loyal subscriber} — un abonné fidèle \\ \hline
\cle{unlimited} & illimité & \eng{unlimited calls and texts} — appels/SMS illimités \\ \hline
\cle{capped} / \cle{limited} & plafonné / limité & a \cle{capped plan} — un forfait plafonné \\ \hline
\cle{rollover} / \cle{carry over} & report (crédit/data non utilisé) & \eng{data rollover} — report de data \\ \hline
\cle{auto-renewal} / \cle{auto-renew} & renouvellement automatique & \eng{turn off auto-renewal} — désactiver le renouvellement auto \\ \hline
\cle{bill} / \cle{invoice} & facture & \eng{pay the bill} — payer la facture \\ \hline
\cle{billing cycle} / \cle{billing period} & cycle de facturation / période & \eng{end of billing cycle} — fin du cycle de facturation \\ \hline
\cle{due date} & date d'échéance & \eng{The due date is the 5th.} — L'échéance est le 5. \\ \hline
\cle{out-of-bundle} / \cle{out-of-plan} & hors forfait & \cle{out-of-bundle rates} — tarifs hors forfait \\ \hline
\cle{top-up voucher} / \cle{recharge card} & carte de recharge & \eng{buy a top-up voucher} — acheter une carte de recharge \\ \hline
\cle{term} / \cle{contract term} & durée d'engagement & a \eng{12-month term} — engagement 12 mois \\ \hline
\cle{early termination fee} & frais de résiliation anticipée & \eng{pay an early termination fee} — payer des frais de résiliation anticipée \\ \hline
\end{tabularx}
\end{center}

\courssub{Faux amis et calques du français — \textbf{Attention !}}
\begin{center}
\begin{tabularx}{\linewidth}{|X|X|X|}\hline
\rowcolor{poppinkL} \textbf{Mot anglais} & \textbf{Piège (sens français)} & \textbf{Vrai sens anglais / Correction} \\ \hline
\cle{forfeit} & \eng{forfait} (package) & \eng{forfeit} = perte, pénalité, confiscation. \textbf{Dire :} \eng{plan, package, bundle}. \\ \hline
\cle{abonnement} (n'existe pas) & \eng{abonnement} & \textbf{Dire :} \eng{subscription} (n), \eng{subscribe} (v), \eng{subscriber} (n). \\ \hline
\cle{factory} & \eng{facture} & \eng{factory} = usine. \textbf{Dire :} \eng{bill} ou \eng{invoice}. \\ \hline
\cle{credit} (note) & \eng{crédit} (note à l'école) & \eng{credit} = crédit (argent, téléphone), note = \eng{mark} / \eng{grade}. \\ \hline
\cle{recharge} (v/n) & \eng{recharger} (batterie) & OK pour téléphone (\eng{top-up} plus courant UK). Pour batterie : \eng{charge} / \eng{recharge}. \\ \hline
\cle{consume} / \cle{consumption} & \eng{consommer} (data) & Possible mais technique. \textbf{Préférer :} \eng{use data}, \eng{data usage}. \\ \hline
\cle{expire} / \cle{expiry date} & \eng{expirer} / \eng{date d'expiration} & Correct. Attention : \eng{expiry date} (UK), \eng{expiration date} (US). \\ \hline
\cle{reclaim} & \eng{réclamer} (remboursement) & \eng{reclaim} = récupérer. \textbf{Dire :} \eng{claim a refund}, \eng{dispute a charge}. \\ \hline
\end{tabularx}
\end{center}

\courssub{Mots de la même famille — Formation des mots (Word Formation)}
\begin{center}
\begin{tabularx}{\linewidth}{|X|X|X|X|}\hline
\rowcolor{popgoldL} \textbf{Verbe} & \textbf{Nom (action/état)} & \textbf{Nom (agent)} & \textbf{Adjectif / Participe} \\ \hline
\eng{subscribe} & \eng{subscription} & \eng{subscriber} & \eng{subscribed} / \eng{subscription-based} \\ \hline
\eng{connect} & \eng{connection} / \eng{connectivity} & — & \eng{connected} / \eng{connectable} \\ \hline
\eng{bill} & \eng{billing} & \eng{biller} (rare) & \eng{billed} / \eng{billable} \\ \hline
\eng{pay} & \eng{payment} & \eng{payer} & \cle{paid} / \eng{payable} \\ \hline
\eng{charge} & \eng{charge} / \eng{charging} & \eng{charger} (appareil) & \eng{charged} / \eng{chargeable} \\ \hline
\eng{renew} & \eng{renewal} & — & \eng{renewed} / \eng{renewable} \\ \hline
\eng{terminate} & \eng{termination} & — & \eng{terminated} / \eng{terminal} (différent) \\ \hline
\eng{activate} & \eng{activation} & — & \eng{active} / \eng{activated} \\ \hline
\eng{deactivate} & \eng{deactivation} & — & \eng{inactive} / \eng{deactivated} \\ \hline
\end{tabularx}
\end{center}

\courssub{Méthodes express}
\begin{methode}[Apprendre le vocabulaire thématique efficacement]
\begin{enumerate}
\item \textbf{Regrouper par champ :} Phone / Internet / Billing (3 listes).
\item \textbf{Associer collocation + traduction :} \eng{run out of credit} $\to$ « ne plus avoir de crédit » (pas *finish my credit*).
\item \textbf{Cartes mentales :} Dessiner la carte (ci-dessus) de mémoire.
\item \textbf{Contexte réel :} Lire une facture (\eng{invoice}) ou un site d'opérateur (MTN, Orange, Camtel) en version anglaise.
\item \textbf{Spaced repetition :} Réviser J+1, J+3, J+7 avec Anki/Quizlet.
\end{enumerate}
\end{methode}

\begin{methode}[Réussir un dialogue « Service client / Recharge »]
\begin{enumerate}
\item \textbf{Saluer \& identifier :} \eng{Good morning, I'd like to top up my line.}
\item \textbf{Préciser le besoin :} \eng{I need a 5000 FCFA recharge card / a 10 GB data bundle.}
\item \textbf{Valider :} \eng{What is the USSD code?} / \eng{Can you activate it for me?}
\item \textbf{Problème fréquent :} \eng{My data is exhausted.} / \eng{I have no network.}
\item \textbf{Clôturer :} \eng{Thank you, have a nice day.}
\end{enumerate}
\end{methode}

\begin{aretenir}[Checklist avant le Bac — Lesson 51]
$\square$ Je distingue \cle{plan}, \cle{bundle}, \cle{package} et je ne dis jamais \eng{forfeit} pour « forfait ».\par
$\square$ Je sais dire « recharger » : \eng{top up} (UK) / \eng{recharge} / \eng{add credit} et le nom \eng{top-up voucher}.\par
$\square$ Je maîtrise les verbes d'action téléphone : \eng{dial, pick up, hang up, leave a voicemail, miss a call}.\par
$\square$ Je connais les termes data : \cle{data cap/allowance}, \cle{throttle}, \cle{stream}, \cle{download/upload speed}, \cle{hotspot}.\par
$\square$ Je sais former la famille de \eng{subscribe} : \eng{subscription, subscriber, subscribed}.\par
$\square$ J'évite les faux amis : \eng{factory} (usine) $\neq$ facture (\eng{bill/invoice}) ; \eng{forfeit} $\neq$ forfait.\par
$\square$ Je sais lire une facture : \eng{billing cycle, due date, invoice number, amount due, VAT}.\par
$\square$ J'utilise les prépositions correctes : \eng{subscribe TO}, \eng{pay FOR}, \eng{charge FOR}, \eng{run OUT OF}.\par
$\square$ Je peux rédiger une plainte formelle (\eng{complaint letter/email}) : objet, référence client, faits, demande précise.\par
$\square$ Je peux jouer un dialogue « Client / Service client » pour activer un forfait ou signaler un problème de connexion.\par
$\square$ Je connais les abréviations clés : \eng{ISP, SIM, USSD, VoLTE, 4G/5G, Wi-Fi, SMS, MMS}.\par
\end{aretenir}

\findecours

\end{document}
